\documentclass[11pt]{article}

\usepackage[T1]{fontenc}
\usepackage[margin=1in]{geometry}
\usepackage{lmodern}
\usepackage{microtype}
\usepackage{amsmath,amssymb,mathtools}
\usepackage{amsthm}
\usepackage{aliascnt}
\usepackage{booktabs}
\usepackage{tabularx}
\usepackage[font=small,labelfont=bf]{caption}
\usepackage[section]{placeins}
\usepackage{float}
\usepackage{array}
\usepackage{xcolor}
\usepackage{enumitem}
\usepackage{needspace}
\usepackage[hidelinks]{hyperref}
\usepackage[nameinlink,noabbrev]{cleveref}
\usepackage{url}
\usepackage[backend=biber,style=numeric,sorting=nyt,maxbibnames=99]{biblatex}
\addbibresource{references.bib}


\newtheorem{theorem}{Theorem}

\newaliascnt{lemma}{theorem}

\newtheorem{lemma}[lemma]{Lemma}
\aliascntresetthe{lemma}

\newaliascnt{proposition}{theorem}

\newtheorem{proposition}[proposition]{Proposition}
\aliascntresetthe{proposition}

\newaliascnt{corollary}{theorem}

\newtheorem{corollary}[corollary]{Corollary}
\aliascntresetthe{corollary}
\crefname{theorem}{Theorem}{Theorems}
\crefname{lemma}{Lemma}{Lemmas}
\crefname{proposition}{Proposition}{Propositions}
\crefname{corollary}{Corollary}{Corollaries}

\definecolor{deepblue}{RGB}{23,69,124}
\definecolor{lightblue}{RGB}{242,247,252}
\definecolor{darkred}{RGB}{145,32,32}
\ifdefined\ANONYMOUS

\newcommand{\pdfpaperauthor}{Anonymous}
\else

\newcommand{\pdfpaperauthor}{Justin Thaler, Kai Zhang, Scott Kominers, Quang Dao}
\fi
\hypersetup{
  colorlinks=true,
  linkcolor=deepblue,
  citecolor=deepblue,
  urlcolor=deepblue,
  pdftitle={The Symmetry-Quotient Attack on Odd-Characteristic SNOVA},
  pdfauthor={\pdfpaperauthor},
  pdfsubject={Cryptanalysis of SNOVA Version 2.3},
  pdfkeywords={SNOVA, multivariate signatures, forgery attack, NIST PQC}
}
\setlist{leftmargin=*,topsep=3pt,itemsep=1.5pt,parsep=0pt,partopsep=0pt}


\newcommand{\F}{\mathbb F}

\newcommand{\Mat}{\operatorname{Mat}}

\newcommand{\Sym}{\operatorname{Sym}}

\newcommand{\rank}{\operatorname{rank}}

\newcommand{\Span}{\operatorname{span}}

\newcommand{\transpose}{\mathsf T}

\newcommand{\SNOVA}{\textsf{SNOVA}}

\newcommand{\bits}{\text{ bits}}

\newcommand{\LoneModelRange}{$2^{128.90}$--$2^{129.51}$}

\newcommand{\LthreeModelRange}{$2^{171.76}$--$2^{183.51}$}

\newcommand{\LfiveModelRange}{$2^{214.20}$--$2^{235.27}$}

\newcommand{\LoneGenericFallback}{$2^{136.11}$}

\newcommand{\LthreeGenericFallback}{$2^{180.66}$}

\newcommand{\LfiveGenericFallback}{$2^{224.75}$}

\newcommand{\LoneAXNRange}{$2^{132.25}$--$2^{133.82}$}

\newcommand{\LthreeAXNRange}{$2^{175.12}$--$2^{187.82}$}

\newcommand{\LfiveAXNRange}{$2^{217.56}$--$2^{239.57}$}

\newcommand{\LoneNANDRange}{$2^{133.12}$--$2^{134.74}$}

\newcommand{\LthreeNANDRange}{$2^{175.99}$--$2^{188.74}$}

\newcommand{\LfiveNANDRange}{$2^{218.42}$--$2^{240.50}$}

\newcommand{\AXNNominalHeadroomRange}{9.18--54.44}

\newcommand{\AXNExpectedAESGapRange}{8.06--52.78}

\newcommand{\AXNFullAESGapRange}{9.06--53.78}

\newcommand{\NANDNominalHeadroomRange}{8.26--53.58}

\newcommand{\NANDExpectedAESGapRange}{8.96--53.75}

\newcommand{\NANDFullAESGapRange}{9.96--54.75}

\newcommand{\attacktitle}{The Symmetry-Quotient Attack on
  Odd-Characteristic \texorpdfstring{\SNOVA}{SNOVA}}

\title{\attacktitle}
\ifdefined\ANONYMOUS
  \author{Anonymous submission}
\else
  \author{%
    Justin Thaler\thanks{a16z crypto research and Georgetown University} \and
    Kai Zhang\thanks{MIT} \and
    Scott Kominers\thanks{Harvard University and a16z crypto research} \and
    Quang Dao\thanks{Carnegie Mellon University and LayerZero Labs}%
  }
\fi
\date{28 July 2026}

\begin{document}
\setlength{\emergencystretch}{2em}
\raggedbottom
\maketitle

\ifdefined\DISCLOSURE
\begin{center}
\fcolorbox{darkred}{white}{\parbox{0.91\textwidth}{\centering\small
\textbf{Confidential coordinated disclosure.} This draft is intended for the
\SNOVA{} submission team and NIST.\@ Please do not redistribute before the
coordinated public-disclosure date.}}
\end{center}
\fi

\begin{abstract}
The odd-characteristic ($q=19$) redesign of \SNOVA{} Version~2.3 has a
structural loss of quadratic dimension that holds for every key.  Each public
base form is built from ordered pairs of powers $(\xi,\eta)$ of a public
matrix, but
symmetry of the base matrices makes the feature labeled $(\xi,\eta)$ equal to
the one labeled $(\eta,\xi)$.  Writing $d=\dim_{\F_q}\F_q[S]$, the $d^2$ ordered
features per form thus factor through $\Sym^2$ and span at most
$\binom{d+1}{2}$ dimensions: $16$ collapse to $10$ when $\ell=4$, and $4$ to
$3$ when $\ell=2$.  The identification happens before the public outer map and
holds for every symmetric public key.

We turn this identity into exact forgery reductions for all nine published
$q=19$ parameter sets.  Each reduction keeps every verifier equation, handles
the symmetric-signature rejection rule, and is unconditional: its correctness
follows from public linear algebra.  Feeding the reduced systems to standard MQ
solvers puts the forgery cost below NIST's Level~I/III/V references
($143/207/272$) at all three levels, with the six $\ell=4$ sets remaining below
reference even under a generic-MQ estimator and under solver branches that
assume no selected Macaulay degree.

Patching the flaw is expensive.  Restoring the quadratic-image dimension that
the ordered-label security analysis assumed, while keeping symmetry and
Version~2.3's parameter coupling, requires $31$--$61\%$ more verifier outputs.
\end{abstract}

\section{Introduction}\label{sec:intro}

\paragraph{The cryptanalytic task.}
A multivariate signature scheme publishes a polynomial map of total degree
at most two,
\[
  \mathcal P:\F_q^{N}\longrightarrow\F_q^{M}.
\]
For a message with hash-derived target $t\in\F_q^M$, the algebraic core of a
forgery is to find $x$ satisfying $\mathcal P(x)=t$; any additional format or
rejection checks must also pass.  This is a multivariate-quadratic (MQ)
problem: $M$ equations of degree at most two in the coordinates of $x$.  Linear and constant terms are
allowed; throughout the paper, ``quadratic'' refers to the maximum degree.
The apparent equation and variable counts are only a first approximation to
security.  Public algebraic structure can make different-looking quadratic
features identical, can expose linear equations, or can split the remaining
variables into blocks that a specialized solver can exploit.

\SNOVA{} is one of the nine digital-signature candidates that NIST advanced to
the third round of its additional post-quantum signature process in
May~2026~\cite{NISTThirdRound2026,NISTIR8610}.  The 2026 wedge attack breaks
the updated Round-2 construction in characteristic two
\cite{BrosEtAl2026SNOVAWedge}.  NIST describes the odd-characteristic system
studied here as a proposed response: it moves to $q=19$ and replaces the
asymmetric bilinear public forms by symmetric quadratic forms.  NIST also
states that \SNOVA{} has not reached a stable form and that more substantial
changes may be considered than for more mature candidates~\cite{NISTIR8610}.
The object of this paper is therefore both concrete and time-sensitive: the
public $q=19$ Version~2.3 construction, not the broken characteristic-two
submission and not every future third-round revision.

\SNOVA{} stores the signature variables in small matrices and builds its public
quadratics from pairs of powers of a public matrix $S$.  Beullens showed that
an attacker can impose an affine relation among the signature columns and
thereby reduce forgery to a smaller underdetermined MQ system
\cite{Beullens2024SNOVA}.  We find an additional, deterministic reduction
inside the quadratic features themselves.

\paragraph{The hidden duplication.}
Let $P_i$ be a public symmetric base matrix, let
$A=\F_q[S]$, and write $\Sigma_\xi=I_n\otimes\xi$ for $\xi\in A$.  The feature
with left label $\xi$ and right label $\eta$ is
\[
  q_{i,\xi,\eta}(u)
  =u^\transpose\Sigma_\xi P_i\Sigma_\eta u.
\]
Because the matrices are symmetric and the result is a scalar,
\begin{equation}\label{eq:intro-swap}
  q_{i,\xi,\eta}(u)=q_{i,\eta,\xi}(u).
\end{equation}
The verifier begins from ordered labels, but the polynomial depends only on
the corresponding unordered pair.  We call the operation of identifying
these duplicate labels the \emph{symmetry quotient}.

For a two-dimensional toy algebra with basis $\{1,S\}$, the four labels
\[
 (1,1),\qquad(1,S),\qquad(S,1),\qquad(S,S)
\]
produce only three polynomials because the middle two are equal.  More
generally, if $\dim_{\F_q}A=d$, then at most
\begin{equation}\label{eq:intro-K}
  K=m_1\binom{d+1}{2}
\end{equation}
independent power-pair features remain across the $m_1$ base forms, rather
than $m_1d^2$.  Thus each $\ell=4$ base form drops from $16$ ordered labels to
$10$ unordered labels; each $\ell=2$ base form drops from $4$ to $3$.  The
loss occurs before the public outer linear map---the final public mixing of
these features into verifier outputs---so later randomization cannot recreate
the discarded directions.  It holds for every symmetric public key
and every attacker-chosen column relation.

\begin{center}
\fcolorbox{deepblue}{lightblue}{\parbox{0.94\textwidth}{\small
\textbf{Core idea.}  The verifier assigns two names to the
same quadratic polynomial whenever it reverses a pair of power labels.  The
attack first removes those duplicate features, then uses ordinary linear
algebra to retain all remaining verifier equations, and finally solves only
the smaller quadratic core.}}
\end{center}

\paragraph{The result.}
The reductions are exact and public.  Under the standard generic-MQ estimator,
forgery falls below NIST's Level~I/III/V references ($143/207/272$) at all
three levels, so Version~2.3 fails its targeted security at every level.

\paragraph{What is exact, and what is modeled.}
The paper separates algebraic statements from concrete-cost inputs.
\Cref{tab:claim-ledger} summarizes the split used throughout.

\begin{table}[tbp]
\centering
\caption{Claim ledger.  A public certificate can be checked from the public key
and a solver transcript; ``model input'' means that the displayed exponent
assumes the cited estimator continues to describe the production instance.}
\label{tab:claim-ledger}
\small
\renewcommand{\arraystretch}{1.08}
\begin{tabularx}{0.97\textwidth}{@{}>{\bfseries\raggedright\arraybackslash}p{0.25\textwidth}>{\raggedright\arraybackslash}p{0.19\textwidth}X@{}}
\toprule
Statement & Status & What a verifier checks or assumes\\
\midrule
Symmetry quotient & Unconditional theorem & Symmetry of $S$ and the public
base matrices; no random-key hypothesis.\\
Frobenius-channel normal form & Unconditional theorem and exact certificate &
$\Sym^2_{\F_q}(\F_{q^4})\cong\F_{q^4}\oplus\F_{q^4}\oplus\F_{q^2}$;
the released $10\times10$ matrix has rank ten over $\F_{19}$.\\
Ordered-cap restoration floor & Unconditional necessary condition & To
restore the ordered-label quadratic-image cap while retaining symmetry, a
design needs both enough outputs and enough symmetric-square features.  Under
the current coupling this forces a $31.25\%$--$60.71\%$ output increase across
the nine rows, before any security margin.\\
Affine reduction and rejection handling & Exact public certificate/theorem &
An invertible target-coordinate split, matrix ranks, fixed-skew identities,
and an accepted-density singleton bound for the square $\ell=4$ rows.\\
Root and retry bounds & Exact theorem; seeded-model theorem or fixed-key
certificate & A reserved-line construction leaves a fresh coefficient block
and bounds the bad aggregate spectrum in the random-XOF idealization.  For
a fixed key, a finite projective rank histogram certifies the same sum exactly.\\
Special-XL extraction & Exact correctness once certified & Macaulay row
combinations, a complete border prebasis, commuting multiplication matrices,
and the original verifier.  Quotient dimension $D>1$ is sound but may require
additional block linear algebra.\\
Implicit Macaulay state & Exact operator; standard black-box solver & Rows and
transpose products are generated on demand.  Random normal preconditioning
uses $O(M)$ or $O(bM)$ field-vector state.  Both the compact $\F_{19^8}$ and
Level-I robust $\F_{19^{12}}$ scalar realizations are priced from exact row
counts.\\
Displayed special-XL exponent & Model input with explicit sensitivity & The
compact rows retain the published $D=1$ prediction.  The robust rows instead
charge a finite-quotient certificate probability only $\sigma\ge1/8$; any
larger block/kernel work still has to fit the displayed ledger.\\
Channel-homotopy branch & Theorem-backed nonsingular-root recovery; explicit
public preflight and sensitivity factor & A nonzero leading-Jacobian witness,
the finite-cardinality homotopy theorem, and exact filtering remove the
selected-degree assumption.  The displayed leading exponent leaves
$\kappa_{\rm hom}$ explicit and is not a production timing or a low-memory
claim.\\
Basis-complete/square-slice homotopy & Theorem-backed recovery without a fixed
core or selected degree & The aggregate spectrum bounds full-Jacobian
singularity.  Every possible family-count pattern is swept at Levels I/V; the
complete square residual system is used at Level III.  All soft-$O$ and
implementation overhead remains the explicit multiplier $\kappa_{\rm hom}$.\\
Generic-MQ fallback & Model input & The same exact quotient system, but the
pinned generic semi-regularity and linear-algebra estimator.\\
Circuit-unit conversion & Sensitivity calculation & Only the field arithmetic
charged by the estimator; not a complete attack circuit or category claim.\\
\bottomrule
\end{tabularx}
\end{table}

\paragraph{The attack as a data-flow.}
The proof is easiest to follow as a sequence of exact transformations.  The
word \emph{lift} below means the attacker-chosen affine rule that turns one
column vector into a full signature candidate.

\begin{center}
\renewcommand{\arraystretch}{1.25}
\begin{tabularx}{0.96\textwidth}{@{}>{\bfseries\raggedright\arraybackslash}p{0.22\textwidth}>{\centering\arraybackslash}p{0.04\textwidth}X@{}}
Full public verifier & $\longrightarrow$ & impose one affine common-column
relation; quadratic, linear, and constant terms become explicit\\
& $\downarrow$ & \\
Symmetry quotient & $\longrightarrow$ & identify $(\xi,\eta)$ with
$(\eta,\xi)$ and keep only the distinct quadratic features\\
& $\downarrow$ & \\
Exact affine elimination & $\longrightarrow$ & use a left-kernel projection
to expose every output direction unreachable by the quadratics as an ordinary
linear constraint\\
& $\downarrow$ & \\
Structured residual MQ & $\longrightarrow$ & choose either a stable
$A$-linear slice or a full-dimensional ordinary slice for $\ell=4$, with
exact rejection-density control, or a rejection-compatible relation for
$\ell=2$\\
& $\downarrow$ & \\
Complementary solvers & $\longrightarrow$ & stream the multigraded Macaulay
operator; solve a Frobenius-channel core; or sweep every square
multidegree-family projection needed to cover a full-Jacobian root\\
& $\downarrow$ & \\
Candidate signature & $\longrightarrow$ & map back to the original
coordinates and run the unmodified public verifier
\end{tabularx}
\end{center}

Every transformation through the residual MQ system is exact and uses only
public algebra and linear algebra.  The solver model enters only when we
estimate the cost of solving that residual system.  The final public-verifier
check makes the attack Las Vegas: a trial may be discarded, but an invalid
signature is never returned.

\paragraph{Four jobs that a complete attack must perform.}
The equality \eqref{eq:intro-swap} is only the first job.  The construction
also has to preserve all public outputs, ensure that targets have roots, and
make the residual equations compatible with a practical solver:
\begin{enumerate}[label=\textbf{\arabic*.},itemsep=3pt]
\item \textbf{Constrain the signature columns.}  Every signature column is
      written as an affine function of one common vector.  Substitution
      separates the verifier into quotient quadratic features, offset-linear
      terms, and constants.
\item \textbf{Eliminate the affine part without dropping equations.}  Output
      directions outside the quadratic image become linear constraints.  For
      each target satisfying those affine consistency conditions, Gaussian
      elimination leaves a smaller MQ system with exactly the same solutions
      as the constrained verifier.
\item \textbf{Choose a solver-friendly search space.}  For $\ell=4$, a
      structured choice of offsets leaves a large $A$-linear kernel and an
      injective local skew map makes rejection cost less than $0.0014$ bits.  Over a
      splitting field this kernel becomes four variable blocks, preserving the
      multidegrees used by special XL.\@  The same slice also has an exact
      Frobenius-channel form with two $A$-valued channels of degrees two and
      twenty.  For $\ell=2$, fixed-skew and filtered relations handle the
      verifier's rejection rule directly.
\item \textbf{Retry, solve, and verify.}  The low-state branch uses streamed
      multigraded linear algebra.  Two independent branches use symbolic
      homotopy: one on a Frobenius-channel core, and one on a complete family of
      square projections that is guaranteed to contain a nonsingular core
      whenever the full restricted Jacobian has rank $h$.  Exact moment bounds
      control accepted nonsingular roots directly, without requiring a unique
      root.  Every recovered candidate is transformed back and checked by the
      original verifier.
\end{enumerate}

\paragraph{A complete dimension walk-through.}
The Level-I parameter set $(v,o,\ell,r)=(28,5,4,4)$ illustrates where the
reduction comes from: it has 28 vinegar blocks, 5 oil blocks, and blocks with
4 rows and 4 columns.  Here $n=v+o=33$, one stacked signature column has
$N_0=n\ell=132$ coordinates, the verifier has $M=or\ell=80$ scalar outputs,
and there are $m_1=5$ public base-form groups.

\begin{table}[H]
\centering
\caption{The dimensions in the $(28,5,4,4)$ attack.  This table is only a
worked example; the general formulas appear in later sections.}
\label{tab:running-example}
\small
\renewcommand{\arraystretch}{1.12}
\begin{tabularx}{0.96\textwidth}{@{}>{\raggedright\arraybackslash}p{0.28\textwidth}>{\centering\arraybackslash}p{0.18\textwidth}X@{}}
\toprule
Stage & Dimension & Meaning\\
\midrule
Raw public feature vector & $N_{\rm raw}=1{,}280$ & $m_1\ell^2r^2$ labeled slots before
the column relation\\
After the common-column relation & $80$ ordered quadratics & one feature for
each of $m_1\ell^2=5\cdot16$ ordered power pairs\\
After the symmetry quotient & $K=50$ distinct quadratics &
$m_1\binom{\ell+1}{2}=5\cdot10$ unordered power pairs\\
Affine output directions & $M-K=30$ & linear equations that determine an
affine search space\\
Full residual domain & $102$ coordinates & $N_0-(M-K)=132-30$ after affine
elimination\\
Stable kernel & $52$ base-field dimensions & an $A$-linear space of
$A$-dimension $13$ on which offset-linear motion vanishes\\
Selected special-XL slice & $h=40$ & four blocks of $s=10$ variables; hence
$\tau=K-h=10$\\
\bottomrule
\end{tabularx}
\end{table}

The final slice has fewer variables than equations.  Full cross-polar rank is
not needed; the condition that matters is the aggregate collision mass
$\bar\eta_L\le19^{-1}$.  The minimum-rank certificate $\rank T_y\ge h+1=41$ for
every nonzero slice direction implies that bound.  Either certificate makes a random target--slice trial contain a
unique base-field root with probability at least
$19^{-10}(1-19^{-1})$, so the expected retry count is at most
$19^{10}\cdot19/18$.  The factor $19/18$ costs only $0.0781$ bits.  Paying
these retries is worthwhile because reducing the Macaulay matrix from the full
residual domain to four blocks of ten variables saves much more.  This is why
the paper analyzes both the algebraic reduction and the root-retry probability,
rather than quoting an MQ dimension alone.

\paragraph{A fixed-core-free selected-degree-free frontier.}
The aggregate spectrum also bounds the density at which the full restricted
Jacobian loses rank.  At a descended root, $19$-Frobenius cyclically permutes
the four eigenblocks and carries every nonzero Jacobian minor to a nonzero
conjugate minor.  The $55$ near-square family-count patterns therefore collapse
to $16$ Frobenius orbits, and one representative per orbit suffices.  This
removes both a production selected-degree premise and any preferred square-core
Jacobian premise.

\begin{table}[H]
\centering
\caption{Fixed-core-free, selected-degree-free homotopy frontier.  ``Per key''
charges the robust aggregate-spectrum and acceptance bounds.  ``Normalized''
also charges the inverse of the exact random-XOF structural-preflight lower
bound $0.9423222679$.  Headroom is measured against $143/207/272$.}
\label{tab:selected-homotopy-frontier}
\small
\setlength{\tabcolsep}{4.0pt}
\begin{tabular}{@{}c c c r r r@{}}
\toprule
Level & $(h,K)$ & patterns $\to$ orbits & per key & normalized & headroom\\
\midrule
I   & $(48,50)$ & $55\to16$ & $138.324$ & $138.410$ & $4.590$\\
III & $(68,70)$ & $55\to16$ & $179.758$ & $179.844$ & $27.156$\\
V   & $(88,90)$ & $55\to16$ & $220.832$ & $220.918$ & $51.082$\\
\bottomrule
\end{tabular}
\end{table}

The artifact enumerates all $55$ patterns, computes their $16$ exact cyclic
orbits, and sums the exact multihomogeneous B\'ezout and homotopy ledgers for
one representative from each orbit.  Every unused equation and every final
verification check is charged.  The straightforward largest one-core output
ceilings are about $1.890$~PiB, $2.076$~ZiB, and $2416.552$~YiB; hence this
route is included to remove logical assumptions, not as a resident-memory
claim.  The separately conditioned eigenblock-core and streamed-XL branches
provide lower-output and lower-state alternatives.

\paragraph{Conditional concrete estimates.}
\Cref{tab:main} reports the base-2 expected-cost exponents produced by the
\emph{pinned} solver models.  ``Same-quotient generic'' applies the pinned
generic-MQ estimator to exactly the same symmetry-quotient residual system;
``gain'' is therefore a like-for-like solver-structure comparison, not a
comparison with the submission's different gate methodology.  Every repeated
solve includes its explicit target or target--slice retry factor.

\begin{table}[H]
\centering
\caption{Primary conditional estimates for the public Version~2.3 parameter
sets.  The tuple is $(v,o,\ell,r)$.  All numerical columns are base-2 expected
cost exponents or differences of such exponents in the cited solver model.
The exact reductions determine the systems; the rank-spectrum and production
solver conditions are listed immediately below.}
\label{tab:main}
\footnotesize
\setlength{\tabcolsep}{3.2pt}
\renewcommand{\arraystretch}{1.08}
\begin{tabular}{@{}c l l r r r@{}}
\toprule
Level & Parameters & Attack mode & Model & \shortstack{same-quotient\\generic} & Gain\\
\midrule
I   & $(28,5,4,4)$  & structured XL & 128.90 & 136.11 &  7.21\\
I   & $(48,16,2,2)$ & fixed-skew MQ & 129.51 & 129.51 &  0.00\\
I   & $(28,4,4,5)$  & structured XL & 128.90 & 136.11 &  7.21\\
\midrule
III & $(40,7,4,4)$  & structured XL & 171.76 & 180.66 &  8.90\\
III & $(72,24,2,2)$ & fixed-skew MQ & 183.51 & 183.51 &  0.00\\
III & $(38,5,4,5)$  & structured XL & 171.76 & 180.66 &  8.90\\
\midrule
V   & $(50,9,4,4)$  & structured XL & 214.20 & 224.75 & 10.55\\
V   & $(96,32,2,2)$ & filtered MQ   & 235.27 & 237.92 &  2.65\\
V   & $(52,6,4,6)$  & structured XL & 214.20 & 224.75 & 10.55\\
\bottomrule
\end{tabular}
\end{table}

The six $\ell=4$ rows stay below the numerical values $143$, $207$, and
$272$ even after special XL is discarded and the exact quotient systems are
fed to the pinned generic-MQ model.  This is a model comparison: NIST evaluates
categories across all the security metrics it regards as relevant, not by one
arithmetic exponent alone~\cite{NISTPQCCriteria}.

Space is not hidden inside those exponents.  Minimizing arithmetic alone would
demand one-vector special-XL states of $65.4$~GiB, $118$~PiB, and $614$~EiB;
the last two are physically implausible, so the body instead prices low-state
alternatives (\Cref{sec:special-xl,sec:three-plus-one}).  Their best low-state
rows cost $2^{141.29}$, $2^{200.49}$, and $2^{263.40}$ AXN operations in packed
states of $1.53$~TiB, $55.4$~GiB, and $9.55$~GiB, with Level~I enlarging only
the preconditioning field to $\F_{19^{12}}$ and charging arithmetic over
$\F_{19^4}$.

Two further routes reach the same conclusion without the selected-degree
assumption that special XL uses.  The three-plus-one decomposition
(\Cref{sec:three-plus-one}) drops one block from the nonlinear solve; the
Frobenius-channel route (\Cref{sec:channel-homotopy}), built from the normal
form
$\Sym^2_{\F_{19}}(\F_{19^4})\cong\F_{19^4}\oplus\F_{19^4}\oplus\F_{19^2}$,
solves a square mix of two channels and filters the rest.  Each is derived with
its own exponents and memory envelopes in the body, and each stays below the
$143/207/272$ references.  All of these are proved operator-and-accounting
statements, not production timings.

\paragraph{Conditions behind the table.}
For the structured rows, the random-XOF theorem of
\Cref{thm:seeded-spectrum} gives $\bar\eta_L\le19^{-1}$ except with probability
below $2^{-55.22}$, $2^{-123.19}$, and $2^{-191.16}$ for the compact profiles,
provided the reserved-line slice construction is used.  This is an
information-theoretic theorem in the random-XOF idealization of the official
expansion, not an independence theorem for the deterministic bytes produced by
a fixed SHAKE seed.  Instantiating the XOF by SHAKE is the usual computational
modeling step.  A fixed key can instead publish the exact projective rank
histogram.  The pointwise certificates $\rank T_y\ge h+1$ are stronger
sufficient conditions, not requirements.

At the selected special-XL degree, the compact scalar-Wiedemann rows retain the
published one-dimensional-kernel prediction.  The robust explicit rows charge
only $\sigma\ge1/8$ in the Las Vegas wrapper and allow
$\bar\eta_L\le1/2$.  A certified finite border quotient of larger dimension
still gives zero-error extraction by \Cref{thm:finite-quotient-recovery}, but
its additional block/kernel work is not free and must fit the stated work and
state envelope.  For the three-plus-one rows, $\sigma$ jointly covers a complete finite
quotient for the three-block subsystem, its extraction and enumeration within
the stated work--state reserve, and certified linear completion of the fourth
block.  A quotient or completion that exceeds the reserve, or a consistent
rank-deficient completion, causes a restart unless handled by a separately
priced fallback.

The channel-homotopy rows do not use this selected-degree assumption.  They
instead require one exact nonzero leading-Jacobian witness for the chosen
public slice and core.  Their arithmetic theorem has a soft-$O$ factor, so the
paper reports the normalized leading term together with the full break-even
budget $\kappa_{\rm hom}$ and a dense coefficient-volume ceiling; it does
not promote that term to a production timing.

The basis-complete and complete-square rows also avoid the selected-degree
assumption, and they do not require a fixed-core Jacobian witness.  Their root
success follows from the same aggregate spectrum used for collision control.
Their displayed costs sum every family-count core (or use the complete square
system), charge candidate filtering, and leave the soft-$O$/implementation
factor $\kappa_{\rm hom}$ explicit.  Their principal caveat is output volume,
not logical root completeness.

The fixed-skew rows assume $\Theta_{\rm pol}\le19^{-1}$, and the filtered
Level-V row assumes $\Psi_{\rm pol}\le19^{-1}$; each costs less than $0.075$
bits.  Minimum polar ranks $49$, $73$, and $194$ are simple sufficient
certificates for the three headline rows, not requirements.
The stronger thresholds $96$, $144$, and $240$ make every target solvable.
The released finite sampling supports these conditions but does not certify
the complete spectra.  Likewise, the small-profile Macaulay experiments do
not certify the selected production degree.

The later AXN and NAND conversions change only the arithmetic unit charged by
the estimators.  They are sensitivity calculations, not a complete gate-level
attack cost.

\paragraph{How to read the paper.}
\Cref{sec:mq-background} fixes the background and notation.  The exact attack is
in \Cref{sec:prelim,sec:symmetry,sec:structured-lift}.  Root existence and
retry probabilities are separated in \Cref{sec:slices}; the solver and gate
models are isolated in \Cref{sec:complexity}, including the independent
Frobenius-channel route of \Cref{sec:channel-homotopy}; and
\Cref{sec:evidence} states exactly which assumptions are tested rather than
proved.  All long proofs and exhaustive cost-search certificates are deferred
to the appendices.

The submitted $q=16$ matrices are not symmetrized and therefore do not satisfy
the hypothesis of our quotient theorem.  Returning unchanged to that design,
however, would restore the known characteristic-two attack surface discussed
in \Cref{sec:related}.  A future odd-characteristic revision that removes
public-matrix symmetry, changes the left/right actions, or changes the affine
geometry falls outside the theorem until its new public forms are analyzed.

\section{Mathematical Background and Relation to Prior Work}\label{sec:mq-background}

This section fixes the notation and background used throughout: underdetermined
MQ, polar forms, Macaulay matrices, and multigraded XL.

\subsection{MQ systems, roots, and affine slices}\label{sec:mq-primer}

Write $\F_q$ for the finite field with $q$ elements.  An MQ system over
$\F_q$ is a vector-valued polynomial map of total degree at most two,
\[
  g=(g_1,\ldots,g_K):\F_q^N\longrightarrow\F_q^K,
\]
together with a target $z\in\F_q^K$.  We continue to call $g$ quadratic when
some coordinates also contain linear or constant terms.  A \emph{root} is an
$x$ with $g(x)=z$.
The set
\[
  g^{-1}(z)=\{x:g(x)=z\}
\]
is the \emph{fiber} above $z$: it is the set of all candidate assignments
that produce the target.  After the affine reconstruction described later,
those assignments become full signature candidates.  All arithmetic is over
the stated finite field; every concrete \SNOVA{} calculation in this paper uses
$q=19$, so additions and multiplications are performed modulo $19$.

As a toy example, two equations such as
\[
 x_1x_2+x_3^2=z_1,
 \qquad
 x_1^2+x_2x_3+x_1=z_2
\]
form an MQ system in three variables.  Having more variables than equations
makes the system \emph{underdetermined}, but does not make it linear or
necessarily easy.  Modern MQ attacks often spend most of their time on very
large linear-algebra matrices built from the quadratic equations.

An \emph{affine space} is a translate $a+V=\{a+v:v\in V\}$ of a linear
subspace.  An \emph{affine slice} $a+L\subseteq a+V$ restricts the search to a
smaller direction space $L$.  If $h=\dim L$, choosing a basis of $L$ turns the
slice into an MQ system in $h$ coordinates.  Smaller $h$ makes the algebraic
solve cheaper, but it also lowers the chance that the slice contains a root;
\Cref{sec:slices} quantifies this tradeoff.

Three logically distinct events recur throughout the paper:
\begin{enumerate}[label=(\roman*),itemsep=2pt]
\item the chosen target and slice contain at least one $\F_q$-rational root;
\item the slice contains exactly one such root; and
\item the selected Macaulay degree exposes enough ideal relations to build a
      finite quotient algebra and extract its roots.
\end{enumerate}
The first two are properties of the quadratic map.  The third is a statement
about a specific solver representation.  A proof of root existence is not, by
itself, a proof that one solver invocation reaches a complete enough degree.
Conversely, the quotient algebra need not be one-dimensional: once its
multiplication tables are available, standard zero-dimensional solving can
enumerate its roots and the public verifier can select a valid one.

\subsection{Rank, polar forms, and collisions}

For a scalar quadratic polynomial $f$, define its polar form by
\[
  \beta_f(x,y)=f(x+y)-f(x)-f(y)+f(0).
\]
In odd characteristic, $\beta_f$ is bilinear.  It is the finite-field analogue
of the part of a directional derivative that comes from the quadratic terms.
Its matrix rank measures how many independent directions the quadratic part
couples.

For a vector-valued map $g=(g_1,\ldots,g_K)$, every nonzero coefficient vector
$b\in\F_q^K$ selects a scalar combination $b\cdot g$.  The minimum polar rank
must be taken over \emph{all} such output directions.  Multiplying $b$ by a nonzero scalar does not change the rank, so these
vectors can be grouped into projective directions: one-dimensional lines
through the origin.

Rank matters because finite-field Fourier (character-sum) identities imply
that a scalar quadratic with high polar rank is well balanced.  Applied to all
nonzero combinations $b\cdot g$, this bounds every target fiber rigorously:
when the ranks are high enough relative to $K$, each fiber has size close to
the random-map prediction $q^{N-K}$.  A related \emph{cross-polar rank}
measures how strongly a displacement inside a slice interacts with the full
residual domain; high cross-polar rank suppresses pairs of distinct slice
points that map to the same target.  In summary:
\begin{center}
\begin{tabularx}{0.91\textwidth}{@{}>{\bfseries}p{0.28\textwidth}X@{}}
Polar rank & controls whether full-space target fibers are nonempty and close
to their expected size.\\
Cross-polar rank & controls collisions inside an affine slice and therefore
unique-root retry probabilities.\\
Macaulay quotient & commuting border tables certify that row reduction
exposes a finite quotient algebra; its dimension controls the cost of root
extraction.
\end{tabularx}
\end{center}
High rank here is not a ``randomness'' heuristic: \Cref{sec:slices} gives exact
first- and second-moment formulas.  The production task is to certify the
relevant weighted rank spectrum.  A lower bound in every nonzero direction is a
stronger sufficient route, but the exact moment formulas tolerate a small
exceptional locus.  Public sampling gives evidence; an exact spectrum bound or
a proof is still needed.

\subsection{XL and Macaulay matrices}\label{sec:xl-primer}

XL (``eXtended Linearization'') turns a nonlinear system into a larger linear
system.  Starting from quadratic equations $f_i=0$, it multiplies each $f_i$
by selected monomials.  For example,
\[
  f=x_1x_2+x_3^2+x_1+1
  \quad\Longrightarrow\quad
  x_1f=x_1^2x_2+x_1x_3^2+x_1^2+x_1.
\]
Each resulting polynomial becomes a row.  Each monomial---for example
$x_1^2x_2$ or $x_1x_3^2$---becomes a column.  The coefficient matrix is a
\emph{Macaulay matrix}.  Gaussian elimination on this matrix can reveal enough
relations among the monomials to recover a root.  Treating monomials as columns
is only a linear-algebra device: the monomials are not independent variables in
the original problem, and the solver must eventually recover coordinates that
satisfy all of their multiplicative relations.

The cost comes from the size and rank of this matrix, not from evaluating the
original equations once.  Generic XL tracks only total degree.\@  \emph{Special
XL} exploits a variable-block structure.  A monomial then has a vector of
degrees, one entry per block, called its \emph{multidegree}.  Building only
selected multidegrees can make the Macaulay matrix much smaller.

Our $\ell=4$ reduction produces four blocks after a change to eigenblock
coordinates.  A quadratic using one block twice has degree pattern
$(2,0,0,0)$ up to permutation; a bilinear equation using two blocks has pattern
$(1,1,0,0)$.  The stable-lift construction in
\Cref{sec:structured-lift} exists precisely to preserve these patterns after
affine elimination.  A \emph{homogenizer} is one extra variable per block that
allows the constant and linear terms to be written with the same
multidegree; it is set to one when the root is read out.

If $x$ is a root, the vector of all selected monomial values
\[
  \nu(x)=(1,x_1,\ldots,x_N,x_1^2,x_1x_2,\ldots)
\]
lies in the right kernel of the dehomogenized Macaulay matrix.  A
one-dimensional right kernel is the easiest case: normalizing the constant
coordinate reveals the root immediately.  Larger kernels are also recoverable.
If row reduction instead expresses every border monomial $x_jm$ in terms of a
finite monomial basis $\mathcal B$, those relations give the multiplication
matrices of the quotient algebra.  The FGLM algorithm converts these tables to
a lexicographic Gr\"obner basis and extracts every algebraic
root~\cite{FaugereEtAl1993FGLM}, so a commuting-border certificate gives
publicly checkable correctness beyond exact Macaulay corank one.  The scalar-Wiedemann exponent reported later
still prices the published one-dimensional-kernel case; a larger quotient
requires separately charged block/nullspace work.

\subsection{Relation to prior attacks}\label{sec:related}

Beullens rewrote the \SNOVA{} verifier in ordered bilinear power-pair
coordinates and derived affine-column forgery attacks from rank defects in the
resulting outer matrix~\cite{Beullens2024SNOVA}.  Our quotient acts one layer
earlier: public-matrix symmetry identifies the power-pair polynomials
\emph{before} the outer map sees them.

The danger of odd-characteristic symmetrization had been flagged before.  In a
2025 security presentation, the \SNOVA{} team based its odd-$q$ discussion on
the $o\binom{\ell}{2}$ alternating-form count, observed that a wedge-style
variant might apply, and said that further cryptanalysis was needed
\cite{WangEtAl2025Round2Security}.  NIST subsequently described the symmetric
odd-characteristic construction as a proposed response rather than a stable
third-round design~\cite{NISTIR8610}.  Version~2.3 also includes a public
block-symmetry rejection check~\cite{SNOVA23}.  We therefore do not claim
novelty for the qualitative warning that symmetrization can create an attack
surface or for the rough alternating-form count.


\begin{table}[H]
\centering
\caption{Boundary with the closest prior analyses.  The distinction is by
public construction and cryptanalytic deliverable, not merely terminology.}
\label{tab:prior-art-boundary}
\small
\renewcommand{\arraystretch}{1.13}
\begin{tabularx}{0.97\textwidth}{@{}>{\raggedright\arraybackslash}p{0.20\textwidth}>{\raggedright\arraybackslash}p{0.23\textwidth}X@{}}
\toprule
Work & Setting & What is established\\
\midrule
SNOVA team, 2025~\cite{WangEtAl2025Round2Security} & Proposed odd-$q$ response &
Qualitative warning that a wedge-style variant may work, the
$o\binom{\ell}{2}$ alternating-form sizing criterion, and rejection of
symmetric signature blocks; no instantiated direct forgery or concrete cost.\\
Bros et al., 2026~\cite{BrosEtAl2026SNOVAWedge} & Submitted construction over
$q=2^k$ & Characteristic-two wedge attack breaking the updated Round-2
parameters; not an analysis of the $q=19$ symmetric redesign.\\
This work & Version~2.3 over $q=19$ & Exact $\Sym^2(A)$ affine-column
factorization, complete public-key forgery reductions, retry accounting,
space-aware costs, and constructive bypass of the proposed rejection rule.\\
\bottomrule
\end{tabularx}
\end{table}

The novelty is the missing cryptanalytic instantiation: an exact
symmetric-square factorization inside the affine-column attack; complete
residual systems for every public $q=19$ shape; exact root and retry accounting;
an $A$-linear lift preserving the special-XL grading; and constructive
fixed-skew and filtered bypasses of Version~2.3's block-symmetry check.  The
presentation identified a possible attack surface but gave neither this
forgery reduction nor concrete costs.  This paper supplies both and separately
proves that the deployed rejection condition does not remove the reduction.

Cabarcas et al.\ exploit stability under the public matrix algebra with a
multigraded XL algorithm~\cite{Cabarcas2025SNOVA}.  Furue et al.\ expose related
multihomogeneous systems through matrix transformations~\cite{FurueEtAl2025}.  Our attack starts from the complete affine forgery
system, uses a public column lift to retain every output equation, and then
restricts to an invariant kernel on which the affine offsets preserve the
block grading needed by special XL.\@

The 2026 wedge attack is already a full break of the updated Round-2
parameters, but its algebra is explicitly characteristic two: it works over
$q=2^k$ and exploits the alternating polar form of the block-ring system
\cite{BrosEtAl2026SNOVAWedge}.  Later work studies that wedge map further
\cite{TsengEtAl2026Wedge}.  Our target is the proposed odd-characteristic
response.  Its public matrices are symmetric rather than alternating, and our
attack factors the public affine-column feature map through $\Sym^2(A)$; it is
not an application of the characteristic-two wedge kernel.  Other analyses
reinterpret scalar-expanded \SNOVA{} as Unbalanced Oil and Vinegar (UOV), or
lift it to extension-field UOV, for key recovery~%
\cite{IkematsuAkiyama2024,LiDing2024,NakamuraEtAl2025}.  Symmetric and
exterior algebras also appear in recent UOV key-recovery attacks~%
\cite{JinEtAl2026,Ran2026}.  Our attack is a direct forgery attack: it does not
recover the hidden oil subspace that constitutes the UOV-style secret key.

For the generic $\ell=2$ cores, we reproduce Hashimoto's optimization,
including its published boundary value $a=1$~%
\cite{Hashimoto2023,Beullens2024SNOVA}.  When this leaves the elementary system
$\operatorname{MQ}(q,1,1)$, the artifact solves it by exhaustive evaluation and
charges that work conservatively.  May, Ostuzzi, and Ressler's \emph{Just
Guess} algorithm and the generalized variable partitioning of Asanuma et al.\ are more recent approaches to underdetermined MQ~%
\cite{MayOstuzziRessler2026,AsanumaEtAl2026}.  Evaluating the released
classical Just-Guess formulas on our residual dimensions gives larger estimates
than the selected Beullens--Hashimoto modes.  We have not implemented the newer
generalized partition search.  The ``generic MQ'' numbers below are therefore
pinned, reproducible baselines, not a claim to have exhausted every generic
solver.

Solver engineering is also moving quickly.  Sakata and Takagi use
Hilbert-series information to reduce the matrices built by an $F_4$
Gr\"obner-basis solver and report a new MQ-challenge
record~\cite{SakataTakagi2026}.  That is an implementation and challenge
result, not a plug-in cost formula for the underdetermined residual systems
here, so we do not substitute it into \Cref{tab:main}.


\subsection{Notation at a glance}\label{sec:notation}

\begin{table}[H]
\centering
\caption{Recurring notation.  Generic background notation such as $N$ is local;
$N_0$ is the specific common-column dimension in the \SNOVA{} reduction.}
\label{tab:notation}
\small
\renewcommand{\arraystretch}{1.12}
\begin{tabularx}{0.96\textwidth}{@{}>{\raggedright\arraybackslash}p{0.23\textwidth}X@{}}
\toprule
Symbol & Meaning\\
\midrule
$(v,o,\ell,r)$ & Public parameter tuple: $v$ vinegar blocks, $o$ oil blocks,
block height $\ell$, and $r$ columns per signature block.\\
$n=v+o$ & Number of matrix-valued signature blocks.\\
$M=or\ell$ & Number of scalar public output equations.\\
$N$ & Generic domain dimension for an MQ map $g:\F_q^N\to\F_q^K$; its
meaning is local to the section in which the map is defined.\\
$N_0=n\ell$ & Length of one stacked signature column after the common-column
relation.\\
$N_{\rm raw}=m_1\ell^2r^2$ & Number of labeled verifier feature slots before imposing
the relation.\\
$A=\F_q[S]$ & Public commutative matrix algebra generated by $S$; in the
Version~2.3 rows, $\dim_{\F_q}A=\ell$.\\
$m_1$ & Number of public base-form groups in the pinned implementation.\\
$K=m_1\binom{\ell+1}{2}$ & Maximum number of distinct quadratic power-pair
features after the symmetry quotient.\\
$Q_R$ & Matrix mapping the $K$ quotient features to public output directions.\\
$G_R$ & A public left inverse of $Q_R$ when $Q_R$ has full column rank.\\
$C_{R,\mathbf v}$ & Linear constraint matrix obtained by projecting away the quadratic
image.\\
$g_{R,\mathbf v}$ & Canonical $K$-coordinate residual map after the exact output split.\\
$H_{R,\mathbf v}$ & Stable $A$-linear kernel used to choose structured solver slices.\\
$L$, $h$ & Slice direction space and its base-field dimension.\\
$\tau=K-h$ & Slice deficit.  It enters the exact full-rank retry formula;
for the selected $\tau=10,14$ slices, the retry count is essentially $q^\tau$.\\
$r_*$ & Minimum full-space polar rank over all nonzero scalar output
combinations.\\
$e_L(b)$ & Cross-polar rank of output direction $b$ relative to slice $L$.\\
$\mu$, $\bar\eta$ & Expected roots on a random target--slice pair and an upper
bound on its collision mass.\\
\bottomrule
\end{tabularx}
\end{table}

\section{The Affine-Column Framework}\label{sec:prelim}

This section rewrites the public verifier after an attacker-imposed affine
relation among signature columns.  The result has three visibly separate
parts: quotient quadratic features, linear terms created by the offsets, and a
constant term.  \Cref{prop:affine-reduction} then eliminates the affine part
exactly.

Three vector spaces appear and should not be confused.  The raw public feature
vector has length $N_{\rm raw}$; the verifier outputs $M$ scalar equations; and the
common column variable $u$ has length $N_0$.  The matrices below simply move
between these three spaces.

\subsection{Parameters and public forms}\label{sec:parameters}

The public parameter tuple is $(v,o,\ell,r)$.  There are $v$ vinegar blocks and
$o$ oil blocks, for a total of
\[
  n=v+o
\]
matrix-valued signature blocks.  The names ``vinegar'' and ``oil'' come from
SNOVA's UOV ancestry; in this public forgery attack they serve only to specify
dimensions, and the attacker never learns or reconstructs the secret partition.
Each block has $\ell$ rows and $r$ columns.
Stacking the $j$th column of all $n$ blocks gives a vector of length
$N_0=n\ell$; the verifier has $M=or\ell$ scalar outputs.  The pinned
Version~2.3 implementation defines
\begin{equation}\label{eq:parameters-basic}
 m_1=\left\lceil\frac{or}{\ell}\right\rceil,
 \qquad M=or\ell,
 \qquad N_0=n\ell.
\end{equation}
Here $m_1$ is the number of public base-form groups used to assemble the
outputs.  The ceiling is used by the reference implementation; a public
analysis script uses a floor and therefore disagrees when $\ell\nmid or$.  All
dimensions in this paper follow the pinned implementation~\cite{SNOVA23}.

Let $S\in\Mat_\ell(\F_q)$ be the public symmetric matrix with irreducible
characteristic polynomial.  Its powers generate the commutative matrix
algebra
\[
 A=\F_q[S],\qquad \Sigma_\xi=I_n\otimes\xi\quad(\xi\in A).
\]
For the Version~2.3 parameters, $A\cong\F_{q^\ell}$ and
$\dim_{\F_q}A=\ell$.  Multiplication by $\Sigma_\xi$ applies the same
$\ell\times\ell$ matrix $\xi$ independently to every signature block.

The scalar-expanded public base matrices are
$P_1,\ldots,P_{m_1}\in\Mat_{N_0}(\F_q)$.  In the public $q=19$ construction,
each $P_i$ is symmetric.  For $\xi,\eta\in A$, define the bilinear form
\begin{equation}\label{eq:bilinear-form}
 B_i^{\xi,\eta}(x,y)=x^\transpose\Sigma_\xi P_i\Sigma_\eta y.
\end{equation}
Setting $x=y$ makes this a quadratic feature.  The power basis
$1,S,\ldots,S^{\ell-1}$ gives the ordered power-pair coordinates used in the
prior affine-column attack.

\subsection{Imposing one common column variable}\label{sec:affine-columns}

Write the stacked signature columns as
\[
  U=[u_0|\cdots|u_{r-1}],\qquad u_j\in\F_q^{N_0}.
\]
The attacker chooses multipliers $R_j\in A$ and offsets
$v_j\in\F_q^{N_0}$, with $R_0=1$.  Write
$R=(R_0,\ldots,R_{r-1})$ and
$\mathbf v=(v_0,\ldots,v_{r-1})$, and impose
\begin{equation}\label{eq:relation}
 u_j=\Sigma_{R_j}u+v_j,
 \qquad 0\le j<r.
\end{equation}
Thus all $r$ columns are affine functions of one vector $u$.  This defines a
publicly checkable subset in which the attacker searches for a forgery, with no
assumption about honest signatures.

Before the relation is imposed, the feature vector has
\[
  N_{\rm raw}=m_1\ell^2r^2
\]
coordinates: $m_1$ base forms, $\ell^2$ ordered power pairs, and $r^2$ ordered
column pairs.  The official verifier flattens each $r\times\ell$ block in
$(i,j)$ order, whereas our feature formulas use the block-transposed $(j,i)$
order.  Let $t_{\rm off}\in\F_q^M$ and
$\mathcal E_{\rm off}\in\F_q^{M\times N_{\rm raw}}$ denote the target and public outer
map in the official order.  If $\Pi_{\rm out}$ is the corresponding public
permutation, define
\[
  t=\Pi_{\rm out}t_{\rm off},\qquad
  \mathcal E=\Pi_{\rm out}\mathcal E_{\rm off}.
\]
This is only a coordinate reordering; it changes neither ranks nor solutions.

After substituting \eqref{eq:relation}, products of two variable parts give
quadratic terms in $u$, products of one variable part and one offset give
linear terms, and products of two offsets give constants.  For a fixed
relation $R=(R_j)_{j=0}^{r-1}$, let
$\mathcal R_R\in\F_q^{N_{\rm raw}\times m_1\ell^2}$ expand the remaining ordered
quadratic features into the full feature vector, and let
$\mathcal F_{R,\mathbf v}\in\F_q^{N_{\rm raw}\times N_0}$ collect the offset-linear
coefficients.  The complete substituted verifier is
\begin{equation}\label{eq:ordered-system}
 \mathcal E\bigl(\mathcal R_Rz_{\rm ord}(u)+\mathcal F_{R,\mathbf v}u\bigr)
 +c_{R,\mathbf v}=t,
\end{equation}
where $z_{\rm ord}(u)$ contains the $m_1\ell^2$ values
$B_i^{S^a,S^b}(u,u)$ and $c_{R,\mathbf v}$ is the public constant term.  Put
\begin{equation}\label{eq:ER}
 E_R=\mathcal E\mathcal R_R\in\F_q^{M\times m_1\ell^2},
 \qquad \Lambda_{R,\mathbf v}=\mathcal E\mathcal F_{R,\mathbf v}\in\F_q^{M\times N_0}.
\end{equation}

\subsection{Separating quadratic and affine output directions}

By \Cref{thm:symmetric-square}, the ordered feature vector duplicates every
off-diagonal power pair.  Let $z_{\rm sym}(u)$ contain one coordinate for each
unordered pair $0\le a\le b<\ell$, and let $D$ copy each off-diagonal
coordinate into its two ordered positions.  Then
$z_{\rm ord}(u)=Dz_{\rm sym}(u)$.  Define
\begin{equation}\label{eq:QR}
 Q_R=E_RD\in\F_q^{M\times K},
 \qquad K=m_1\binom{\ell+1}{2}.
\end{equation}
The columns of $Q_R$ span all output directions that the residual quadratic
features can reach.

Now choose $W_R$ whose rows form a basis of the left kernel of $Q_R$ and put
\begin{equation}\label{eq:C}
 C_{R,\mathbf v}=W_R\Lambda_{R,\mathbf v}.
\end{equation}
Multiplying by $W_R$ kills every quadratic feature.  What remains is therefore
an ordinary affine system in $u$.  This elementary projection is the key to
retaining the \emph{complete} verifier rather than dropping output equations
that do not lie in the quadratic image.

\begin{center}
\begin{tabularx}{0.92\textwidth}{@{}>{\bfseries}p{0.18\textwidth}X@{}}
$Q_R$ & sends the $K$ distinct quadratic features to the $M$ public output
coordinates; its column space is the quadratic image.\\
$W_R$ & keeps precisely the output directions orthogonal to that image, so
$W_R Q_R=0$.\\
$C_{R,\mathbf v}=W_R\Lambda_{R,\mathbf v}$ & records the ordinary linear equations that remain in
those output directions after the offsets are introduced.
\end{tabularx}
\end{center}

\begin{lemma}[Canonical residual coordinates]
\label{lem:canonical-residual-coordinates}
Assume $\rank Q_R=K$.  Choose a public left inverse
$G_R\in\F_q^{K\times M}$ satisfying $G_RQ_R=I_K$.  Then
\begin{equation}\label{eq:stacked-output-change}
 \begin{bmatrix}G_R\\ W_R\end{bmatrix}\in\F_q^{M\times M}
\end{equation}
is invertible.  Writing $b=t-c_{R,\mathbf v}$, the complete substituted
verifier is equivalent to the pair
\begin{equation}\label{eq:canonical-residual-system}
 C_{R,\mathbf v}u=W_Rb,
 \qquad
 g_{R,\mathbf v}(u)=G_Rb,
 \qquad
 g_{R,\mathbf v}(u)
   \coloneqq z_{\rm sym}(u)+G_R\Lambda_{R,\mathbf v}u.
\end{equation}
In particular, the polar map of $g_{R,\mathbf v}$ is exactly the polar map of
$z_{\rm sym}$; offsets and affine elimination do not change it.  If $b$ is
uniform in $\F_q^M$, then the two right-hand sides in
\eqref{eq:canonical-residual-system} are independent and uniform in
$\F_q^{M-K}$ and $\F_q^K$, respectively.
\end{lemma}

\begin{proof}
If a vector $v\in\F_q^M$ is killed by both $G_R$ and $W_R$, then
$W_Rv=0$ implies $v=Q_Ra$ for some $a\in\F_q^K$.  Applying $G_R$ gives
$a=G_RQ_Ra=0$, so $v=0$ and the square matrix in
\eqref{eq:stacked-output-change} is invertible.  Applying its two block rows
to $Q_Rz_{\rm sym}(u)+\Lambda_{R,\mathbf v}u=b$ gives
\eqref{eq:canonical-residual-system}; conversely, invertibility recovers the
original equation.  Affine maps have zero polar.  The final statement follows
because an invertible linear change sends a uniform vector to a uniform vector
on the product space.
\end{proof}

\begin{center}
\fcolorbox{deepblue}{lightblue}{\parbox{0.92\textwidth}{\small
\textbf{A two-output example.}  Suppose that, after the column substitution,
two verifier outputs over $\F_{19}$ are
\[
  y_1=q(u)+\ell_1(u),\qquad y_2=2q(u)+\ell_2(u),
\]
where $q$ is quadratic and $\ell_1,\ell_2$ are affine.  The combination
$2y_1-y_2=2\ell_1(u)-\ell_2(u)$ cancels the quadratic term.  Solving this
affine equation and substituting it back leaves one quadratic equation; neither
original output was dropped.  The rows of $W_R$ perform the same cancellation
simultaneously for the complete verifier.}}
\end{center}

\begin{proposition}[Exact affine reduction]\label{prop:affine-reduction}
Let $\rho=\rank Q_R$, $\lambda=\rank C_{R,\mathbf v}$, and
$\delta=M-\rho-\lambda$.  Then
\[
 \rank[Q_R\ \Lambda_{R,\mathbf v}]=\rho+\lambda,
\]
where $[Q_R\ \Lambda_{R,\mathbf v}]$ denotes horizontal matrix concatenation.  After row
reduction, exactly $\delta$ independent affine consistency conditions remain
on the target.  For every target satisfying those conditions, Gaussian
elimination leaves $\rho$ equations of degree at most two in
$N_0-\lambda$ variables, with exactly the same solution set as
\eqref{eq:ordered-system}.

In particular, if $\rho=K$ and $\lambda=M-K$, every target passes the
\emph{affine consistency test}, and the residual system has $K$ equations of
degree at most two on an affine translate of $\ker C_{R,\mathbf v}$.  This conclusion
does not assert that the residual system has a root.
\end{proposition}

\begin{proof}
The row space of $W_R$ is $\ker Q_R^\transpose$.  Since $W_R$ has full row
rank,
\[
 \ker W_R=\bigl(\ker Q_R^\transpose\bigr)^\perp=\operatorname{im}Q_R.
\]
Thus $W_R$ realizes the quotient map
$\F_q^M\to\F_q^M/\operatorname{im}Q_R$.  The image of the columns of
$\Lambda_{R,\mathbf v}$ in this quotient has dimension
$\rank(W_R \Lambda_{R,\mathbf v})=\lambda$, which proves
$\rank[Q_R\ \Lambda_{R,\mathbf v}]=\rho+\lambda$.

Write $b=t-c_{R,\mathbf v}$.  Projecting the substituted verifier by $W_R$ gives
\[
 C_{R,\mathbf v}u=W_R b.
\]
The codomain of $C_{R,\mathbf v}$ has dimension $M-\rho$, while its image has
dimension $\lambda$.  Solvability therefore imposes exactly
$(M-\rho)-\lambda=\delta$ independent affine conditions on the target.
For a target satisfying the affine consistency conditions, parameterize the
solution space of these affine equations by $N_0-\lambda$ free variables and
substitute into a complementary set of $\rho$ rows.  This leaves $\rho$
equations of degree at most two and uses only invertible row operations and
exact affine elimination, so the solution set is unchanged.  If $\rho=K$ and
$\lambda=M-K$, then $\delta=0$ and $C_{R,\mathbf v}$ is surjective onto its
$M-K$-dimensional codomain; hence every target passes the affine consistency
test.
\end{proof}

\paragraph{Operational meaning.}
The attacker computes $Q_R$, $G_R$, $W_R$, $C_{R,\mathbf v}$, and the canonical
residual map $g_{R,\mathbf v}$ from the public key before running an MQ solver.  If the ranks are unfavorable, the relation or offsets
are changed.  When $\rank Q_R=K$ and $\rank C_{R,\mathbf v}=M-K$, every hash target determines a
nonempty affine space $a(t)+\ker C_{R,\mathbf v}$ on which the canonical map
$g_{R,\mathbf v}$ supplies exactly $K$ equations of degree at most two.  This is an exact reformulation that discards no verifier
equation; whether the residual equations actually have a root is analyzed in
\Cref{sec:slices}.  Conversely, every root of the residual system maps back
through the column relation to a solution of the complete substituted verifier.

\section{The Symmetry Quotient}\label{sec:symmetry}

The central reduction is the formal version of
\eqref{eq:intro-swap}.  Readers may interpret $\Sym^2(A)$ simply as the vector
space spanned by \emph{unordered} pairs from $A$; the tensor notation records
why the identification is basis-independent.

For a fixed base form $P_i$, define
\[
  q_{i,a,b}(u)=u^\transpose\Sigma_{S^a}P_i\Sigma_{S^b}u.
\]
Transposition gives $q_{i,a,b}=q_{i,b,a}$.  In coordinates, the feature vector
therefore contains two copies of every off-diagonal polynomial.  The theorem
below shows that this is not an accident of the power basis: every alternating
label $\xi\otimes\eta-\eta\otimes\xi$ vanishes.  Write
$\F_q[u]^{\mathrm{hom}}_2$ for the vector space of homogeneous quadratic
polynomials in the coordinates of $u$.

\begin{theorem}[Symmetric-square factorization]\label{thm:symmetric-square}
Let $A=\F_q[S]$ have dimension $d$, and assume $S^\transpose=S$ and
$P_i^\transpose=P_i$.  For each $i$, the map
\[
 A\otimes_{\F_q}A\longrightarrow \F_q[u]^{\mathrm{hom}}_2,
 \qquad
 \xi\otimes\eta\longmapsto
 u^\transpose\Sigma_\xi P_i\Sigma_\eta u
\]
is invariant under $\xi\otimes\eta\mapsto\eta\otimes\xi$.  It therefore
factors through
\[
 \Sym^2_{\F_q}(A)
 =\bigl(A\otimes_{\F_q}A\bigr)
   \big/\Span\{\xi\otimes\eta-\eta\otimes\xi\}.
\]
Consequently, the direct sum over the $m_1$ base forms has rank at most
$m_1\binom{d+1}{2}$.  After any public linear map to $M$ outputs---in
particular, for every relation-dependent matrix $Q_R$---the rank is at most
\begin{equation}\label{eq:symmetry-rank}
 \min\left\{M,m_1\binom{d+1}{2}\right\}.
\end{equation}
If the minimal polynomial of $S$ has degree $\ell$, then $d=\ell$ and, in the
power basis,
\begin{equation}\label{eq:duplication}
 z_{\rm ord}(u)=Dz_{\rm sym}(u).
\end{equation}
\end{theorem}

\begin{proof}
Every element of $A$ is symmetric.  Since the displayed expression is a
scalar,
\[
 u^\transpose\Sigma_\xi P_i\Sigma_\eta u
 =\bigl(u^\transpose\Sigma_\xi P_i\Sigma_\eta u\bigr)^\transpose
 =u^\transpose\Sigma_\eta P_i\Sigma_\xi u.
\]
Thus every alternating label $\xi\otimes\eta-\eta\otimes\xi$ lies in the
kernel.  Taking the direct sum over $i$ gives the first rank bound, and a
linear map to $M$ public outputs gives the additional codomain cap.  When the
powers of $S$ form a basis of $A$, the factorization is exactly the coordinate
identity \eqref{eq:duplication}.
\end{proof}

\paragraph{Why the outer map cannot repair the loss.}
The public outer map receives only the values of these quadratic features.  If
two labeled inputs already define the same polynomial, no later linear
combination can separate them.  The outer map may have full rank on the
$K$-dimensional quotient, but it cannot enlarge that quotient.

For $\ell=4$, each base form loses six of its sixteen ordered coordinates:
$\ell^2=16$ becomes $\binom{\ell+1}{2}=10$.  For $\ell=2$, four ordered
coordinates become three.  These reductions occur for every relation and every
symmetric public key.

The theorem is an upper bound: additional public dependencies could make the
quadratic image even smaller.  The concrete attacks separately check
$\rank Q_R=K$ before solving, so the value $K$ used in the cost tables is the
exact rank observed in the public preflight for the chosen relation.

The factorization itself uses only transposition symmetry and is valid in every
characteristic.  Odd characteristic is needed only for the familiar direct-sum
interpretation
$A\otimes A=\Sym^2(A)\oplus\bigwedge^2(A)$.  The submitted $q=16$ system is
outside this theorem because its public base matrices are not symmetrized.

\subsection{A canonical Frobenius-channel normal form}
\label{sec:frobenius-channels}

The dimension statement $\dim_{\F_q}\Sym^2(A)=\binom{d+1}{2}$ can be sharpened
into an explicit coordinate system.  This matters algorithmically: the
coordinates have different Frobenius degrees, so an attacker need not treat all
ten quotient coordinates as unrelated quadratics over $\F_q$.

For $x,y\in A=\F_{q^d}$, write $[x\otimes y]_{\rm sym}$ for the class of
$x\otimes y$ in $\Sym^2_{\F_q}(A)$.  For $j>0$, define the symmetric
$\F_q$-bilinear map
\[
 \phi_j([x\otimes y]_{\rm sym})
 =xy^{q^j}+x^{q^j}y,
 \qquad
 \phi_0([x\otimes y]_{\rm sym})=xy.
\]
When $d$ is even and $j=d/2$, the value of $\phi_j$ is fixed by
$q^{d/2}$-Frobenius and therefore lies in $\F_{q^{d/2}}$.

\begin{theorem}[Frobenius-distance normal form]
\label{thm:frobenius-normal-form}
Let $A=\F_{q^d}$.  The direct sum of the maps above is an $\F_q$-linear
isomorphism
\[
 \Sym^2_{\F_q}(A)\xrightarrow{\ \sim\ }
 \begin{cases}
   A^{(d+1)/2}, & d\text{ odd},\\[2mm]
   A^{d/2}\oplus\F_{q^{d/2}}, & d\text{ even}.
 \end{cases}
\]
For even $d$, the $A$-coordinates are $\phi_0,\ldots,\phi_{d/2-1}$ and the
last coordinate is $\phi_{d/2}$.
\end{theorem}

\begin{proof}
Extend scalars to an algebraic closure and identify
$A\otimes_{\F_q}\overline{\F}_q\cong\overline{\F}_q^d$.  Index the primitive
idempotents by $i\in\mathbb Z/d\mathbb Z$, so Frobenius cyclically shifts the
indices.  In these coordinates, $\phi_0$ records the $d$ diagonal unordered
pairs $\{i,i\}$.  For $0<j<d/2$, the $i$th embedding of $\phi_j$ records the
unordered pair $\{i,i+j\}$; these form one orbit of size $d$.  When $d$ is
even and $j=d/2$, the pair $\{i,i+d/2\}$ is repeated at $i+d/2$, leaving one
orbit of size $d/2$, exactly the coordinates of $\F_{q^{d/2}}$.  The possible
Frobenius distances partition all unordered pairs.  The displayed map is
therefore a coordinate permutation after scalar extension, hence an
isomorphism over $\F_q$.
\end{proof}

For the public $\ell=4$ parameters, put $A=\F_{q^4}$ and
$B=\F_{q^2}\subset A$.  The theorem specializes to
\begin{equation}
\label{eq:frobenius-channel-isomorphism}
 \Sym^2_{\F_q}(A)\xrightarrow{\ \sim\ }A\oplus A\oplus B,
 \qquad
 [x\otimes y]_{\rm sym}\longmapsto
 \bigl(xy,\ xy^q+x^qy,\ xy^{q^2}+x^{q^2}y\bigr).
\end{equation}
The companion certificate constructs the resulting $10\times10$ matrix in
$A=\F_{19}[t]/(t^4-t-1)$ and verifies rank ten and determinant eight modulo
$19$.

\begin{corollary}[Polynomial shape of the three channels]
\label{cor:frobenius-channel-shape}
Let $L\cong A^s$ be an $A$-linear slice.  After the invertible dual output
change induced by \eqref{eq:frobenius-channel-isomorphism}, the ten quotient
coordinates belonging to each public base form become
\[
 F_{i,0}:A^s\to A,
 \qquad F_{i,1}:A^s\to A,
 \qquad F_{i,2}:A^s\to B.
\]
Their homogeneous parts are, after choosing $A$-coordinates on $L$, of the
forms
\begin{align*}
 F^{(2)}_{i,0}(x)&=\sum_{j,k}c^{(0)}_{i,j,k}x_jx_k,\\
 F^{(q+1)}_{i,1}(x)&=\sum_{j,k}c^{(1)}_{i,j,k}x_jx_k^q,\\
 F^{(q^2+1)}_{i,2}(x)&=\sum_{j,k}c^{(2)}_{i,j,k}x_jx_k^{q^2}.
\end{align*}
Translation to a target-dependent affine coset adds only terms in
$(x,x^q)$ to $F_{i,1}$ and in $(x,x^{q^2})$ to $F_{i,2}$.  Consequently their
ordinary total degrees over $A$ are at most $2$, $q+1$, and $q^2+1$.
Moreover, a uniform quotient target becomes an independent uniform element of
$A^{m_1}\oplus A^{m_1}\oplus B^{m_1}$.
\end{corollary}

\begin{proof}
Over a splitting field, \Cref{lem:spectral-quotient} groups the ten unordered
eigenblock features by cyclic distance: diagonal, adjacent, and opposite.
Descent around the diagonal orbit gives an $A$-valued ordinary quadratic;
descent around the adjacent orbit evaluates one variable block at $x$ and the
next at $x^q$; and the opposite orbit evaluates at $(x,x^{q^2})$ and has only
two conjugates.  Translation gives the stated lower-degree terms.  The target
claim follows because the dual of an invertible $\F_q$-linear coordinate
change is invertible.
\end{proof}

\section{Structured Stable Lifts}\label{sec:structured-lift}

\Cref{prop:affine-reduction} permits arbitrary offsets $v_j$.  Generic offsets
are useful for making the affine equations full rank, but they create linear
terms that mix all coordinates.  That mixing destroys the block grading used
by the selected published $\ell=4$ special-XL model.  We therefore need offsets
that do two jobs at once:
\begin{enumerate}[label=(\roman*),itemsep=2pt]
\item span the missing affine output directions, so that the full verifier is
      retained; and
\item confine all offset-linear terms to a small public row space, so that a
      large invariant kernel remains available for special XL.\@
\end{enumerate}

The following construction achieves both.  Choose one vector
$w\in\F_q^{N_0}$ and multipliers $\Gamma_0,\ldots,\Gamma_{r-1}\in A$, and set
\begin{equation}\label{eq:structured-offsets}
 v_j=\Sigma_{\Gamma_j}w.
\end{equation}
All offsets are therefore generated from the same $w$ by the public algebra
$A$.  Define the row space
\begin{equation}\label{eq:Uw}
 \mathcal U_w=
 \Span_{\F_q}\left\{
 w^\transpose\Sigma_\xi P_i\Sigma_\eta:
 1\le i\le m_1,\ \xi,\eta\in A
 \right\}
 \subseteq(\F_q^{N_0})^*.
\end{equation}
This space contains every linear form that can arise from crossing a variable
term with one of the structured offsets.  Its a priori bound is $m_1d^2$, not
$m_1\binom{d+1}{2}$: one argument is the fixed offset vector $w$ rather than
the variable $u$ on both sides, so reversing the two algebra labels need not
leave this \emph{linear} form unchanged.

A subspace $H\subseteq\F_q^{N_0}$ is called $A$-linear if
$\Sigma_\xi H\subseteq H$ for every $\xi\in A$.  Since
$A\cong\F_{q^\ell}$, this is extension-field linearity.  The earlier version
of this analysis took the $A$-orbit only of the projected affine constraints
and then required a maximum-rank equality to infer that all offset-linear rows
were killed.  That equality is unnecessary.  Take instead the complete right
$A$-orbit of the \emph{entire} offset-linear feature matrix.  For the basis
$1,S,\ldots,S^{d-1}$, form
\begin{equation}\label{eq:orbit-matrix}
 \mathcal O^{\mathcal F}_{R,\mathbf v}=
 \begin{bmatrix}
 \mathcal F_{R,\mathbf v}\\
 \mathcal F_{R,\mathbf v}\Sigma_S\\
 \vdots\\
 \mathcal F_{R,\mathbf v}\Sigma_{S^{d-1}}
 \end{bmatrix},
 \qquad
 \rho_{\mathcal F}=\rank\mathcal O^{\mathcal F}_{R,\mathbf v}.
\end{equation}

\begin{theorem}[Unconditional stable-lift kernel]\label{thm:stable-kernel}
Retain $S^\transpose=S$ and $P_i^\transpose=P_i$, use structured offsets
\eqref{eq:structured-offsets}, and suppose $A\cong\F_{q^d}$.  Then:
\begin{enumerate}[label=(\roman*)]
\item $\mathcal U_w$ is closed under right multiplication by $A$ and
$\dim_{\F_q}\mathcal U_w\le m_1d^2$;
\item every row of $\mathcal F_{R,\mathbf v}$ and
$\mathcal O^{\mathcal F}_{R,\mathbf v}$ lies in $\mathcal U_w$;
\item $\operatorname{row}(\mathcal O^{\mathcal F}_{R,\mathbf v})$ is an
$A$-subspace, so $d\mid\rho_{\mathcal F}$; and
\item
\begin{equation}\label{eq:automatic-stable-kernel}
 H^{\mathcal F}_{R,\mathbf v}
   =\ker\mathcal O^{\mathcal F}_{R,\mathbf v}
\end{equation}
is $A$-linear, satisfies
$\mathcal F_{R,\mathbf v}H^{\mathcal F}_{R,\mathbf v}=0$, and has
\begin{equation}\label{eq:automatic-stable-dimension}
 \dim_A H^{\mathcal F}_{R,\mathbf v}
   =n-\frac{\rho_{\mathcal F}}d
   \ge n-m_1d.
\end{equation}
\end{enumerate}
No orbit-rank equality or genericity assumption is required.  A rank smaller
than $m_1d^2$ only enlarges the attacker's stable search space.
\end{theorem}

\begin{proof}
See \Cref{app:proof-stable-kernel}.
\end{proof}

\paragraph{Interpretation.}
The stable slice now exists for every structured-offset instance.  Moving
inside $H^{\mathcal F}_{R,\mathbf v}$ changes the quadratic features but none
of the offset-linear features.  The offsets can therefore cover the missing
affine output directions, while the nonlinear solver retains a clean
$A$-linear block decomposition.  The former full-orbit-rank test has become a
harmless dimension measurement rather than an attack condition.

\begin{center}
\fcolorbox{deepblue}{lightblue}{\parbox{0.94\textwidth}{\small
\textbf{The stable-lift logic in four unconditional arrows.}
Structured offsets put every offset-linear row in $\mathcal U_w$; taking the
complete right-$A$ orbit of those rows makes their row space $A$-stable; its
kernel is therefore $A$-linear and automatically annihilates
$\mathcal F_{R,\mathbf v}$; and the bound
$\rho_{\mathcal F}\le m_1d^2$ guarantees at least $n-m_1d$ extension-field
search dimensions.  Any unexpected dependency helps rather than hurts.\@}}
\end{center}

\subsection{From an \texorpdfstring{$A$}{A}-linear kernel to four variable blocks}

Special XL needs variables divided into blocks with controlled
multidegrees.  The algebra $A$ supplies those blocks after extending scalars to
a field $\Omega$ in which the minimal polynomial of $S$ splits.  Over
$\Omega$, multiplication by $S$ diagonalizes into $\ell$ eigenspaces.  The
primitive idempotent $e_a$ is simply the projection onto the $a$th eigenspace.
This scalar extension is a solver coordinate change, not a relaxation of the
signature problem: the recovered point must still descend to $\F_q$ and pass
the original verifier.

Let $\mathbf e_a\in\mathbb Z^\ell$ denote the $a$th standard unit vector.  A
quadratic using only block $a$ has multidegree $2\mathbf e_a$; a bilinear term
using blocks $a$ and $b$ has multidegree
$\mathbf e_a+\mathbf e_b$.

\begin{lemma}[Eigenblock coordinates of the quotient]\label{lem:spectral-quotient}
Let $\Omega/\F_q$ be a splitting field of the minimal polynomial of $S$, and
let $e_1,\ldots,e_\ell$ be the primitive idempotents of
$A\otimes_{\F_q}\Omega$.  If $L$ is an $A$-linear space of $A$-dimension
$s$, then
\[
 L\otimes_{\F_q}\Omega=L_1\oplus\cdots\oplus L_\ell,
 \qquad L_a=\Sigma_{e_a}(L\otimes_{\F_q}\Omega),
 \qquad \dim_\Omega L_a=s.
\]
For each $i$, the eigenblock features
\[
 x^\transpose\Sigma_{e_a}P_i\Sigma_{e_b}x,
 \qquad 1\le a\le b\le\ell,
\]
are an invertible linear recombination of the symmetric power-pair features.
They have multidegree $2\mathbf e_a$ when $a=b$ and
$\mathbf e_a+\mathbf e_b$ when $a<b$.
\end{lemma}

\begin{proof}
Finite fields are perfect, so the irreducible minimal polynomial of $S$ is
separable.  Over $\Omega$, the power basis and the primitive-idempotent basis
of $A\otimes\Omega\cong\Omega^\ell$ are related by an invertible Vandermonde
matrix.  Passing to the symmetric square gives an invertible change between the
unordered power-pair features and the unordered eigenblock features.  The
idempotent $e_a$ projects $L\otimes\Omega$ onto $L_a$.  Hence a diagonal
feature uses only variables in $L_a$, while an off-diagonal feature is bilinear
in the variables from $L_a$ and $L_b$, giving the stated multidegrees.
\end{proof}

\begin{corollary}[Stable affine slices]\label{cor:stable-slices}
Assume the public checks
\begin{equation}\label{eq:stable-preflight}
 \rank Q_R=K,\qquad
 \rank C_{R,\mathbf v}=M-K
\end{equation}
pass, put $H_{R,\mathbf v}=H^{\mathcal F}_{R,\mathbf v}$, and let
$L\subseteq H_{R,\mathbf v}$ be an $A$-linear subspace of $A$-dimension $s$
(equivalently, base-field dimension $h=\ell s$).  By
\Cref{thm:stable-kernel}, such an $L$ is always available whenever
$s\le n-m_1\ell$.  For every
target, the affine equations in \Cref{prop:affine-reduction} have a
solution space $a(t)+\ker C_{R,\mathbf v}$ for a publicly computed base point
$a(t)$.  For every $y\in\ker C_{R,\mathbf v}$, the affine
slice $a(t)+y+L$ is contained in that space, and the full offset-linear feature
$\mathcal F_{R,\mathbf v}u$ is constant on the slice.

After extension to a splitting field, use the eigenblock coordinates of
\Cref{lem:spectral-quotient}.  Translation by $a(t)+y$ creates only linear and
constant terms in the same one or two blocks as the corresponding quadratic
term.  Introducing one homogenizing variable per block---an auxiliary variable that
is set to one after solving---therefore turns each complete affine equation
into a homogeneous equation of multidegree
$2\mathbf e_a$ or $\mathbf e_a+\mathbf e_b$.  If the restricted quadratic
feature map has total rank $K$, then the disjoint monomial supports of these
unordered multidegrees force rank $m_1$ in every component: there are
$\binom{\ell+1}{2}$ components, each has rank at most $m_1$, and their ranks
sum to $K=m_1\binom{\ell+1}{2}$.
\end{corollary}

\begin{proof}
See \Cref{app:proof-stable-slices}.
\end{proof}

\paragraph{What the preflight certifies.}
The two ranks in \eqref{eq:stable-preflight} say that the quotient has its full
expected rank and that the offsets supply every remaining affine output
direction.  The stable kernel and the identity
$\mathcal F_{R,\mathbf v}H_{R,\mathbf v}=0$ are now theorem consequences, not
rank hypotheses.  Computing $\rho_{\mathcal F}$ merely reports the actual
kernel dimension, which can only exceed the conservative lower bound.  The
restricted quadratic rank then certifies that no multidegree component was
lost on the chosen slice.  All remaining tests are public and occur before an
expensive solver run.

\subsection{Decoupling slice selection from fresh public coefficients}
\label{sec:seeded-decoupling}

The public $q=19$ expansion has more structure than an arbitrary symmetric
matrix.  In the coordinate order used by the reference implementation, write
the public variable space as an $A$-linear direct sum
\begin{equation}\label{eq:public-coordinate-split}
 \mathcal V_{\rm pub}=W\oplus V'\oplus O_{\rm pub},
 \qquad \dim_AW=1,
 \quad \dim_AV'=v-1,
 \quad \dim_AO_{\rm pub}=o.
\end{equation}
Here $W$ is one reserved public vinegar line, $V'$ is the sum of the remaining
public vinegar lines, and $O_{\rm pub}$ is the public oil-coordinate summand.
These are public coordinate labels; the attack does not recover the secret oil
space.

Choose a nonzero $w\in W$ for \eqref{eq:structured-offsets}.  Define the
\emph{cross data} to consist of the public outer-map data together with the
restrictions
\begin{equation}\label{eq:cross-data}
 P_{11,i}|_{W\times(W\oplus V')},
 \qquad
 P_{12,i}|_{W\times O_{\rm pub}},
 \qquad 1\le i\le m_1,
\end{equation}
and exclude both $P_{11,i}|_{V'\times V'}$ and the compressed $P_{22}$ block.
The relation $R$, the multipliers $\Gamma_j$, and a canonical bounded list of
$A$-linear candidate slices are chosen as deterministic functions of this
cross data.

\begin{lemma}[Coordinate separation in the pinned public expansion]
\label{lem:reference-expansion-separation}
In the pinned $q=19$ reference implementation, every independent scalar entry
of an upper-triangular diagonal block of $P_{11}$ is read from a distinct
public-expansion byte and only then copied to its transposed position.
Off-diagonal $P_{11}$ blocks are likewise filled entry by entry and their
transposes are copied afterward.  The bytes used for $P_{12}$, $P_{21}$, and
the public mixing data occupy disjoint later positions.  Finally,
\texttt{convert\_bytes\_to\_GF} reduces each byte coordinatewise modulo $19$.
Consequently, the raw byte coordinates underlying
$P_{11,i}|_{V'\times V'}$ are pairwise distinct and disjoint from every byte
coordinate in the cross data of \eqref{eq:cross-data}.  Under the random-XOF
idealization they are therefore mutually independent, independent of the
cross data, and each has maximum atom $14/256$.
\end{lemma}

\begin{proof}
This is a direct audit of the loops in \texttt{expand\_public}: the expansion
cursor is incremented once for each independent scalar before the symmetry
copy is performed, and the cursor is never reused.  The conversion routine
then applies reduction modulo $19$ separately to each occupied byte.  The
last sentence is the random-XOF interpretation of these disjoint input
coordinates, not a claim of literal statistical independence for a fixed
SHAKE seed~\cite{SNOVA23}.
\end{proof}
\begin{lemma}[Biased-product affine-subspace bound]
\label{lem:biased-affine-subspace}
Let $X=(X_1,\ldots,X_N)$ have independent coordinates in $\F_q$, each with
maximum atom at most $\rho$.  If $T:\F_q^N\to\F_q^r$ has rank $r$, then for
every $b\in\F_q^r$,
\begin{equation}\label{eq:biased-affine-bound}
 \Pr[T(X)=b]\le \rho^r.
\end{equation}
Equivalently, every affine subspace of codimension $r$ has probability at most
$\rho^r$ under this product distribution.
\end{lemma}

\begin{proof}
Choose $r$ pivot coordinates for $T$.  After conditioning on every nonpivot
coordinate, row reduction forces the pivot coordinates successively to one
specified value each.  Independence and the maximum-atom bound give a factor
at most $\rho$ per pivot.
\end{proof}

The following observation removes the former determinant-preflight loss for
the affine source.  It is the alternating analogue of
\Cref{lem:cross-tensor-injection}.

\begin{lemma}[Alternating affine-source injection]
\label{lem:alternating-source-injection}
Assume $q$ is odd, choose
\begin{equation}\label{eq:preflight-relation-offsets}
 R_j=S^j,\qquad \Gamma_0=1,\qquad \Gamma_1=0,
\end{equation}
and let $w\ne0$ span the reserved $A$-line $W$.  In the raw rows with column
pair $(j,k)=(1,0)$, quotienting the quadratic columns
$\mathcal R_RD$ leaves one copy of $\bigwedge^2_{\F_q}(A)$ per public base
form.  If $0\ne c=\sum_\nu a_\nu\wedge b_\nu\in\bigwedge^2(A)$, the
coefficient map of the corresponding offset-linear form on $V'$ has transpose
\begin{equation}\label{eq:Psi-alternating}
 \Psi_{c,w}(x)=
 \sum_\nu\bigl((b_\nu w)\mathbin\odot(a_\nu x)
                    -(a_\nu w)\mathbin\odot(b_\nu x)\bigr).
\end{equation}
The map $x\mapsto\Psi_{c,w}(x)$ is injective on $V'$.  Consequently, in the
random-XOF idealization,
\begin{equation}\label{eq:source-rank-failure}
 \Pr\!\left[
  \rank\bigl(\pi_{\rm alt}\mathcal F_{R,\mathbf v}|_{V'}\bigr)
       <m_1\binom d2
 \right]
 \le
 \frac{q^{m_1\binom d2}-1}{q-1}\,\rho^{d(v-1)},
\end{equation}
where $\pi_{\rm alt}$ denotes the displayed alternating quotient.  On the
complementary event,
\begin{equation}\label{eq:source-rank-lower}
 \rank[\mathcal R_RD\ \mathcal F_{R,\mathbf v}|_{V'}]
 \ge K+m_1\binom d2.
\end{equation}
\end{lemma}

\begin{proof}
As the power labels vary, the $(1,0)$ raw rows identify with $A\otimes A$;
the extra factor $R_1=S$ is invertible.  Their quadratic columns are ordinary
symmetrization, whose image and kernel in odd characteristic are respectively
$\Sym^2(A)$ and $\bigwedge^2(A)$.  The structured offset contributes
$B_i(bw,ax)$ on a pure tensor $a\otimes b$, which gives
\eqref{eq:Psi-alternating} on an alternating tensor.

For $x\ne0$ in $V'$, the $A$-lines $W$ and $Ax$ are disjoint.  Project the
symmetric tensor in \eqref{eq:Psi-alternating} to the cross summand
$W\otimes Ax$ and identify both lines with $A$.  The result is $-c$ after
invertible multiplication in the two factors, and is therefore nonzero.
Thus the transpose coefficient map has rank $\dim_{\F_q}V'=d(v-1)$.

Now take a nonzero dual tuple
$(c_1,\ldots,c_{m_1})\in(\bigwedge^2A)^{m_1}$ and choose $i$ with
$c_i\ne0$.  Condition on every public-form block except
$P_i|_{W\times V'}$.  The event that this tuple lies in the left kernel of the
alternating affine map imposes a rank-$d(v-1)$ affine system on those fresh
symbols, so \Cref{lem:biased-affine-subspace} bounds it by
$\rho^{d(v-1)}$.  Union-bound over projective nonzero dual tuples.  Trivial
intersection between the symmetric and alternating summands then gives
\eqref{eq:source-rank-lower}.
\end{proof}

\begin{theorem}[Random-XOF density of the complete structured preflight]
\label{thm:structured-preflight-density}
Continue with \eqref{eq:preflight-relation-offsets}, put
$J=W\oplus O_{\rm pub}$, and define
\begin{align}
 \delta_{\rm src}
 &=\frac{q^{m_1\binom d2}-1}{q-1}\rho^{d(v-1)},
 \label{eq:delta-src}\\
 \delta_{\rm proj}
 &=\rho^{m_1\binom{d+1}{2}}
   +\left(\frac{|A|^{1+o}-1}{|A|-1}-1\right)\rho^{m_1d^2},
 \label{eq:delta-proj}\\
 p_M(\rho)&=\prod_{c=1}^{M}(1-\rho^c).
 \label{eq:outer-full-rank-prob}
\end{align}
Under the random-XOF idealization, with probability at least
\begin{equation}\label{eq:complete-preflight-prob}
 p_{\rm pre}=p_M(\rho)(1-\delta_{\rm src})(1-\delta_{\rm proj}),
\end{equation}
all of the following hold simultaneously:
\begin{enumerate}[label=(\roman*)]
\item $\rank Q_R=K$ and $\rank C_{R,\mathbf v}=M-K$;
\item $\operatorname{row}(\mathcal F_{R,\mathbf v})=\mathcal U_w$ and
      $H^{\mathcal F}_{R,\mathbf v}\cap J=\{0\}$; and
\item projection of the entire stable kernel to $V'$ is injective, so every
      subspace of it automatically passes \eqref{eq:projection-preflight}.
\end{enumerate}
For the six official $\ell=d=4$ shapes, the exact bounds are:
\begin{center}
\begin{tabular}{@{}c c c c c@{}}
\toprule
shape & $m_1\binom42$ proved/needed & $\log_2\delta_{\rm src}$ &
$\log_2\delta_{\rm proj}$ & $p_{\rm pre}$\\
\midrule
I-a   & $30/30$ & $-329.54$ & $-209.63$ & $>0.9423222679-2^{-209}$\\
I-b   & $30/30$ & $-329.54$ & $-209.63$ & $>0.9423222679-2^{-209}$\\
III-a & $42/42$ & $-479.81$ & $-293.49$ & $>0.9423222679-2^{-293}$\\
III-b & $42/30$ & $-446.27$ & $-293.49$ & $>0.9423222679-2^{-293}$\\
V-a   & $54/54$ & $-596.54$ & $-377.34$ & $>0.9423222679-2^{-377}$\\
V-b   & $54/54$ & $-630.08$ & $-377.34$ & $>0.9423222679-2^{-377}$\\
\bottomrule
\end{tabular}
\end{center}
All decimals are rounded downward at the displayed precision; the artifact
stores the exact rational values.
\end{theorem}

\begin{proof}
Every row of $\mathcal F_{R,\mathbf v}$ lies in $\mathcal U_w$.  Conversely,
the $(1,0)$ raw column pair in the proof of
\Cref{lem:alternating-source-injection}, as its two power labels vary, exposes
every generator $w^\transpose\Sigma_\xi P_i\Sigma_\eta$ of $\mathcal U_w$.
Thus the two row spaces agree; since $\mathcal U_w$ is right-$A$ stable, the
complete orbit matrix has the same row space and
$H^{\mathcal F}_{R,\mathbf v}=\mathcal U_w^\perp$.

If $0\ne x\in W$, then $Ax=W$.  The condition $x\in\mathcal U_w^\perp$
forces every symmetric restriction $P_i|_{W\times W}$ to vanish, imposing
$m_1\binom{d+1}{2}$ independent scalar conditions.  If
$x\in J\setminus W$, then $Ax$ is an $A$-line disjoint from $W$; the
$W\times O_{\rm pub}$ symbols map surjectively to
$\operatorname{Bil}(W,Ax)$, so the same condition imposes $m_1d^2$
independent scalar conditions.  The event depends only on the $A$-line $Ax$.
There is one line $W$ and
$(|A|^{1+o}-1)/(|A|-1)-1$ other lines in $J$.  Apply
\Cref{lem:biased-affine-subspace} and union-bound to obtain
\eqref{eq:delta-proj}.  The source event of
\Cref{lem:alternating-source-injection} uses only $W\times V'$ symbols,
whereas this projection event uses only the symmetric $W\times W$ symbols and
$W\times O_{\rm pub}$ symbols.  They are independent by
\Cref{lem:reference-expansion-separation}.

On the source event, choose $M-K$ affine columns whose alternating images are
independent and adjoin all $K$ quadratic columns.  This gives an
$M$-dimensional source subspace with the quadratic columns as a distinguished
subset.  Conditional on that public-form data, each independent outer-map row
has a surjective image in $\F_q^M$.  If the previous rows span a subspace of
codimension $c$, \Cref{lem:biased-affine-subspace} bounds the chance that the
next row remains in it by $\rho^c$.  Multiplying the complementary
probabilities gives $p_M(\rho)$.  Invertibility on the chosen source columns
implies both $\rank Q_R=K$ and
$\rank[Q_R\ \Lambda_{R,\mathbf v}]=M$; then
\Cref{prop:affine-reduction} gives $\rank C_{R,\mathbf v}=M-K$.
The outer-map bytes are disjoint from both public-form regions, so the three
success probabilities multiply.  Finally,
$H^{\mathcal F}_{R,\mathbf v}\cap J=\{0\}$ is exactly injectivity of its
projection to the complementary summand $V'$.
\end{proof}

\begin{lemma}[Reserved-line decoupling]\label{lem:reserved-line-decoupling}
With $w\in W$, the matrices $Q_R$, $C_{R,\mathbf v}$,
$\mathcal O^{\mathcal F}_{R,\mathbf v}$, the stable kernel $H_{R,\mathbf v}$, and every
candidate slice selected canonically from that kernel depend only on the cross
data in \eqref{eq:cross-data}.  In particular, they are independent of the
internal public coefficients $P_{11,i}|_{V'\times V'}$ in the random-XOF
idealization.

A candidate used in the seeded analysis additionally passes the exact public
projection preflight
\begin{equation}\label{eq:projection-preflight}
 L\cap(W\oplus O_{\rm pub})=\{0\},
 \qquad\text{equivalently}\qquad
 \rank(\pi_{V'}|_L)=\dim_{\F_q}L.
\end{equation}
The full-key restricted-quadratic-rank and verifier preflights may be applied
after the bounded candidate list is fixed.
\end{lemma}

\begin{proof}
The quotient matrix $Q_R$ is determined by the public outer map and the
relation.  Every row created by a structured offset is a linear combination of
forms $(\Sigma_\xi w)^\transpose P_i\Sigma_\eta$.  Its left argument remains
inside $W$, so only the $W$-row restrictions in \eqref{eq:cross-data} occur:
$P_{22}$ and the internal $V'\times V'$ block never appear.  The same remains
true after taking the complete right $A$-orbit in
\eqref{eq:orbit-matrix}.  Kernels, canonical row-reduced bases, and candidate
subspaces derived from these matrices are therefore functions of the cross
data alone.  The last checks are ordinary public rank tests.  Selecting the
first full-key candidate that passes them from a list of size $J$ costs only a
union factor $J$ in the seeded failure bound below.
\end{proof}

The purpose of this construction is not to hide information from the
attacker---all coefficients are public after key generation.  It is to make a
valid probability statement: the slice is fixed using one disjoint portion of
the public seed expansion, while the internal $P_{11}$ coefficients that force
its polar rank remain fresh.  \Cref{thm:seeded-spectrum} turns this data
separation into an explicit aggregate-spectrum theorem.

\subsection{Passing the verifier's block-symmetry check}\label{sec:rejection}

The public verifier does more than evaluate the quadratic map.  It also rejects
signatures containing too many symmetric blocks.  A root of the reduced MQ
system is useful only if it passes this final check.  We therefore build
rejection compatibility into the reduction rather than treating it as an
unmodeled success probability.

Version~2.3 rejects a square odd-characteristic signature if one of its blocks
is symmetric when $\ell>2$, and rejects an $\ell=2$ signature if more than
$\lfloor n/4\rfloor$ blocks are symmetric~\cite{SNOVA23}.  Rectangular blocks
are not subject to this rule.

\paragraph{The structured $\ell=4$ rows.}
Only the square rows with $r=\ell=4$ are subject to the verifier's strict
``no symmetric block'' rule; the rectangular $r=5,6$ rows have no such check.
For a relation $R=(R_0,R_1,R_2,R_3)$, define the local skew map
\begin{equation}\label{eq:local-skew-map}
 \mathsf S_R:\F_q^4\longrightarrow\bigwedge^2\F_q^4,
 \qquad
 \mathsf S_R(x)_{a,c}=(R_cx)_a-(R_ax)_c
 \quad(0\le a<c<4).
\end{equation}
For any offsets, one signature block is symmetric exactly when
$\mathsf S_R(x)$ equals a fixed offset-dependent vector.  Hence the public
preflight $\rank\mathsf S_R=4$ makes the symmetry equation have either zero or
one solution in that block.  Because the $n$ block variables are disjoint, the
accepted subset of the full common-column space has density
\begin{equation}\label{eq:l4-accepted-density}
 a_{4,n}\ge(1-q^{-4})^n.
\end{equation}
No independence from the MQ equations is assumed; the accepted-singleton
theorem below uses only this global density.

\begin{lemma}[An exact injective-skew relation for the official $q=19$ matrix]
\label{lem:l4-skew-rank}
For the official $\ell=4$ matrix
\[
 S=\begin{pmatrix}
 1&2&3&0\\2&3&0&1\\3&0&1&2\\0&1&2&15
 \end{pmatrix}
 \quad\text{over }\F_{19},
\]
the powers relation $R_j=S^j$, $0\le j<4$, has
\[
 [\mathsf S_R]=
 \begin{pmatrix}
 1&1&3&0\\14&8&5&8\\10&3&7&6\\
 5&14&7&16\\3&18&9&0\\18&12&7&1
 \end{pmatrix},
 \qquad \rank_{\F_{19}}\mathsf S_R=4.
\]
Thus the injective-skew preflight is nonvacuous.  The attack retains a relation
only when this rank and the independent $Q_R$, affine, orbit, and restricted
quadratic-rank preflights all pass.
\end{lemma}

\begin{proof}
The six rows are the coefficients of
$(S^cx)_a-(S^ax)_c$ for $0\le a<c<4$.  Row reduction modulo $19$ gives four
pivots.  The companion artifact recomputes $S$, its powers, the displayed
matrix, and its rank from integer data.
\end{proof}

For the three square rows, $n=33,47,59$.  Combining
\eqref{eq:l4-accepted-density} with \Cref{thm:accepted-singleton} changes the
compact retry denominator from $1-19^{-1}$ to $a_{4,n}-19^{-1}$, and the
robust denominator from $1/2$ to $a_{4,n}-1/2$.  The largest resulting
surcharge is only $0.001307$ bits.  It is included in the artifact and is
smaller than the rounding precision of the displayed cost tables.  This is a
probabilistic bypass with an exact density bound: unlike the fixed-skew
$\ell=2$ construction below, it does not claim that every point of the search
space is accepted.

\paragraph{A deterministic $\ell=2$ construction.}
The first two parameter levels use offsets that force a fixed nonzero skew in
every block.

\begin{proposition}[Fixed-skew $\ell=2$ lift]\label{prop:fixed-skew}
For the Version~2.3 $\ell=2$ matrix
\[
 S=\begin{pmatrix}1&2\\2&15\end{pmatrix},
 \qquad R_0=I,
 \qquad R_1=9I+10S=\begin{pmatrix}0&1\\1&7\end{pmatrix},
\]
choose offsets satisfying
\begin{equation}\label{eq:fixed-skew-offset}
 (v_1)_{b,0}-(v_0)_{b,1}=1
\end{equation}
for every block $b$.  Then every block of every signature satisfying the
column relation has
\begin{equation}\label{eq:fixed-skew-identity}
 U_b[0,1]-U_b[1,0]=1.
\end{equation}
Consequently every solution passes the $\ell=2$ symmetric-block rejection
check.  If the public preflight also gives
$\rank Q_R=K$ and $\rank C_{R,\mathbf v}=M-K$, every target passes the
affine consistency test and leaves $K$ equations of degree at most two in
$N_0-M+K$ variables.  Every root of that residual system is a valid signature.
\end{proposition}

\begin{proof}
The first row of $R_1$ is the second row of $R_0=I$.  Hence the variable part
of $U_b[0,1]-U_b[1,0]$ vanishes.  The remaining offset difference is one by
\eqref{eq:fixed-skew-offset}, proving \eqref{eq:fixed-skew-identity}.  No block
is symmetric.  The final statement is the full-lift case of
\Cref{prop:affine-reduction}.
\end{proof}

\begin{corollary}[Symmetric-block rejection is structurally insufficient]
\label{cor:symmetry-rejection-insufficient}
For the public $\ell=2$ Version~2.3 verifier, any countermeasure whose only
additional condition is that a signature contain sufficiently few symmetric
blocks leaves the fixed-skew affine family of
\Cref{prop:fixed-skew} untouched: every block in every candidate has the
prescribed nonzero skew $1$.  Such a rule can reject honest-looking symmetric
blocks, but it cannot exclude this attack family without testing some further
property.
\end{corollary}

\begin{proof}
Equation~\eqref{eq:fixed-skew-identity} holds identically on the entire affine
family, before any MQ equation is solved.  Hence none of its blocks is
symmetric, independently of the target or root.
\end{proof}

Thus the fixed-skew construction has no rejection probability: every MQ root
is acceptable.  The Level-V headline attack uses a faster alternative in which
targets are filtered before solving.

\paragraph{The zero-offset $\ell=2$ filter.}
Set $R_0=I$, $R_1=0$, and use zero offsets.  Each block then has the form
\[
 \begin{pmatrix}x&0\\ y&0\end{pmatrix},
\]
so it is symmetric exactly when $y=0$.  These rejection events are nonzero
hyperplanes on disjoint block coordinates.  This relation is chosen
carefully: a relation that forced the two off-diagonal entries to agree would
make every block symmetric and would be useless.

Let $N$ be the residual-domain dimension ($N=N_0$ for this zero-offset lift)
and let $g:\F_q^N\to\F_q^K$ be the residual map of degree at most two.
For each nonzero
$b\in\F_q^K$, take the polar rank of the scalar quadratic $b\cdot g$, and let
$r_*$ be the minimum of these ranks.  Thus $r_*$ measures the worst output
direction; a large $r_*$ forces every target fiber to have nearly its expected
size.

Put
\[
 t_{\rm rej}=\lfloor n/4\rfloor+1,
 \qquad B=\binom n{t_{\rm rej}},
 \qquad \alpha=1-q^{-K},
\]
and let $\mathcal X_{\rm acc}$ be the assignments accepted by the block-symmetry filter.
At least $t_{\rm rej}$ symmetric blocks are needed for rejection.  Because
the $n$ hyperplanes use disjoint coordinates, the exact density of accepted
assignments is
\begin{equation}\label{eq:l2-accepted-density}
 a_n\coloneqq\frac{|\mathcal X_{\rm acc}|}{q^N}
 =q^{-n}\sum_{k=0}^{\lfloor n/4\rfloor}\binom nk(q-1)^{n-k}.
\end{equation}

\begin{proposition}[Accepted targets for the $\ell=2$ filter]\label{prop:l2-rejection}
Every target has an accepted root if
\begin{equation}\label{eq:l2-rejection}
 1-Bq^{-t_{\rm rej}}>\alpha(1+B)q^{K-r_*/2}.
\end{equation}
For a uniform target $z\in\F_q^K$,
\begin{equation}\label{eq:l2-accepted-target}
 \Pr\bigl[g^{-1}(z)\cap\mathcal X_{\rm acc}\ne\varnothing\bigr]
 \ge \frac{a_n}{1+\alpha q^{K-r_*/2}}.
\end{equation}
Moreover, if $1-\alpha q^{K-r_*/2}>0$, every target $z$ satisfies
\begin{equation}\label{eq:l2-accepted-fraction}
 \frac{|g^{-1}(z)\cap\mathcal X_{\rm acc}|}{|g^{-1}(z)|}
 \ge
 1-\frac{Bq^{-t_{\rm rej}}+\alpha Bq^{K-r_*/2}}
          {1-\alpha q^{K-r_*/2}},
\end{equation}
whenever the right-hand side is positive.
\end{proposition}

\begin{proof}
See \Cref{app:proof-l2-filter}.
\end{proof}

The three conclusions answer different questions.  Readers interested only
in the concrete attack may retain the following takeaway.  The global
condition $r_*\ge240$ gives an accepted root for every Level-V target.  The
much weaker condition $r_*\ge194$ makes the expected number of uniform targets
to an accepted root less than $2^{0.075}$ by
\Cref{cor:target-cushions}; this is the condition charged in the headline.
More generally, \eqref{eq:l2-rejection} is a sufficient worst-case condition,
\eqref{eq:l2-accepted-target} controls expected retries, and
\eqref{eq:l2-accepted-fraction} says that, target by target, almost every root
is accepted when the rank is high enough.

\section{Root Existence, Uniqueness, and Retry Bounds}\label{sec:slices}

After the exact affine reduction, the structured attack repeatedly chooses a
hash target and an affine coset of the selected slice direction space.  We call
one such pair a \emph{target--slice trial}.  The attack asks whether the
resulting MQ system has a recoverable root.  Three events must not be conflated:
\begin{enumerate}[label=(\roman*),itemsep=2pt]
\item the slice contains at least one base-field root;
\item it contains exactly one base-field root; and
\item the selected Macaulay degree exposes a finite quotient algebra from
      which all algebraic roots can be extracted.
\end{enumerate}
This section proves statements about the first two events.  The third is the
separate special-XL solving condition in \Cref{sec:special-xl}; exact corank
one is only its simplest special case.

Let $V$ be the full residual direction space, let $L\subseteq V$ be the chosen
slice direction space, and write $h=\dim L$.  For a map of degree at most two
$g:V\to\F_q^K$, a random-map heuristic predicts
\[
  \mu=q^{h-K}
\]
roots in a random affine coset of $L$ above a random target.  When $h=K-\tau$ with $\tau\ge0$, this mean is $q^{-\tau}$, suggesting
about $q^\tau$ trials.  The obstruction detected by a second-moment analysis is excess
collisions: distinct points in the same slice mapping to the same target.
Polar derivatives count those collisions exactly, so the argument below is a
theorem rather than a random-system heuristic.  Full cross-polar rank is a clean special case, but far stronger than the attack needs.  The
headline condition is the aggregate bound $\bar\eta_L\le q^{-1}$.
The pointwise condition $\rank T_y\ge h+1$ for every nonzero slice direction
is a sufficient certificate; it keeps the unique-root retry count
within a factor $q/(q-1)$ of $q^{K-h}$.  The theorem first states the exact formulas and
then quantifies both sparse defects and weaker uniform rank cushions.

Write the vector-valued polar map as
\[
 \beta(x,y)=g(x+y)-g(x)-g(y)+g(0).
\]
For $b\in\F_q^K\setminus\{0\}$, define the cross-polar map
\begin{equation}\label{eq:cross-polar-map}
 \Phi_{b,L}:V\longrightarrow L^*,
 \qquad \Phi_{b,L}(x)(y)=b\cdot\beta(x,y),
\end{equation}
and let $e_L(b)=\rank\Phi_{b,L}$.  Put
\begin{equation}\label{eq:theta-def}
 e_* = \min_{b\ne0}e_L(b),
 \qquad
 \Theta_L(g)=\sum_{b\ne0}q^{-e_L(b)}.
\end{equation}
For $y\in L\setminus\{0\}$, also define the directional derivative
\begin{equation}\label{eq:directional-map}
 T_y:V\longrightarrow\F_q^K,
 \qquad T_y(x)=\beta(x,y),
\end{equation}
write $r(y)=\rank T_y$, and let
$\kappa(y)$ be one when $-(g(y)-g(0))\in\operatorname{im}T_y$ and zero
otherwise.  The indicator $\kappa(y)$ records whether two points separated by
$y$ can collide at all; the rank $r(y)$ then determines how many such pairs
exist.

\begin{theorem}[Root collisions and the rank-spectrum bound]
\label{thm:coset-success}
Choose an affine coset $C$ uniformly from $V/L$ and choose
$z\in\F_q^K$ uniformly.  Let
\[
 Z=|\{x\in C:g(x)=z\}|,
 \qquad \mu=q^{h-K},
\]
and define
\begin{equation}\label{eq:eta-def}
 \eta_0=\sum_{y\in L\setminus\{0\}}\kappa(y)q^{-r(y)},
 \qquad
 \bar\eta=\sum_{y\in L\setminus\{0\}}q^{-r(y)}.
\end{equation}
Then
\begin{align}
 \mathbb E Z&=\mu,\label{eq:EZ}\\
 \mathbb E[Z(Z-1)]&=\mu\eta_0,\label{eq:factorial-moment}\\
 \eta_0&\le\bar\eta
 =\mu\bigl(1+\Theta_L(g)\bigr)-1.\label{eq:eta-spectrum}
\end{align}
Consequently,
\begin{align}
 \operatorname{Var}(Z)&\le\mu^2\Theta_L(g),\label{eq:theta-var}\\
 \Pr[Z>0]&\ge\frac{\mu}{1+\eta_0}
 \ge\frac{1}{1+\Theta_L(g)}.\label{eq:coset-success}
\end{align}
If $h=K-\tau$ and $e_*\ge h-\Delta$, then
\begin{equation}\label{eq:defect-retry}
 \Theta_L(g)<q^{\tau+\Delta},
 \qquad
 \mathbb E[\textnormal{trials to a root-containing target--coset pair}]
 <1+q^{\tau+\Delta}.
\end{equation}
This last bound controls the search for a root-containing target--coset pair,
not uniqueness of the root.
Here $\Delta$ is the maximum loss from full cross-polar rank $h$.
\end{theorem}

\begin{proof}
See \Cref{app:proof-coset-success}.
\end{proof}

\paragraph{How to read the theorem.}
The first moment $\mathbb EZ=\mu$ is exactly the random-map value.  The second
factorial moment $\mathbb E[Z(Z-1)]$ counts ordered pairs of distinct roots and
therefore measures collisions.  The compatible collision mass $\eta_0$ is
exact; $\bar\eta$ drops the compatibility test and is a convenient upper
bound.  The sum $\Theta_L(g)$ describes the entire cross-polar rank spectrum,
not merely its minimum.  A few isolated rank defects may have little effect,
whereas low rank in many output directions can substantially increase retries.

\begin{center}
\begin{tabularx}{0.92\textwidth}{@{}>{\bfseries}p{0.27\textwidth}X@{}}
$\mu=q^{h-K}$ & expected number of roots in a random target--slice pair\\
$\eta_0$ & exact expected number of compatible collisions per planted root\\
$\bar\eta$ & rank-only upper bound on that collision mass\\
directional rank $r(y)\ge h+1$ & unique-root trial count at most
$q^{K-h}/(1-q^{-1})$
\end{tabularx}
\end{center}

\begin{corollary}[Unique-root trials and the planted distribution]
\label{cor:singleton-coset}
In the setting of \Cref{thm:coset-success},
\begin{align}
 \Pr[Z=1]&\ge\mu(1-\eta_0)
             \ge\mu(1-\bar\eta),\label{eq:singleton-coset}\\
 \Pr[Z>0]&\ge\mu\left(1-\frac{\bar\eta}{2}\right),
 \qquad (\bar\eta<2),\label{eq:bonferroni-success}\\
 \Pr[Z\ge2\mid Z>0]&\le\frac{\bar\eta}{2-\bar\eta},
 \qquad (\bar\eta<2).\label{eq:conditional-multiple}
\end{align}
In particular, if $\bar\eta<1$, the expected number of independent trials
until $Z=1$ is at most
\begin{equation}\label{eq:unique-retry-general}
 \frac{1}{\mu(1-\bar\eta)}.
\end{equation}

Let $\mathsf U$ be the uniform target--coset pair conditioned on $Z>0$, and
let $\mathsf P$ be the planted distribution obtained by choosing a uniform
coset, a uniform point in it, and the point's image as target.  If
$\bar\eta<1$, then
\begin{equation}\label{eq:planted-tv}
 \Pr_{\mathsf P}[Z\ge2]\le\eta_0\le\bar\eta,
 \qquad
 d_{\rm TV}(\mathsf U,\mathsf P)\le\bar\eta.
\end{equation}
\end{corollary}

\begin{proof}
See \Cref{app:proof-singleton-coset}.
\end{proof}

\begin{theorem}[Accepted singleton before affine elimination]
\label{thm:accepted-singleton}
Let $U$ be a finite-dimensional $\F_q$-space, let
$\mathsf C:U\to\F_q^c$ be surjective, let
$L\subseteq\ker\mathsf C$ have dimension $h$, and let
$g:U\to\F_q^K$ have degree at most two.  Choose $d\in\F_q^c$ and
$z\in\F_q^K$ independently and uniformly, then choose uniformly an affine
coset $\mathcal A$ of $L$ contained in $\mathsf C^{-1}(d)$.  For an arbitrary
accepted set $\mathcal X\subseteq U$ of density
$a=|\mathcal X|/|U|$, put
\[
 Z=|\{x\in\mathcal A:g(x)=z\}|,
 \qquad
 Z_{\mathcal X}=|\{x\in\mathcal A\cap\mathcal X:g(x)=z\}|,
 \qquad \mu=q^{h-K}.
\]
Define $T_y(x)=g(x+y)-g(x)-g(y)+g(0)$ and
$\bar\eta_L=\sum_{0\ne y\in L}q^{-\rank T_y}$, as above.  Then
\begin{align}
 \mathbb E Z_{\mathcal X}&=a\mu,\label{eq:accepted-first-moment}\\
 \mathbb E[Z(Z-1)]&=\mu\eta_0\le\mu\bar\eta_L,\label{eq:accepted-pair-bound}\\
 \Pr[Z=Z_{\mathcal X}=1]
 &\ge \mu(a-\eta_0)
 \ge \mu(a-\bar\eta_L).\label{eq:accepted-singleton-bound}
\end{align}
Consequently, if $a>\bar\eta_L$, the expected number of independent trials to
an accepted root that is the unique base-field root of its trial is at most
\begin{equation}\label{eq:accepted-singleton-retry}
 \frac{q^{K-h}}{a-\bar\eta_L}.
\end{equation}
The bound requires no independence between acceptance and the polynomial map.
\end{theorem}

\begin{proof}
Because $L\subseteq\ker\mathsf C$ and $\mathsf C$ is surjective, sampling
$d$ and then a uniform $L$-coset in its fiber is exactly uniform over all
cosets in $U/L$.  Every accepted point is therefore selected with probability
$q^{h-\dim U}$, and its image matches uniform $z$ with probability $q^{-K}$,
which proves \eqref{eq:accepted-first-moment}.  The ordered-pair calculation is
identical to \Cref{thm:coset-success} and gives
\eqref{eq:accepted-pair-bound}.  Pointwise,
\[
 \mathbf 1_{\{Z=Z_{\mathcal X}=1\}}
 \ge Z_{\mathcal X}-Z(Z-1):
\]
for $Z=0$ both sides vanish, for $Z=1$ the right side is one exactly when the
root is accepted, and for $Z\ge2$ the right side is nonpositive.  Taking
expectations proves \eqref{eq:accepted-singleton-bound} and then
\eqref{eq:accepted-singleton-retry}.
\end{proof}

\begin{theorem}[Accepted nonsingular roots from the same aggregate spectrum]
\label{thm:accepted-nonsingular-roots}
Assume the setting of \Cref{thm:accepted-singleton}.  Write the residual map as
$g=q_2+q_1+q_0$, where $q_i$ is homogeneous of degree $i$, and define the
restricted formal Jacobian
\[
 D_Lg(x):L\longrightarrow\F_q^K,
 \qquad D_Lg(x)[y]=\beta(x,y)+q_1(y).
\]
Let $\mathcal X\subseteq U$ be the signatures accepted by the public verifier,
with density $\alpha$, and let $Y$ be the number of points in the sampled
fiber--coset trial that lie in $\mathcal X$, map to the sampled residual
target, and satisfy $\rank D_Lg(x)=h$.  Then
\begin{equation}\label{eq:accepted-nonsingular-success}
 \Pr[Y>0]
 \ge
 \mu\,
 \frac{\left(\alpha-\bar\eta_L/(q-1)\right)^2}
      {1+\bar\eta_L},
 \qquad \mu=q^{h-K},
\end{equation}
whenever $\alpha>\bar\eta_L/(q-1)$.  Hence the expected number of independent
trials to an accepted root at which the full restricted residual Jacobian has
rank $h$ is at most the reciprocal of the right-hand side.

The theorem does not require that the root be unique, and it assumes no
independence among acceptance, Jacobian rank, the unused residual equations,
or the polynomial map.
\end{theorem}

\begin{proof}
For a nonzero direction $y\in L$, the condition $D_Lg(x)[y]=0$ is an affine
system in $x$ whose linear part is $T_y:x\mapsto\beta(x,y)$.  It therefore has
at most $q^{\dim U-\rank T_y}$ solutions.  Union-bound over projective
directions $[y]\in\mathbb P(L)(\F_q)$ to obtain
\begin{equation}\label{eq:jacobian-defect-mass}
 \frac{|\{x\in U:\rank D_Lg(x)<h\}|}{|U|}
 \le \sum_{[y]\in\mathbb P(L)(\F_q)}q^{-\rank T_y}
 =\frac{\bar\eta_L}{q-1}.
\end{equation}
Averaging over the uniform fiber--coset and residual target exactly as in
\Cref{thm:accepted-singleton} gives
\[
 \mathbb E Y\ge
 \mu\left(\alpha-\frac{\bar\eta_L}{q-1}\right).
\]
Moreover $Y(Y-1)\le Z(Z-1)$, where $Z$ counts all residual roots in the trial,
so \Cref{thm:coset-success} gives
$\mathbb E[Y(Y-1)]\le\mu\bar\eta_L$.  Since $Y\le Z$ and
$\mathbb E Z=\mu$,
\[
 \mathbb E[Y^2]
 =\mathbb E Y+\mathbb E[Y(Y-1)]
 \le \mu(1+\bar\eta_L).
\]
Cauchy--Schwarz in the form
$\Pr[Y>0]\ge(\mathbb EY)^2/\mathbb E[Y^2]$ proves
\eqref{eq:accepted-nonsingular-success}.
\end{proof}

Total variation distance is the largest discrepancy that the two distributions
can assign to any event.  Thus, when $\bar\eta$ is tiny, an experiment that
plants a known root samples essentially the same target--coset distribution as
a uniformly chosen trial conditioned on being solvable.  This justifies using
planted instances to probe solver behavior; it does not remove the separate
finite-size-to-production extrapolation.

\begin{corollary}[A one-rank cushion suffices]\label{cor:rank-cushion}
Let
\[
 r_{\min}(L)=\min_{y\in L\setminus\{0\}}\rank T_y.
\]
If $r_{\min}(L)\ge h+c$ for an integer $c\ge1$, then
\begin{equation}\label{eq:rank-cushion-eta}
 \bar\eta\le(q^h-1)q^{-(h+c)}<q^{-c}.
\end{equation}
Consequently, the expected number of independent target--slice trials until a
unique base-field root is at most
\begin{equation}\label{eq:rank-cushion-retry}
 \frac{q^{K-h}}{1-q^{-c}}.
\end{equation}
For the headline choice $q=19$ and $c=1$, the extra exponent over $q^{K-h}$
is at most
\[
 \log_2\frac{19}{18}=0.0780026\ldots\ \text{bits}.
\]
\end{corollary}

\begin{proof}
There are $q^h-1$ nonzero vectors in $L$.  Substitute
$r(y)\ge h+c$ into the definition of $\bar\eta$, and then apply
\eqref{eq:unique-retry-general}.
\end{proof}

\begin{corollary}[Full cross-polar rank]\label{cor:full-cross-rank}
The following are equivalent:
\begin{enumerate}[label=(\roman*)]
\item $e_L(b)=h$ for every $b\ne0$;
\item $T_y$ is surjective for every $y\in L\setminus\{0\}$.
\end{enumerate}
Under these conditions,
\[
 \eta_0=\bar\eta=\mu-q^{-K}.
\]
If $h=K-\tau$ with $\tau\ge0$, the expected number of trials until a unique
base-field root is at most
\begin{equation}\label{eq:unique-retry}
 \mathsf{Retry}_{K,\tau}\coloneqq
 \frac{q^\tau}{1-q^{-\tau}+q^{-K}}.
\end{equation}
\end{corollary}

\begin{proof}
If some $T_y$ is not surjective, a nonzero $b$ annihilates its image, so the
transpose of $\Phi_{b,L}$ has the nonzero kernel vector $y$ and $e_L(b)<h$.  The
converse is the same argument in reverse.  Under either equivalent condition,
$r(y)=K$ and $\kappa(y)=1$ for every nonzero $y$, giving
$\eta_0=\bar\eta=(q^h-1)q^{-K}=\mu-q^{-K}$.  Substitute this value into
\eqref{eq:unique-retry-general}.
\end{proof}

Full cross-polar rank gives the slightly sharper factor
$\mathsf{Retry}_{K,\tau}$, but the headline does not rely on it.  In the
running Level-I example, $h=40$, $K=50$, and $\tau=10$.  The weaker global
condition $r_{\min}(L)\ge41$ gives at most
$19^{10}\cdot19/18=2^{42.56}$ target--slice trials.  The former full-rank
assumption required rank $50$ in every nonzero direction; the new condition
requires only rank $41$.

Equation~\eqref{eq:eta-spectrum} distinguishes sparse rank defects from a
pervasive loss.  If at most $T$ nonzero coefficient vectors $b$ are defective
and each has defect $h-e_L(b)\le\Delta$, then
\begin{equation}\label{eq:defect-spectrum}
 \bar\eta\le
 \mu-q^{-K}+Tq^{-K}(q^\Delta-1).
\end{equation}
If defects are counted projectively, $T$ in this formula must include all
$q-1$ nonzero scalar multiples of each projective direction.  Under full
cross-polar rank and $q=19$, \eqref{eq:conditional-multiple} is below
$2^{-43.47}$ for $\tau=10$ and below $2^{-60.47}$ for $\tau=14$; the
corresponding total-variation bounds are below $2^{-42.47}$ and $2^{-59.47}$.


\begin{proposition}[Exact averaging over random $A$-linear slices]
\label{prop:random-A-slice}
Let $A\cong\F_{q^\ell}$, let $H\subseteq V$ be an $A$-linear space of
$A$-dimension $D$, and choose $L$ uniformly from the $A$-linear subspaces of
$H$ having $A$-dimension $s$.  Define $r(y)=\rank T_y$ using the full residual
domain $V$.  Then
\begin{equation}\label{eq:random-slice-average}
 \mathbb E_L[\bar\eta_L]
 =\frac{q^{\ell s}-1}{q^{\ell D}-1}
   \sum_{y\in H\setminus\{0\}}q^{-r(y)}.
\end{equation}
If a fresh $L$ is sampled for each trial and the right-hand side is
$\eta_{\rm av}<1$, then the average probability of a unique base-field root is
at least $q^{\ell s-K}(1-\eta_{\rm av})$.
\end{proposition}

\begin{proof}
Every fixed nonzero $y\in H$ belongs to the same fraction of the
$s$-dimensional $A$-subspaces, namely
$(q^{\ell s}-1)/(q^{\ell D}-1)$.  Average the definition of $\bar\eta_L$
term by term.  The unique-root statement follows by averaging
\eqref{eq:singleton-coset} over $L$.
\end{proof}

This identity gives a second exact certification route.  A small exceptional
rank-defect locus in the stable kernel can be harmless even when it prevents a
worst-direction theorem for one fixed slice.  The proposition concerns retry
probability only; a randomly selected slice must still pass the public
restricted-rank and rejection preflights before its special-XL cost model is
used.

\begin{corollary}[Most random slices are good]\label{cor:random-slice-markov}
In the setting of \Cref{prop:random-A-slice}, for every $\varepsilon>0$,
\begin{equation}\label{eq:random-slice-markov}
 \Pr_L[\bar\eta_L>\varepsilon]
 \le \frac{\eta_{\rm av}}{\varepsilon}.
\end{equation}
In particular, if $\eta_{\rm av}\le\delta\varepsilon$, at least a
$1-\delta$ fraction of the $A$-linear slices satisfy
$\bar\eta_L\le\varepsilon$.
\end{corollary}

\begin{proof}
This is Markov's inequality applied to the nonnegative random variable
$\bar\eta_L$ and the exact expectation in
\eqref{eq:random-slice-average}.
\end{proof}

\subsection{A random-XOF theorem for the structured rank spectrum}
\label{sec:seeded-spectrum}

In the pinned $q=19$ public expansion, each generated field symbol is a
distinct XOF output byte reduced modulo $19$~\cite{SNOVA23}.  In the
random-XOF model these bytes are independent and uniform; hence every symbol
has maximum point probability
\begin{equation}\label{eq:rho-mod19}
 \rho=\max_{a\in\F_{19}}\Pr[X\bmod19=a]=\frac{14}{256}.
\end{equation}
The slight modulo bias is retained rather than replaced by a uniform-field
heuristic.  The theorem below is information-theoretic under this idealization;
it does not assert literal independence among bytes generated from one fixed
SHAKE seed.

Write $a\mathbin\odot b$ for the class of $a\otimes b$ in a symmetric square.

\begin{lemma}[Cross-tensor injection]\label{lem:cross-tensor-injection}
Let $U$ be an $A$-vector space, let $0\ne y\in U$, and let
$0\ne c\in\Sym^2_{\F_q}(A)$.  If
$c=\sum_\nu a_\nu\mathbin\odot b_\nu$, define the linear map
\begin{equation}\label{eq:Xi-map}
 \Xi_{c,y}(x)=
 \sum_\nu\bigl((a_\nu y)\mathbin\odot(b_\nu x)
              +(b_\nu y)\mathbin\odot(a_\nu x)\bigr)
 \in\Sym^2_{\F_q}(U).
\end{equation}
When $q$ is odd,
\begin{equation}\label{eq:Xi-kernel}
 \ker\Xi_{c,y}\subseteq Ay.
\end{equation}
\end{lemma}

\begin{proof}
If $x\notin Ay$, the $A$-lines $Ay$ and $Ax$ are disjoint.  Project
$\Sym^2(U)$ onto the cross summand between these two lines and identify that
summand with $Ay\otimes Ax$.  Under the identifications $Ay\cong A\cong Ax$,
the projection of \eqref{eq:Xi-map} is the ordinary symmetrization of $c$.
Symmetrization is injective on $\Sym^2(A)$ in odd characteristic: diagonal
tensors acquire the nonzero factor $2$, and the off-diagonal symmetric tensors
remain distinct.  Since $c\ne0$, the projection is nonzero.  Thus
$\Xi_{c,y}(x)\ne0$ whenever $x\notin Ay$.
\end{proof}

\begin{theorem}[Decoupled seeded-spectrum bound]\label{thm:seeded-spectrum}
Assume the random-XOF model above and the reserved-line construction of
\Cref{sec:seeded-decoupling}.  Suppose
$\rank Q_R=K$, $\rank C_{R,\mathbf v}=M-K$, and the
$\F_q$-linear slice $L\subseteq\ker C_{R,\mathbf v}$ has base-field dimension
$h$ and passes \eqref{eq:projection-preflight}.  Put
\begin{equation}\label{eq:seeded-D0-t}
 D_0=V'\cap\ker C_{R,\mathbf v},
 \qquad
 t_0=\dim_{\F_q}D_0-\ell,
 \qquad
 t=\ell(v-2)-(M-K).
\end{equation}
Then $t_0\ge t$, and, over the fresh internal coefficients
$P_{11,i}|_{V'\times V'}$,
\begin{align}
 \Pr[b\in\ker T_y^\transpose]&\le\rho^t
 &&(0\ne y\in L,\ 0\ne b\in\F_q^K),
 \label{eq:seeded-direction-output}\\
 \mathbb E[\bar\eta_L]
 &\le(q^h-1)(q^{-K}+\rho^t),
 \label{eq:seeded-eta-mean}\\
 \Pr[\bar\eta_L>\varepsilon]
 &\le\frac{(q^h-1)(q^{-K}+\rho^t)}{\varepsilon}
 &&(\varepsilon>0).
 \label{eq:seeded-eta-tail}
\end{align}
More sharply, after subtracting the unavoidable full-rank baseline,
\begin{equation}\label{eq:seeded-excess-tail}
 \Pr\!\left[\bar\eta_L>(q^h-1)q^{-K}+\varepsilon\right]
 \le\frac{(q^h-1)(1-q^{-K})\rho^t}{\varepsilon}
 \qquad(\varepsilon>0).
\end{equation}
In particular,
\begin{equation}\label{eq:seeded-headline-tail}
 \Pr[\bar\eta_L>q^{-1}]
 \le q(q^h-1)(q^{-K}+\rho^t).
\end{equation}
If a list of at most $J$ candidates is fixed from the cross data and the first
candidate passing all later full-key preflights is selected, the probability
that the selected candidate violates either tail is at most $J$ times the
corresponding right-hand side.  The slice need not be $A$-linear;
$A$-linearity is required later only by the multigraded XL and eigenblock
solvers.
\end{theorem}

\begin{proof}
The dimension bound follows from
\[
 \dim D_0\ge\dim V'-(M-K)=\ell(v-1)-(M-K).
\]
Fix $y\ne0$ and $b\ne0$.  By \Cref{lem:canonical-residual-coordinates}, the canonical residual map is
$z_{\rm sym}$ plus an affine term.  Therefore the nonzero residual output
coefficient $b$ is itself a nonzero tuple
$(c_1,\ldots,c_{m_1})\in\bigoplus_i\Sym^2(A)$ in the fixed quotient-feature
basis.  Choose $i$ with $c_i\ne0$.
The projection $\bar y=\pi_{V'}(y)$ is nonzero by
\eqref{eq:projection-preflight}.

For $x\in D_0$, the coefficient vector of the fresh symmetric matrix
$P_{11,i}|_{V'\times V'}$ in the scalar polar equation
$b\cdot T_y(x)=0$ is precisely the tensor
$\Xi_{c_i,\bar y}(x)$.  Contributions involving the $W$ component of $y$ or
the public-oil component use only the already fixed cross data; $P_{22}$ does
not occur because $x$ is supported on $V'$.  By
\Cref{lem:cross-tensor-injection}, the kernel of this coefficient map on $D_0$
has dimension at most $\ell$.  Its rank is therefore at least
$t_0\ge t$.

Select $t$ test vectors in $D_0$ whose coefficient vectors are independent.
The event $b\in\ker T_y^\transpose$ implies the resulting rank-$t$ affine
system in the independent upper-triangular $P_{11,i}|_{V'\times V'}$ symbols.
After row reduction, condition on all nonpivot symbols.  Each pivot symbol is
forced to one value and has probability at most $\rho$; independence gives
probability at most $\rho^t$.  This proves
\eqref{eq:seeded-direction-output}.

Finally,
\[
 q^{-\rank T_y}=q^{-K}|\ker T_y^\transpose|
 =q^{-K}\left(1+\sum_{b\ne0}
       \mathbf1_{\{b\in\ker T_y^\transpose\}}\right).
\]
Taking expectation and summing the indicators gives both
\[
 \mathbb E[q^{-\rank T_y}]\le q^{-K}+\rho^t
 \quad\text{and}\quad
 \mathbb E[q^{-\rank T_y}-q^{-K}]
 \le(1-q^{-K})\rho^t.
\]
Sum over the $q^h-1$ nonzero $y\in L$ to obtain
\eqref{eq:seeded-eta-mean}.  Markov's inequality applied first to
$\bar\eta_L$ and then to its nonnegative excess above
$(q^h-1)q^{-K}$ gives \eqref{eq:seeded-eta-tail} and
\eqref{eq:seeded-excess-tail}.  For a bounded candidate list, a bad selected
candidate implies that at least one member of the fixed list is bad, so the
union bound applies even when the later exact preflights use the full key.
\end{proof}

\begin{corollary}[Numerical bad-spectrum bounds]\label{cor:seeded-numerics}
For the six structured Version~2.3 rows, the lower bound $t$ and the resulting
probabilities are as follows.  A slash in the $t$ column gives the values for
the two shapes at that level.  ``Compact'' refers to
\Cref{tab:explicit-stream}; ``robust'' refers to
\Cref{tab:robust-explicit}.

\begin{center}
\small
\begin{tabular}{@{}c c c r c r@{}}
\toprule
Level & $t$ & compact $(s,\tau)$ &
$\log_2\Pr[\bar\eta>1/19]$ & robust $(s,\tau)$ &
$\log_2\Pr[\bar\eta>1/2]$\\
\midrule
I   & $74/74$   & $(9,14)$  & $<-55.22$  & $(10,10)$ & $<-41.48$\\
III & $110/114$ & $(10,30)$ & $<-123.19$ & $(11,26)$ & $<-109.45$\\
V   & $138/146$ & $(11,46)$ & $<-191.16$ & $(12,42)$ & $<-177.41$\\
\bottomrule
\end{tabular}
\end{center}
For the compact Level-I row, even a cross-data-generated list of $2^{20}$
candidates leaves the union bound below $2^{-35.22}$.  These probabilities are
under that idealized key distribution.  They do not replace a deterministic certificate when a
claim is made about one already fixed public seed.
\end{corollary}

\begin{proof}
Substitute $q=19$, $\rho=14/256$, and the public dimensions into
\eqref{eq:seeded-eta-tail}.  The exact high-precision evaluations are included
in the companion artifact.
\end{proof}

\begin{corollary}[Full-dimensional fresh ordinary slices]
\label{cor:full-ordinary-slices}
For an $\ell=4$ row, let
$D_0=V'\cap\ker C_{R,\mathbf v}$.  The rank of $C_{R,\mathbf v}$ gives the
deterministic lower bound
\begin{equation}\label{eq:D0-lower-bound}
 \dim_{\F_{19}}D_0
 \ge 4(v-1)-(M-K).
\end{equation}
For all six public rows this lower bound is at least $K$.  Hence the attacker
can select, from cross data alone, a $K$-dimensional ordinary
$\F_{19}$-linear slice $L\subseteq D_0$.  Put
\[
 B_K=1-19^{-K}+\frac1{19}.
\]
In the random-XOF idealization,
$\bar\eta_L\le B_K$ except with the probabilities in the final column below.
The bound uses \eqref{eq:seeded-excess-tail} with $\varepsilon=1/19$.

\begin{center}
\small
\setlength{\tabcolsep}{3.6pt}
\begin{tabular}{@{}c l r r r r@{}}
\toprule
Level & parameters & $K$ & $\dim D_0\ge$ & $t\ge$ &
$\log_2\Pr[\bar\eta_L>B_K]$\\
\midrule
I   & $(28,5,4,4)$  & 50 & 78  & 74  & $<-93.61$\\
I   & $(28,4,4,5)$  & 50 & 78  & 74  & $<-93.61$\\
III & $(40,7,4,4)$  & 70 & 114 & 110 & $<-159.59$\\
III & $(38,5,4,5)$  & 70 & 118 & 114 & $<-176.36$\\
V   & $(50,9,4,4)$  & 90 & 142 & 138 & $<-192.02$\\
V   & $(52,6,4,6)$  & 90 & 150 & 146 & $<-225.56$\\
\bottomrule
\end{tabular}
\end{center}

For the worst accepted density at each level,
\Cref{thm:accepted-nonsingular-roots} therefore gives an accepted
full-Jacobian root in one target--coset trial with probability at least
$0.43163$, $0.43153$, and $0.43145$, respectively.  The corresponding
expected trial counts are below $2.318$.  These ordinary slices are used by the
complete-square homotopy route; they need not preserve the $A$-linear grading
required by special XL.
\end{corollary}

\begin{proof}
The intersection inequality
$\dim(V'\cap\ker C)\ge\dim V'-\rank C$ gives
\eqref{eq:D0-lower-bound}.  The table substitutes the six public dimensions,
uses $t=\dim D_0-4$, and applies \eqref{eq:seeded-excess-tail}.  Because
$h=K$, the deterministic baseline is
$(19^K-1)19^{-K}=1-19^{-K}$.  The final root probabilities follow from
\eqref{eq:accepted-nonsingular-success} with $\bar\eta_L=B_K$ and the exact
square-row acceptance lower bounds of \Cref{sec:rejection}.
\end{proof}

\begin{proposition}[Exact projective rank-histogram certificate]
\label{prop:projective-rank-certificate}
Choose an $\F_q$-basis $y_0,\ldots,y_{h-1}$ of a fixed slice $L$ and write
\[
 T(Y)=\sum_{i=0}^{h-1}Y_iT_{y_i}\in
 \F_q[Y_0,\ldots,Y_{h-1}]^{K\times N}.
\]
For $0\le j<h$, use the disjoint normalized projective chart
$Y_0=\cdots=Y_{j-1}=0$, $Y_j=1$, put
$R_j=\F_q[Y_{j+1},\ldots,Y_{h-1}]$, and let $T^{(j)}$ be the resulting
matrix over $R_j$.  For $t\ge0$, define
\begin{equation}\label{eq:rank-chart-ideal}
 I_{j,t}=\left\langle
  Y_i^q-Y_i\ (i>j),\quad
  \textnormal{all $(t+1)\times(t+1)$ minors of }T^{(j)}
 \right\rangle\subseteq R_j.
\end{equation}
Set $n_{j,\le t}=\dim_{\F_q}(R_j/I_{j,t})$, put
$n_{j,\le-1}=0$, and define
\begin{equation}\label{eq:projective-rank-count}
 N_t=\sum_{j=0}^{h-1}
       \bigl(n_{j,\le t}-n_{j,\le t-1}\bigr).
\end{equation}
Then $N_t$ is exactly the number of projective directions $[y]\in
\mathbb P(L)(\F_q)$ with $\rank T_y=t$, and
\begin{equation}\label{eq:projective-eta-histogram}
 \bar\eta_L=(q-1)\sum_tN_tq^{-t}.
\end{equation}
Consequently, the quotient dimensions in
\eqref{eq:projective-rank-count} give a finite, independently checkable
certificate of the exact aggregate collision mass.  In particular,
$\rank T_y\ge r+1$ for every $y\ne0$ if and only if every
$I_{j,r}$ is the unit ideal.

The same construction applied to the polar matrix pencil indexed by
$b\in\F_q^K$ gives exact projective counts $N_t^{\rm pol}$ and hence
\begin{equation}\label{eq:projective-polar-histograms}
 \Theta_{\rm pol}=(q-1)\sum_tN_t^{\rm pol}q^{-t},
 \qquad
 \Psi_{\rm pol}=(q-1)\sum_tN_t^{\rm pol}q^{-t/2}.
\end{equation}
A certificate may consist of reduced Gr\"obner bases and their standard
monomial sets, or of polynomial reduction transcripts proving the unit-ideal
claims and quotient dimensions; verification needs no probabilistic sampling.
\end{proposition}

\begin{proof}
The chosen charts contain exactly one representative of each projective
$\F_q$-direction.  The field equations identify
\[
 R_j/\langle Y_i^q-Y_i:i>j\rangle
 \cong \prod_{a\in\F_q^{h-j-1}}\F_q
\]
by evaluation.  Every ideal in this product of fields is radical, and its
quotient dimension is the number of coordinates on which all of its generators
vanish.  The minors in \eqref{eq:rank-chart-ideal} vanish exactly when
$\rank T^{(j)}(a)\le t$.  Thus $n_{j,\le t}$ counts the rank-at-most-$t$
points in chart $j$, and differencing gives \eqref{eq:projective-rank-count}.
Every projective direction has $q-1$ nonzero scalar representatives and rank
is invariant under base-field scaling, which proves
\eqref{eq:projective-eta-histogram}.  The polar formulas are identical.
\end{proof}

\begin{proposition}[A fully priced exhaustive spectrum computation]
\label{prop:exhaustive-rank-cost}
Let $P_h=(q^h-1)/(q-1)$ be the number of projective directions in an
$h$-dimensional slice, and suppose $K\le N$.  The exact histogram in
\Cref{prop:projective-rank-certificate} can also be computed directly, without
Gr\"obner bases, in at most
\begin{equation}\label{eq:exhaustive-rank-cost}
 4P_hK^2N
\end{equation}
base-field operations and $O(hKN)$ stored base-field elements.  The algorithm
enumerates the disjoint normalized charts, updates $T_y$ along a $q$-ary Gray
code, row-reduces each $K\times N$ matrix, and accumulates its rank.

For the model-state profiles in \Cref{tab:xl-optima}, charging the paper's
$b_1=22.29$ bit operations per base-field operation gives the following
worst-case exponents, rounded upward to the displayed precision:
\begin{center}
\begin{tabular}{@{}c c c c c@{}}
\toprule
Level & $(h,K,N)$ & exact spectrum & structured solve & combined\\
\midrule
I   & $(32,50,102)$ & $156.21$ & $140.13$ & $156.21$\\
III & $(36,70,146)$ & $174.69$ & $205.47$ & $205.47$\\
V   & $(44,90,182)$ & $209.71$ & $262.53$ & $262.53$\\
\bottomrule
\end{tabular}
\end{center}
Here ``combined'' is the logarithm of the sum of the certification and solve
work.  Thus, for the displayed Level-III and Level-V low-state attacks, an
exact favorable value of $\bar\eta_L$ is not an unpriced production step: the
complete enumeration is dominated by the modeled solve.  At Level I, direct
enumeration is too expensive for a below-$143$ claim, so the structured
certificate of \Cref{prop:projective-rank-certificate} or another proof remains
necessary.  The computation can of course return an unfavorable spectrum; the
proposition prices certification, not the probability that a randomly selected
slice passes it.
\end{proposition}

\begin{proof}
Precompute the $h$ coefficient matrices $T_{y_i}$.  A standard $q$-ary Gray
code changes one coordinate at a time, so forming the next matrix costs at most
$KN$ additions.  Gaussian elimination on a $K\times N$ matrix uses fewer than
$K^2N$ multiply-add pairs and $KN$ normalizations.  Counting each addition,
multiplication, or division separately and allowing slack for chart changes
gives \eqref{eq:exhaustive-rank-cost}.  Substitution of the displayed
parameters gives the table; the machine-readable artifact performs the integer
and logarithmic calculations.
\end{proof}

\begin{corollary}[Full-space target availability]\label{cor:full-space}
Let $g:\F_q^N\to\F_q^K$ have degree at most two.  Let $r_*$ be the
minimum polar rank of a nonzero output combination, and put
$\alpha=1-q^{-K}$.  For a uniform target,
\begin{equation}\label{eq:full-space-success}
 \Pr[g^{-1}(z)\ne\varnothing]
 \ge \frac{1}{1+\alpha q^{K-r_*}}.
\end{equation}
Hence the expected number of target trials is at most
$1+\alpha q^{K-r_*}$.  If $r_*\ge2K$, every target has a root.
\end{corollary}

\begin{proof}
See \Cref{app:proof-full-space}.
\end{proof}

\begin{corollary}[Aggregate full-space rank spectra]
\label{cor:aggregate-full-space}
For $b\in\F_q^K\setminus\{0\}$, let $\rho(b)$ be the polar rank of the scalar
quadratic $b\cdot g$, and put
\begin{equation}\label{eq:aggregate-polar-spectra}
 \Theta_{\rm pol}(g)=\sum_{b\ne0}q^{-\rho(b)},
 \qquad
 \Psi_{\rm pol}(g)=\sum_{b\ne0}q^{-\rho(b)/2}.
\end{equation}
For a uniform target $z$,
\begin{equation}\label{eq:aggregate-full-space-success}
 \Pr[g^{-1}(z)\ne\varnothing]\ge\frac{1}{1+\Theta_{\rm pol}(g)}.
\end{equation}
More generally, if $\mathcal X\subseteq\F_q^N$ has density $a$, then
\begin{equation}\label{eq:aggregate-accepted-success}
 \Pr[g^{-1}(z)\cap\mathcal X\ne\varnothing]
 \ge\frac{a}{1+\Psi_{\rm pol}(g)}.
\end{equation}
Thus the expected target counts are at most $1+\Theta_{\rm pol}(g)$ and
$(1+\Psi_{\rm pol}(g))/a$, respectively.  These bounds depend on the complete
rank spectra, not their minima; a small exceptional low-rank locus can be
charged exactly rather than controlling the whole pencil.
\end{corollary}

\begin{proof}
Take $L=V$ in \Cref{thm:coset-success}.  Then $e_V(b)=\rho(b)$, so
\eqref{eq:coset-success} is exactly
\eqref{eq:aggregate-full-space-success}.  For the second statement, the
quadratic character-sum bound gives, for every target,
\[
 |g^{-1}(z)|
 \le q^{N-K}\left(1+\sum_{b\ne0}q^{-\rho(b)/2}\right)
 =q^{N-K}(1+\Psi_{\rm pol}(g)).
\]
Double-counting the $a q^N$ points of $\mathcal X$ by their images shows that
at least $a q^K/(1+\Psi_{\rm pol}(g))$ targets have a preimage in
$\mathcal X$, proving \eqref{eq:aggregate-accepted-success}.
\end{proof}

\begin{corollary}[Small expected-target cushions]\label{cor:target-cushions}
For every integer $c\ge1$,
\begin{align}
 \Theta_{\rm pol}(g)\le q^{-c}
 &\quad\Longrightarrow\quad
 \mathbb E[\textnormal{uniform targets to a nonempty fiber}]
 \le1+q^{-c},\label{eq:full-space-spectrum-cushion}\\
 \Psi_{\rm pol}(g)\le q^{-c}
 &\quad\Longrightarrow\quad
 \mathbb E[\textnormal{uniform targets to an accepted root}]
 \le \frac{1+q^{-c}}{a_n}.
 \label{eq:filter-spectrum-cushion}
\end{align}
The convenient minimum-rank conditions
\begin{align}
 r_*\ge K+c
 &\quad\Longrightarrow\quad \Theta_{\rm pol}(g)<q^{-c},
 \label{eq:full-space-small-cushion}\\
 r_*\ge2K+2c
 &\quad\Longrightarrow\quad \Psi_{\rm pol}(g)<q^{-c}
 \label{eq:filter-small-cushion}
\end{align}
are sufficient, but not necessary.
In particular, at $q=19$ and $c=1$, either aggregate condition adds less than
$\log_2(20/19)<0.075$ bits.  For the Level-V filter,
$n=128$, $K=96$, and
\[
 1-a_{128}\le\binom{128}{33}19^{-33}<2^{-38.4},
\]
so $\Psi_{\rm pol}(g)\le19^{-1}$---in particular, $r_*\ge194$---makes the
total expected-target overhead less than $0.075$ bits as well.
\end{corollary}

\begin{proof}
Apply \Cref{cor:aggregate-full-space}.  Moreover,
\[
 \Theta_{\rm pol}(g)
 \le(q^K-1)q^{-r_*}<q^{K-r_*},
 \qquad
 \Psi_{\rm pol}(g)
 \le(q^K-1)q^{-r_*/2}<q^{K-r_*/2},
\]
which proves the minimum-rank implications.  The final estimate is the union
bound for at least $33$ zero hyperplane coordinates among $128$ independent
blocks.
\end{proof}

For a structured $\ell=4$ lift, an invertible output transformation separates
the $M-K$ affine coordinates from the remaining $K$ coordinates of a uniform
target.  Conditioning on the affine coordinates fixes a base point; the other
coordinates remain uniform, and translation does not change the polar maps.
We take $V=\ker C_{R,\mathbf v}$ and let $L$ be the selected stable slice, of dimension
$h=\ell s=K-\tau$.  The root ledger uses the aggregate bound
$\bar\eta_L\le q^{-1}$, so \eqref{eq:unique-retry-general} bounds the
unique-root retry factor by $q^\tau/(1-q^{-1})$.  The minimum-rank condition
$r_{\min}(L)\ge h+1$ is sufficient, and full cross-polar rank would improve
the exponent by less than $0.079$ bits.  The bound can now be justified in
three distinct ways: the random-XOF idealized-expansion theorem
\Cref{thm:seeded-spectrum}, a fixed-key projective histogram, or the random
$A$-slice average of \Cref{prop:random-A-slice}.  Reaching the degree at which
the solver exposes a finite quotient remains a separate condition.

For the fixed-skew $\ell=2$ lift, take $L=V$ to be the complete residual space.
By \Cref{prop:fixed-skew}, every root is accepted.  The headline uses
$\Theta_{\rm pol}\le19^{-1}$ from \Cref{cor:aggregate-full-space}, which
adds less than $0.075$ bits of expected target search.  The minimum-rank values
$r_*\ge K+1$, namely $49$, $73$, and $97$ at Levels I, III, and V, are
sufficient but not necessary.
The stronger condition $r_*\ge2K$ makes every target solvable.  Under the AXN
conversion, the nominal arithmetic gaps remain positive for global ranks at least
46, 68, and 89.  These budget and all-target minimum ranks concern the converted operation
model, not a complete solver circuit; the sampled evidence in
\Cref{sec:evidence} is not a certificate.  The Level-V filter attack instead
uses the accepted-root version of \Cref{cor:target-cushions}.

\begin{theorem}[End-to-end Las Vegas wrapper]
\label{thm:las-vegas-wrapper}
Fix a public key and an attacker-chosen relation that pass the exact affine and
rejection preflights.  Suppose a residual solver has the following certified
behavior: when it reports success, it returns a finite candidate set containing
all base-field roots of that residual system; every returned point is then
checked by the original public verifier.  Let $\sigma\in(0,1]$ be a lower bound
on the solver's conditional probability of producing such a completeness
certificate, conditioned on the root event stated below.

\begin{enumerate}[label=(\alph*),itemsep=3pt]
\item For a structured slice $L$ of dimension $h$, let $a$ be a
      lower bound on the density of assignments accepted by the original
      verifier in the full common-column space.  If $\bar\eta_L<a$, the
      expected number of target--slice solver trials before a certified valid
      signature is at most
      \begin{equation}\label{eq:wrapper-structured}
        \frac{q^{K-h}}{(a-\bar\eta_L)\sigma}.
      \end{equation}
      Here the conditioning event for $\sigma$ is that the trial has exactly
      one base-field root and that root is accepted.  Take $a=1$ for the
      rectangular rows and $a=(1-q^{-4})^n$ for the square $\ell=4$ rows.
\item For a full-space fixed-skew lift, if the aggregate polar spectrum is
      $\Theta_{\rm pol}$, the expected number of target--solver trials is at
      most
      \begin{equation}\label{eq:wrapper-fixed-skew}
        \frac{1+\Theta_{\rm pol}}{\sigma}.
      \end{equation}
      Here the conditioning event is that the target fiber is nonempty.
\item For a filtered lift whose accepted assignments have density $a$, if the
      aggregate square-root spectrum is $\Psi_{\rm pol}$, the expected number
      of target--solver trials is at most
      \begin{equation}\label{eq:wrapper-filtered}
        \frac{1+\Psi_{\rm pol}}{a\sigma}.
      \end{equation}
      Here the conditioning event is that the target has an accepted root.
\end{enumerate}
In every case the algorithm is Las Vegas: it may restart, but it never returns
an invalid signature.
\end{theorem}

\begin{proof}
For the structured case, the invertible target split of
\Cref{lem:canonical-residual-coordinates} makes the affine right-hand side and
residual target independent and uniform.  \Cref{thm:accepted-singleton} then
gives an accepted unique-root probability at least
$q^{h-K}(a-\bar\eta_L)$.  Multiplying by the conditional certificate
probability $\sigma$ and inverting gives \eqref{eq:wrapper-structured}.  The other two root-event probabilities are at
least $1/(1+\Theta_{\rm pol})$ and $a/(1+\Psi_{\rm pol})$ by
\Cref{cor:aggregate-full-space}; the same calculation proves
\Cref{eq:wrapper-fixed-skew,eq:wrapper-filtered}.  Completeness of a reported
candidate set prevents a missed root from being mistaken for a proof of
infeasibility, and the original verifier removes all algebraic or
reconstruction artifacts.
\end{proof}

\paragraph{Where the numerical model enters.}
The compact special-XL rows set $\sigma=1$ at the selected degree.  The robust
explicit rows set only $\sigma=1/8$ and charge the resulting three bits.  A
transcript satisfying \Cref{thm:finite-quotient-recovery} certifies success on
an individual trial.  For any other measured or proved value, the correction
is transparent: add $\log_2(1/\sigma)$ bits.  The seeded root-spectrum theorem
and the production solver-certificate rate are therefore separate inputs and
cannot be silently conflated.

\paragraph{Certified attack procedure.}
The forger (i) runs the public rank and rejection preflights; (ii) samples an
exactly uniform official target as described next; (iii) forms the exact
residual system and, for the structured rows, a target--slice trial; (iv) runs
the selected solver and accepts a solver claim only with a completeness
certificate; and (v) reconstructs the full signature and invokes the
unmodified verifier.  A failed preflight, empty fiber, absent certificate, or
failed verifier check causes a restart, never an unchecked output.

\subsection{Sampling exactly uniform targets with the official conversion}

The probability theorems above assume a uniform target in $\F_q^K$.  Reducing
binary hash chunks modulo $19$ introduces a small bias, so uniformity cannot
simply be assumed.  The attacker removes it without changing the verifier: it
repeatedly chooses a message--salt input, runs the official hash-to-field
conversion, and keeps the input only when each underlying binary chunk lies in
a range that makes the resulting base-$19$ digits exactly uniform.

Concretely, the official $q=19$ conversion reads at most fifteen base-19 digits
from each 64-bit chunk.  If a uniform $b$-bit chunk $X$ is to supply $a$
digits, the attacker accepts that message--salt candidate only if
\begin{equation}\label{eq:uniform-chunk}
 X<\left\lfloor\frac{2^b}{q^a}\right\rfloor q^a.
\end{equation}
Conditioned on this event, $X\bmod q^a$ is uniform, so the extracted
digits are jointly uniform.  Applying \eqref{eq:uniform-chunk} independently
to every chunk makes the retained official target exactly uniform.  The
invertible output-coordinate change used in \Cref{sec:affine-columns}
preserves uniformity.  Conditioning on the values of the $M-K$ affine-output coordinates---or on
their passing values in a filtered lift---then leaves the remaining $K$
residual target coordinates exactly uniform, which is the distribution
assumed in the root theorems.  This is
ordinary rejection sampling performed by the forger over message--salt
candidates; the verifier continues to apply the unmodified official
conversion.  Across all Version~2.3 lengths the acceptance probability is at
least $0.152462$, or fewer than 6.6 candidates on average.

The Level-V $\ell=2$ filter needs about $19^{32}>2^{128}$ raw targets before
the $M-K=32$ consistency conditions pass, so the 128-bit salt alone does not
supply enough distinct inputs.  An existential forger may vary a chosen-message
prefix as well as the salt.  Distinct prefix--salt pairs are distinct random-oracle
inputs and therefore give independent targets in the model.  The
artifact's deliberately loose target-generation envelope is below $2^{32}$
gates per target retained by the unbiased hash-chunk sampler (the AXN charge
is no larger than the NAND envelope used there).  Even after multiplication by the $19^{32}$ filtering factor,
this term is more than 70 exponent bits below the cost-driving MQ solve.

\section{Concrete Costs}\label{sec:complexity}

The algebraic reduction determines the residual MQ systems exactly.  Turning
those systems into numerical security estimates requires a separate cost
model.  We report that model in layers so that a reader can see which numbers
are structural, which come from an MQ solver estimate, and which are merely a
change of arithmetic units.

\begin{table}[H]
\centering
\caption{How to interpret the cost accounting.}
\label{tab:cost-layers}
\small
\renewcommand{\arraystretch}{1.12}
\begin{tabularx}{0.96\textwidth}{@{}>{\bfseries}p{0.23\textwidth}X@{}}
\toprule
Layer & What it means\\
\midrule
Residual system & Exact numbers of equations, variables, blocks, and public
rank conditions after the reduction.\\
Solver model & Published generic-MQ or special-XL formula for the dominant
finite-field work at those dimensions.\\
AXN conversion & Secondary sensitivity calculation replacing the solver
formula's coarse field-operation factor by an explicit two-input
AND/XOR/XNOR circuit for one signed multiply-accumulate.\\
Not included & A complete implementation of sparse matrix construction,
pivoting and control, memory addressing and traffic, communication, or an
optimized end-to-end attack netlist.\\
\bottomrule
\end{tabularx}
\end{table}

A cost exponent $c$ means modeled work proportional to $2^c$ in the stated
unit.  A difference of ten exponent bits is a factor of about $2^{10}$, not ten
percent.  Because the solver formulas count dominant arithmetic rather than a
complete machine, any difference from a NIST reference is only a
\emph{nominal arithmetic gap}; it is not an implementation margin or a
category claim.

\subsection{Cost units and circuit-unit sensitivity}\label{sec:units}

The cited attacks first count finite-field operations and then apply coarse
bit-cost factors~\cite{Beullens2024SNOVA,Cabarcas2025SNOVA}.  NIST lists
$2^{143}$, $2^{207}$, and $2^{272}$ as nominal classical-gate references, while
also requiring category comparisons across every metric it considers relevant
to practical security~\cite{NISTPQCCriteria}.  We retain each attack's
field-operation count and replace only its dominant field factor by explicit
circuits.  The resulting figures sharpen one metric; they do not by themselves
establish a category break.  Our primary AXN basis assigns unit cost to two-input
AND, XOR, and XNOR gates, with free constants, wires, fan-out, and structural
sharing.  This is the gate family used by the concrete AES-128 circuit on
NIST's circuit list~\cite{NISTCircuitComplexity}.  A two-input-NAND conversion
provides a separate basis-sensitivity check.

For generic MQ, the published factor is
\begin{equation}\label{eq:published-f19-factor}
 b_1=(\log_2 19)^2+\log_2 19=22.29.
\end{equation}
Using canonical five-bit representatives and exact Barrett reduction, the
validated circuits for one signed multiply-accumulate, $c+ab$ or $c-ab$, over
$\F_{19}$ use 402 and 441 AXN gates, respectively; their NAND translations use 781 and 835 gates.  To avoid
assuming a particular sign schedule in elimination, we charge
\begin{equation}\label{eq:axn-f19-factor}
 g_1^{\mathrm{AXN}}=441,
 \qquad g_1^{\mathrm{NAND}}=835.
\end{equation}
Both signs are checked exhaustively on all $19^3$ valid input triples by the
generator and by an independent scalar evaluator.

For special XL, we work in the isomorphic sparse basis
$\F_{19}[t]/(t^4-t-1)$.  In this basis, the element
$13+13t+15t^2+2t^3$ is a root of the official defining polynomial; evaluating
its powers supplies the explicit basis change to $\F_{19}[S]$.  Two-level
Karatsuba uses nine base-field multiplications, and $t^4=t+1$ makes reduction
additive.  The validated circuits for $c\mathbin{\pm}ab$ over $\F_{19^4}$ use
6,077 and 6,081 AXN gates, or 11,086 and 11,102 NAND gates.  We therefore
charge
\begin{equation}\label{eq:axn-f194-factor}
 g_4^{\mathrm{AXN}}=6{,}081,
 \qquad g_4^{\mathrm{NAND}}=11{,}102.
\end{equation}
The artifact checks the complete scalar basis tensor, 4,096 random extension
products, the basis isomorphism, and both translated netlists.  Charging a
signed multiply-accumulate for every multiplication in the special-XL formula
is conservative.

\paragraph{Same-basis AES calibration.}
The release builds variable-key AES-128/192/256 one-pair key-filter circuits
with the NIST 113-gate S-box, the complete key schedule, encryption, and a
final ciphertext comparator.  To identify the actual key rather than merely a
key that matches one pair, surviving candidates are checked lazily on enough
independent plaintext--ciphertext pairs to leave at most $2^{-128}$ expected
wrong survivors in the ideal-cipher model: a total of two pairs for AES-128 and
a total of three for AES-192 and AES-256.  The additional average work is
negligible because only
survivors reach those checks.  Every first-pair circuit is validated on the
FIPS test vector and eight additional keys against an independent AES
implementation.

\begin{table}[tbp]
\centering
\caption{Constructive AES calibration circuits.  Full charges the entire key
space for the first-pair filter; expected search is one bit lower, and the
lazy confirmation overhead is negligible under the stated ideal-cipher model.
These circuits are not lower bounds on optimal AES implementations and do not
replace NIST's nominal references.}
\label{tab:aes-calibration}
\scriptsize
\setlength{\tabcolsep}{3.5pt}
\begin{tabular}{@{}c r r r r r r@{}}
\toprule
Level & pairs & NIST nominal & AXN gates/filter & AXN full & NAND gates/filter & NAND full\\
\midrule
I   & 2 & 143 & 30,081 & 142.88 & 106,616 & 144.70\\
III & 3 & 207 & 34,227 & 207.06 & 121,735 & 208.89\\
V   & 3 & 272 & 41,543 & 271.34 & 147,772 & 273.17\\
\bottomrule
\end{tabular}
\end{table}

The NIST-listed 28,600-gate AES-128 circuit has full-key-space exponent
$128+\log_2(28{,}600)=142.80$, numerically consistent with the nominal value
143.  Across our nine rows, the AXN arithmetic estimates are
\AXNFullAESGapRange{} bits below the constructive full-space AES circuits and
\AXNNominalHeadroomRange{} bits below the nominal NIST references.  The NAND
cross-check gives attack ranges \LoneNANDRange, \LthreeNANDRange, and
\LfiveNANDRange.

\paragraph{Scope of the conversion.}
The AXN and NAND figures price the field-arithmetic schedule counted by the
published estimators; they do not include solver control, memory traffic, or a
complete sparse-linear-algebra implementation.  Public preprocessing and
exact-uniform target generation are bounded separately, so the dominant
unresolved inputs are the global production rank spectra and the
production-scale solving models.

\subsection{Structured special XL for \texorpdfstring{$\ell=4$}{l=4}}\label{sec:special-xl}

The exact reduction supplies an $A$-linear slice of $A$-dimension $s$,
hence base-field dimension $4s$ when $\ell=4$.  After extension to the
splitting field, this becomes four blocks of $s$ variables.  For each unordered pair of blocks there are $m_1$ independent equations
of degree two: diagonal equations use one block twice, and off-diagonal
equations use two blocks once each.  Special XL records a separate degree for
each block and selects a multidegree
at which the associated Macaulay matrix is predicted to have the needed rank.
The series below is only a compact device for choosing the matrix degree and
counting the resulting monomial columns.

After adding one homogenizing variable to each block, the Hilbert series used
by the published special-XL model is
\begin{equation}\label{eq:hilbert-series}
 H_{m_1,s}(\mathbf t)=
 \frac{\prod_{a=1}^{4}(1-t_a^2)^{m_1}
       \prod_{1\le a<b\le4}(1-t_at_b)^{m_1}}
      {\prod_{a=1}^{4}(1-t_a)^{s+1}}.
\end{equation}
For $\mathbf d=(d_1,\ldots,d_4)$, put
\begin{equation}\label{eq:macaulay-size}
 M_s(\mathbf d)=\prod_{a=1}^{4}\binom{s+d_a}{d_a}.
\end{equation}
The search chooses a multidegree whose signed Hilbert coefficient is
negative.  This sign change is a rank-prediction trigger in the cited
special-XL model: a random Macaulay matrix is predicted to have full column
rank, while a system with a planted solution is predicted to have a
one-dimensional right kernel~\cite{Cabarcas2025SNOVA}.  We use the same degree
and row-reduction cost.  \Cref{thm:finite-quotient-recovery} shows that
one-dimensionality is not necessary for correctness: any certified
commuting-table finite border quotient can be solved and checked.  It does
not make a larger kernel free.  Unless extra block/nullspace work is separately
charged, the displayed scalar-Wiedemann exponent retains the published
$D=1$ prediction.  The published upper bound for one target--slice trial is
$3s^2\bigl(M_s(\mathbf d)\bigr)^2$ extension-field multiplications.  This is a
Wiedemann-style sparse-linear-algebra count, not a license to materialize a
dense $M_s(\mathbf d)\times M_s(\mathbf d)$ matrix
\cite{Beullens2024SNOVA,Cabarcas2025SNOVA,ChengEtAl2012XL}.  The next
proposition makes the corresponding state claim precise.

Cabarcas et al.\ handle an underdetermined system by specializing $k$
coordinates, where $4\mid k$, and multiplying the work by $q^k$.  Our attack
uses a different retry loop: it fixes a $4s$-dimensional slice and samples the
target and affine coset.  The headline condition $\bar\eta_L\le q^{-1}$ and
\eqref{eq:unique-retry-general} give the conservative retry factor
\[
 \mathsf{Retry}^{+}_{\tau}=\frac{q^\tau}{1-q^{-1}},
 \qquad \tau=K-4s.
\]
The modeled work is therefore
\begin{equation}\label{eq:special-xl-cost}
 \mathsf{Retry}^{+}_{\tau}\,3s^2M_s(\mathbf d)^2
\end{equation}
extension-field multiplications.  Relative to the old full-rank factor, this
adds at most $\log_2(19/18)=0.0780026\ldots$ bits.  The minimum-rank condition $r_{\min}(L)\ge4s+1$ from
\Cref{cor:rank-cushion} is one sufficient certificate.  The more general
spectral formulas \eqref{eq:eta-def} and \eqref{eq:random-slice-average} can
give a smaller factor and tolerate an exceptional low-rank locus.

The computation takes place over $\F_{19^4}$.  The published convention assigns
each extension-field multiplication
\begin{equation}\label{eq:bitop-conversion}
 b_4=2\bigl(\log_2 19^4\bigr)^2+\log_2 19^4=594.43
\end{equation}
bit operations.  The AXN arithmetic conversion replaces $b_4$ by the conservative
signed multiply-accumulate bound $g_4^{\mathrm{AXN}}$ from
\eqref{eq:axn-f194-factor}.

\begin{proposition}[Streaming the special-XL Macaulay operator]
\label{prop:streaming-macaulay}
Let $A_{\mathbf d}$ be the special-XL Macaulay matrix at multidegree
$\mathbf d$, with $M=M_s(\mathbf d)$ columns.  With the convention that a
binomial coefficient with negative lower argument is zero, its row count is
\begin{align}
 R_s(\mathbf d)
 &=m_1\sum_{a=1}^{4}
   \binom{s+d_a-2}{d_a-2}\prod_{c\ne a}\binom{s+d_c}{d_c}
 \notag\\
 &\quad{}+m_1\sum_{1\le a<b\le4}
   \binom{s+d_a-1}{d_a-1}\binom{s+d_b-1}{d_b-1}
   \prod_{c\notin\{a,b\}}\binom{s+d_c}{d_c}.
 \label{eq:macaulay-row-count}
\end{align}
Every row can be generated from the public quadratic coefficients and its
multiplier monomial using $O(s^2)$ scratch space; no matrix entry table is
needed.

Put $R=R_s(\mathbf d)$.  Let $E/\F_{19^4}$ be a finite extension, choose
the entries of a random diagonal $\Delta\in E^{R\times R}$ independently and
uniformly, and put
\begin{equation}\label{eq:streamed-normal-operator}
 B=A_{\mathbf d}^{\transpose}\Delta A_{\mathbf d}\in E^{M\times M}.
\end{equation}
A product $Bx$ can be evaluated by one pass over the rows of
$A_{\mathbf d}$, accumulating
$\Delta_{jj}\langle a_j,x\rangle a_j^{\transpose}$.  It uses $O(M)$ field
elements of resident vector state and
$O(R_s(\mathbf d)(s+1)^2)$ field operations.  If
$r=\rank A_{\mathbf d}$, then
\begin{equation}\label{eq:preconditioner-success}
 \Pr_{\Delta}[\ker B\ne\ker A_{\mathbf d}]
 \le \frac{r}{|E|}\le\frac{M}{|E|}.
\end{equation}
Consequently scalar Wiedemann needs $O(M)$ field-vector state to seek a
one-dimensional right kernel, and block size $b$ needs $O(bM)$ state to seek a
$b$-dimensional kernel, under the standard generic-projection assumptions of
black-box Wiedemann.  The matrix is never materialized.

For the scalar $b=1$ realization used below, put
$\Omega=\F_{19^4}$ and $E=\F_{19^8}$.  If
$t=(s+1)^2$, one product by $B$ costs fewer than $5Rt$ signed
$\Omega$ multiply--accumulates: an $E$ element times an $\Omega$ scalar uses
two $\Omega$ multiplications, and an $E$ product uses three by Karatsuba.
Producing the scalar Wiedemann sequence and reconstructing one null vector
uses fewer than $3M$ products by $B$.  Thus
\begin{equation}\label{eq:explicit-stream-work}
 16R(s+1)^2M
\end{equation}
signed $\Omega$ multiply--accumulates is a rounded upper bound per solver
trial.  A straightforward seeded implementation ledger reserves eight
length-$M$ arrays over $E$ for the iterate, accumulator, scalar sequence,
reconstruction, and polynomial-work buffers.

For the robust Level-I profile only, take instead
$E'=\F_{19^{12}}$, a cubic extension of $\Omega$.  A sparse dot product and
row accumulation use at most $6t$ $\Omega$ multiplications, and naive cubic
multiplication uses nine more.  Since $6t+9<9t$ for $t=(s+1)^2\ge4$, the same
three-$M$ product ledger gives the rounded bound
\begin{equation}\label{eq:explicit-stream-work-cubic}
 32R(s+1)^2M
\end{equation}
signed $\Omega$ multiply--accumulates per trial.  Storing the three
$\Omega$ components in three bytes each uses nine bytes per $E'$ coordinate;
this is more conservative than bit-packing an $E'$ element.
\end{proposition}

\begin{proof}
The two sums in \eqref{eq:macaulay-row-count} count multiplier monomials for,
respectively, the $m_1$ diagonal equations of multidegree $2e_a$ and the
$m_1$ off-diagonal equations of multidegree $e_a+e_b$.  A homogenized
quadratic row has at most $(s+1)^2$ terms, which proves the streaming row and
transpose products.

For the probability bound, factor $A_{\mathbf d}=UV$, where
$U\in E^{R\times r}$ has full column rank and
$V\in E^{r\times M}$ has full row rank.  If $U^{\transpose}\Delta U$ is
invertible, then
\[
 \ker(A_{\mathbf d}^{\transpose}\Delta A_{\mathbf d})
 =\ker V=\ker A_{\mathbf d}.
\]
By Cauchy--Binet,
$\det(U^{\transpose}\Delta U)$ is a nonzero multilinear polynomial of total
degree $r$ in the diagonal entries: a nonzero $r\times r$ minor of $U$
contributes its squared determinant as a nonzero coefficient.  The
Schwartz--Zippel bound gives \eqref{eq:preconditioner-success}.  Standard
scalar and block Wiedemann recover null vectors from black-box products with
linear vector storage; block variants can recover a basis under their usual
random-projection conditions
\cite{Coppersmith1994BlockWiedemann,ChengEtAl2012XL,JouxPierrot2015NearlySparse}.

For the explicit scalar ledger, write an $E$ element in an $\Omega$-basis.
Each sparse dot product and each row accumulation uses at most $2t$
$\Omega$ products, while multiplying the dot product by the diagonal
preconditioner uses three.  Since $4t+3<5t$ for $t\ge4$, one normal-operator
product costs less than $5Rt$.  The usual scalar sequence requires fewer than
$2M$ products and null-vector reconstruction fewer than $M$ more; rounding
$15$ to $16$ gives \eqref{eq:explicit-stream-work}.  For a cubic extension,
the corresponding row cost is at most $6t+9<9t$; multiplying by the same
three-$M$ product count and rounding $27$ to $32$ gives
\eqref{eq:explicit-stream-work-cubic}.  The eight-array figures are
implementation reserves rather than asymptotic claims.
\end{proof}

\begin{theorem}[Recovery from a certified finite Macaulay quotient]
\label{thm:finite-quotient-recovery}
Fix a target and affine slice, and let
$I\subseteq\Omega[x_1,\ldots,x_h]$ be the ideal of its dehomogenized
residual equations.  Let $\mathcal B$ be a finite divisor-closed set of
monomials containing $1$, with $|\mathcal B|=D$, and let
$\partial\mathcal B$ be its border.  Suppose the selected Macaulay row space
provides, for every $\beta\in\partial\mathcal B$, a certified relation
\begin{equation}\label{eq:border-relation}
 g_\beta=\beta-\sum_{b\in\mathcal B}c_{\beta,b}b\in I.
\end{equation}
Define $D\times D$ formal multiplication matrices
$M_1,\ldots,M_h$ from these rewrite rules and from the tautological products
that remain in $\mathcal B$.  If
\begin{equation}\label{eq:border-commute}
 M_iM_j=M_jM_i\qquad(1\le i,j\le h),
\end{equation}
then $G=\{g_\beta:\beta\in\partial\mathcal B\}$ is a
$\mathcal B$-border basis of the zero-dimensional ideal $J=(G)$,
$\mathcal B$ is an $\Omega$-basis of $\Omega[x]/J$, and $J\subseteq I$.
Consequently $V(I)\subseteq V(J)$ and the attacker can enumerate every
base-field root of the original slice system from the multiplication tables.
A dense FGLM conversion costs $O(hD^3)$ operations in $\Omega$; finite-field
factorization and root extraction add randomized polynomial time in
$D$, $h$, and $\log|\Omega|$.  Every extra point is removed by the original
equations and public verifier.  The one-dimensional-kernel method is the
special case $D=1$.
\end{theorem}

\begin{proof}
The relations in \eqref{eq:border-relation} form a border prebasis for the
order ideal $\mathcal B$.  The commuting-multiplication-matrix
characterization of border bases states that such a prebasis is a border basis
exactly when its formal multiplication matrices commute
\cite{KehreinKreuzer2005}.  For completeness, the mechanism is as follows.
Commutativity makes
\[
 \pi:\Omega[x_1,\ldots,x_h]\longrightarrow
      \operatorname{End}_\Omega(\Omega^{\mathcal B}),
 \qquad x_i\longmapsto M_i,
\]
an algebra homomorphism.  Divisor closure gives
$\pi(b)e_1=e_b$ for every $b\in\mathcal B$.  By construction,
$\pi(g_\beta)e_1=0$ for every border relation.  Since the matrices commute,
for every $b\in\mathcal B$,
\[
 \pi(g_\beta)e_b
 =\pi(g_\beta)\pi(b)e_1
 =\pi(b)\pi(g_\beta)e_1=0.
\]
Thus each border relation acts as zero.  The standard border-division theorem
then gives unique normal forms in $\Span_\Omega\mathcal B$, so
$\Omega[x]/J$ has basis $\mathcal B$ and the $M_i$ are its multiplication
matrices.  The supplied Macaulay row combinations certify
$g_\beta\in I$, hence $J\subseteq I$ and therefore $V(I)\subseteq V(J)$.

FGLM converts the multiplication tables to a lexicographic Gr\"obner basis in
$O(hD^3)$ field operations~\cite{FaugereEtAl1993FGLM}.  Standard
zero-dimensional solving over finite fields then enumerates the geometric
points of $J$; inverse eigenblock and affine transformations reconstruct the
complete signature.  Checking membership in the base field, the original
slice equations, and the public verifier makes the procedure Las Vegas.
\end{proof}

\paragraph{How much quotient dimension is already priced.}
The per-trial row-reduction term in \eqref{eq:special-xl-cost} is
$3s^2M_s(\mathbf d)^2$.  FGLM has arithmetic complexity $O(hD^3)$, but the
hidden universal constant should not silently be set to one.  We therefore
report two break-even scales.  The constant-free scale is
\begin{equation}\label{eq:Dmax}
 D_{\rm raw}
 =\left\lfloor\left(\frac{3s^2M_s(\mathbf d)^2}{h}\right)^{1/3}\right\rfloor.
\end{equation}
For a deliberately extravagant implementation reserve, we also charge
$2^{30}hD^3$ extension-field operations and put
\begin{equation}\label{eq:Dmax-reserve}
 D_{30}
 =\left\lfloor\left(\frac{3s^2M_s(\mathbf d)^2}{2^{30}h}\right)^{1/3}\right\rfloor.
\end{equation}
Any extraction implementation costing at most $2^{30}hD^3$ is arithmetically
subsumed by the modeled row reduction whenever $D\le D_{30}$.  These bounds
answer only whether \emph{postprocessing already-supplied multiplication
tables} changes the leading arithmetic exponent.  They do not price obtaining
a $D$-dimensional nullspace, constructing all border relations, or the extra
block Wiedemann work when $D>1$.  They are also \emph{not} storage-feasibility
claims: dense multiplication tables require $\Theta(hD^2)$ field elements,
and every real run must impose an explicit memory cap (or certify a
sparse/streaming representation).  Thus $D>1$ is a correctness fallback, not
a free extension of the displayed scalar cost.  The arithmetic break-even
values are:

\begin{table}[H]
\centering
\caption{Arithmetic-only break-even dimensions for FGLM postprocessing once
multiplication tables are available.  $D_{\rm raw}$ ignores the fixed
implementation constant; $D_{30}$ reserves a factor $2^{30}$.  The table does
not price block/nullspace construction or the $\Theta(hD^2)$ storage of dense
multiplication tables.}
\label{tab:quotient-extraction}
\small
\setlength{\tabcolsep}{5pt}
\begin{tabular}{@{}c l r r r r@{}}
\toprule
$m_1$ & profile & $h$ & $D_{\rm raw}$ & $D_{30}$ & $\log_2D_{30}$\\
\midrule
5 & time minimum & 40 & $16{,}021{,}697$ & $15{,}646$ & 13.93\\
5 & model-state  & 32 & $110{,}695$ & $108$ & 6.75\\
5 & AXN-state    & 36 & $1{,}151{,}871$ & $1{,}124$ & 10.13\\
\midrule
7 & time minimum & 60 & $280{,}271{,}069{,}473$ & $273{,}702{,}216$ & 28.03\\
7 & state minimum & 40 & $370{,}128$ & $361$ & 8.50\\
\midrule
9 & time minimum & 76 & $92{,}610{,}866{,}101{,}092$ & $90{,}440{,}298{,}926$ & 36.40\\
9 & state minimum & 44 & $224{,}424$ & $219$ & 7.77\\
\bottomrule
\end{tabular}
\end{table}

\paragraph{Las Vegas wrapper.}
For each target--slice trial, row-reduce the Macaulay matrix.  If a
one-dimensional kernel appears, use the direct normalization method.  More
generally, continue when the row space supplies the border relations in
\Cref{thm:finite-quotient-recovery}, construct the quotient multiplication
matrices, and enumerate its roots.  Discard every point that fails base-field
descent or the original verifier.  Failure to obtain a complete finite quotient
at that degree causes a retry or a controlled increase of the multidegree; it
can never cause acceptance of an invalid signature.

\paragraph{Precisely what the displayed special-XL costs assume.}
The exponents in \Cref{tab:xl-optima} use (i) the aggregate condition
$\bar\eta_L\le19^{-1}$, which gives the factor
$\mathsf{Retry}^{+}_{\tau}$, and (ii) the published Hilbert-series prediction
that a successful selected-degree trial has a one-dimensional right kernel.
That is the case priced by scalar Wiedemann.  The finite-quotient theorem
proves zero-error correctness for $D>1$, and \Cref{tab:quotient-extraction}
shows that subsequent FGLM work can be small, but neither statement prices the
additional block/nullspace work needed to obtain all $D$ border relations.
Accordingly a production $D>1$ transcript must charge that work and its
success probability separately.  What remains unproved is the selected-degree
$D=1$ prediction (or a fully priced larger-$D$ replacement).

Before stating the frontier, note that a block degree $d_a=1$ does not
remove block $a$ or any of its affine variables.  Its monomial columns contain
the homogenizer and every block variable.  It only means that no multiple of a
diagonal equation of multidegree $2\mathbf e_a$ fits in that target component;
\eqref{eq:macaulay-row-count} assigns zero rows to that unavailable sector.
Off-diagonal equations still constrain the block, and the same signed Hilbert
coefficient and border-certificate tests apply.  Degree one is therefore the
natural minimum for a recovery matrix containing every variable.

\begin{table}[H]
\centering
\caption{Certified points on the special-XL time--state frontier for the six
$\ell=4$ rows.  ``One vector'' uses three bytes per $\F_{19^4}$ coordinate.
The model-state rows globally minimize $M$ among negative-coefficient profiles
with $d_i\ge1$ and published-model exponent at most $143/207/272$.  The
AXN-state rows are the corresponding minima under the secondary AXN bounds;
at Level~V one row satisfies both.}
\label{tab:xl-optima}
\tiny
\setlength{\tabcolsep}{2.6pt}
\begin{tabular}{@{}c l c c r r r r l@{}}
\toprule
$m_1$ & profile & $(s,\tau)$ & $\mathbf d$ & $\log_2M$ & $R/M$ & Model & AXN & one vector\\
\midrule
5 & time minimum & $(10,10)$ & $(3,3,3,4)$ & 34.45 & 2.69 & 128.90 & 132.25 & $65.4$ GiB\\
5 & model-state  & $(8,18)$  & $(1,2,3,3)$ & 23.39 & 1.99 & 140.13 & 143.49 & $31.5$ MiB\\
5 & AXN-state    & $(9,14)$  & $(2,2,3,4)$ & 28.83 & 2.36 & 134.34 & 137.70 & $1.33$ GiB\\
\midrule
7 & time minimum & $(15,10)$ & $(4,4,6,6)$ & 55.30 & 4.06 & 171.76 & 175.12 & $118$ PiB\\
7 & model-state  & $(9,34)$  & $(1,1,2,4)$ & 21.91 & 1.81 & 205.47 & 208.82 & $11.3$ MiB\\
7 & AXN-state    & $(10,30)$ & $(1,2,3,3)$ & 25.82 & 1.96 & 196.61 & 199.97 & $170$ MiB\\
\midrule
9 & time minimum & $(19,14)$ & $(5,6,6,6)$ & 67.68 & 4.58 & 214.20 & 217.56 & $614$ EiB\\
9 & model/AXN-state & $(11,46)$ & $(1,2,2,3)$ & 24.66 & 1.73 & 262.53 & 265.89 & $76.0$ MiB\\
\bottomrule
\end{tabular}
\end{table}

\begin{proposition}[Certified low-state frontier]
\label{prop:low-state-frontier}
Restrict only to multidegrees with $d_i\ge1$, so every variable block is
represented.  Over all slice dimensions available to both $\ell=4$ parameter
sets at each level and all multidegrees with negative coefficient in
\eqref{eq:hilbert-series}, the three model-state rows of
\Cref{tab:xl-optima} globally minimize $M_s(\mathbf d)$ subject to
published-model exponents at most $143$, $207$, and $272$.  The three
AXN-state rows globally minimize $M_s(\mathbf d)$ subject to AXN arithmetic
exponents at most the same references.
\end{proposition}

\begin{proof}
The exact finite search and monotone tail certificate are given in
\Cref{app:cost-certificates} and in the machine-readable frontier artifact.
\end{proof}

The original exact search checks every admissible $s$ and all 42,356 sorted
multidegrees whose modeled work could match or beat the time incumbent.  It
recovers the three time minima in \Cref{tab:xl-optima}, uniquely up to block
permutation.  A second exact search proves \Cref{prop:low-state-frontier}.
The model-state rows lower the simple sufficient cross-polar ranks from
$41/61/77$ to $33/37/45$; the Level-I and Level-III AXN-state rows require
ranks $37$ and $41$.  For the three model-state rows, random diagonal
preconditioning over $E=\F_{19^8}$ fails with probability at most $0.00065$,
$0.00024$, and $0.00157$ by \eqref{eq:preconditioner-success}.  The Level-I
and Level-III AXN-state bounds are $0.0281$ and $0.00350$ and can be reduced
geometrically by repetition.  The existing full-size rank campaign tested the
time-minimizing slices, not these smaller ones.  The reserved-line construction
therefore selects each low-state $A$-subspace from cross data and applies
\Cref{thm:seeded-spectrum}; a fixed key may additionally publish the exact
histogram.

Thus a below-reference point need not carry an astronomical state vector.
The table's three-byte figures describe vectors over $\F_{19^4}$ and isolate
the column-count frontier.  Over the rigorous preconditioning field
$E=\F_{19^8}$, the three AXN-state rows use packed one-vector states
$2.22$~GiB, $283.2$~MiB, and $126.8$~MiB.  The next result goes beyond an
asymptotic $O(M)$ statement: it prices the exact streamed row count and a
complete scalar-Wiedemann product ledger.  It still assumes that the selected
production degree has the published one-dimensional-kernel behavior.

\begin{table}[H]
\centering
\caption{Explicit streamed normal-operator accounting for the globally
minimal state-aware profiles below the numerical references.  ``Stream''
charges \eqref{eq:explicit-stream-work}, AXN arithmetic over $\F_{19^4}$,
$\bar\eta_L\le19^{-1}$, expected $\F_{19^8}$ preconditioner restarts, and a
factor $1+2^{-10}$ for extraction.  ``Spectrum'' is complete direct
projective enumeration; ``combined'' is the sum of the two work factors.}
\label{tab:explicit-stream}
\tiny
\setlength{\tabcolsep}{2.5pt}
\begin{tabular}{@{}c c c r r r r r@{}}
\toprule
Level & $(s,\tau)$ & $\mathbf d$ & $M$ & $R/M$ & Stream & Spectrum & Combined\\
\midrule
I   & $(9,14)$  & $(2,2,3,4)$ & $475{,}832{,}500$ & $2.358$ & $141.70$ & $173.20$ & $173.20$\\
III & $(10,30)$ & $(1,2,3,3)$ & $59{,}383{,}896$  & $1.956$ & $203.63$ & $191.68$ & $203.63$\\
V   & $(11,46)$ & $(1,2,2,3)$ & $26{,}574{,}912$  & $1.725$ & $269.35$ & $209.71$ & $269.35$\\
\bottomrule
\end{tabular}
\end{table}

\begin{proposition}[Explicit streamed state frontier]
\label{prop:explicit-stream-frontier}
Restrict to $d_i\ge1$, $M<19^8$, and all slice dimensions available to both
$\ell=4$ shapes at the indicated level.  Among every negative-coefficient
profile in \eqref{eq:hilbert-series} whose explicit streamed AXN work from
\eqref{eq:explicit-stream-work} is at most $2^{143}$, $2^{207}$, or
$2^{272}$, respectively, the three rows of \Cref{tab:explicit-stream}
globally minimize $M$.  Their streamed margins are $1.30$, $3.37$, and
$2.65$ bits.  One packed length-$M$ vector over $\F_{19^8}$ occupies
$2.22$~GiB, $283$~MiB, and $127$~MiB; the eight-array implementation reserve
occupies $17.7$~GiB, $2.21$~GiB, and $1.01$~GiB.
\end{proposition}

\begin{proof}
For each admissible $s$, both $M_s(\mathbf d)$ and
$R_s(\mathbf d)$ are coordinatewise nondecreasing in $\mathbf d$.  The
all-one profile therefore gives a monotone lower bound on streamed work.
Every profile with $M$ below the displayed incumbent lies in a finite sorted
degree box determined by
$\binom{s+D}{D}(s+1)^3<M_{\rm incumbent}$.  Exhaustive exact Hilbert
coefficient evaluation over those boxes finds no feasible smaller $M$; the
machine-readable certificate records every box and the monotone tail closure.
\end{proof}

This calculation is deliberately separate from the published
$3s^2M^2$ estimate: it prices the actual streamed row count, the extension
used for normal preconditioning, and the full scalar-Wiedemann product ledger.
At Levels~III and V, complete exact projective-spectrum computation is cheaper
than the compact solve and can therefore be run as public precomputation
without changing the leading exponent.  At Level~I direct enumeration costs
$2^{173.20}$, but the reserved-line theorem supplies the aggregate bound over
the random-XOF idealization.  A compressed determinantal certificate is needed
only for a deterministic fixed-seed claim below the numerical reference.

\paragraph{An implementation-level joint stress frontier.}
The compact profiles above optimize state under the headline values
$\bar\eta_L\le1/19$ and $\sigma=1$.  We now weaken both inputs at once and
price the actual streamed operator rather than the published $3s^2M^2$
estimator.  Every row below allows
\begin{equation}\label{eq:robust-inputs}
 \bar\eta_L\le\frac12,
 \qquad
 \sigma\ge\frac18,
\end{equation}
charges the exact row count, homogenizing coordinates, preconditioner restarts,
and the extraction reserve.  Levels~III and V use
\eqref{eq:explicit-stream-work}.  Level~I uses the cubic-preconditioner bound
\eqref{eq:explicit-stream-work-cubic}; only the random normal preconditioner
moves to $\F_{19^{12}}$.

\begin{table}[H]
\centering
\caption{Robust explicit streamed frontier.  All work columns are base-2 AXN
exponents.  ``State'' is the conservative eight-array envelope.  Exact
Hilbert enumeration and monotone tail closure prove minimum state among every
negative-coefficient profile with $d_i\ge1$ fitting the indicated explicit
bound and preconditioner field.}
\label{tab:robust-explicit}
\scriptsize
\setlength{\tabcolsep}{3.2pt}
\begin{tabular}{@{}c c c c r r r r@{}}
\toprule
Level & $(s,\tau)$ & $\mathbf d$ & preconditioner & $M$ & robust & margin & state\\
\midrule
I   & $(10,10)$ & $(3,3,3,4)$ & $\F_{19^{12}}$ & $23{,}417{,}049{,}656$ & $141.29$ & $1.71$ & $1.53$ TiB\\
III & $(11,26)$ & $(1,2,3,5)$ & $\F_{19^8}$ & $1{,}488{,}195{,}072$ & $200.49$ & $6.51$ & $55.4$ GiB\\
V   & $(12,42)$ & $(2,2,2,3)$ & $\F_{19^8}$ & $342{,}874{,}805$ & $264.03$ & $7.97$ & $12.8$ GiB\\
\bottomrule
\end{tabular}
\end{table}

\begin{proposition}[Robust explicit state frontier]
\label{prop:robust-explicit-frontier}
Restrict to $d_i\ge1$ and to all slice dimensions available to both structured
shapes at each level.  For the Level-I field $\F_{19^{12}}$ and the Level-III/V
field $\F_{19^8}$, the rows of \Cref{tab:robust-explicit} globally minimize
$M_s(\mathbf d)$ among every negative-Hilbert profile satisfying
\eqref{eq:robust-inputs} whose explicit streamed work is at most
$2^{143}$, $2^{207}$, or $2^{272}$.  Their one-vector states are
$196.3$~GiB, $6.93$~GiB, and $1.60$~GiB; the displayed state column reserves
eight such arrays.
\end{proposition}

\begin{proof}
Both $M_s(\mathbf d)$ and $R_s(\mathbf d)$ are coordinatewise nondecreasing.
For fixed largest degree $D$, the profile $(1,1,1,D)$ therefore lower-bounds
streamed work.  The all-one profile excludes the smaller slice dimensions, and
$M<M_{\rm incumbent}$ together with the relevant work bound leaves a finite
sorted degree box for every remaining $s$.  Exhaustive integer Hilbert
coefficient evaluation finds no feasible smaller $M$.  The artifact records
every box, tested profile count, and monotone tail closure.
\end{proof}

This joint stress test is stronger than the earlier estimator-only sensitivity
calculation.  It says that the numerical conclusion does not require a perfect
rank spectrum or a solver certificate on every root-containing trial.  The
seeded theorem is stronger still on the first input: for these robust slices,
$\Pr[\bar\eta_L>1/2]$ is below $2^{-41.48}$, $2^{-109.45}$, and
$2^{-177.41}$ in the random-XOF idealization.

There is also meaningful degree slack at Level~V.  The baseline robust run plus
all four one-coordinate degree-increment neighbors costs $2^{270.95}$ in
total, still below $2^{272}$.  This is a first-shell sensitivity statement,
not a theorem that one of those neighbors must provide a certificate.  A
higher-state Level-V profile $(s,\tau,\mathbf d)=(13,38,(1,2,3,5))$ costs only
$2^{257.02}$ under the same joint stress and uses a $262.8$~GiB eight-array
envelope, leaving $14.98$ bits of margin.

For comparison, once the seeded $\bar\eta_L\le1/19$ event holds, the compact
profiles remain below their references provided only that the production
certificate rates exceed $0.405$, $0.0968$, and $0.159$ at Levels~I, III, and
V.  These are break-even values, not measured production rates.  The 112
planted four-block Macaulay experiments reported in \Cref{sec:evidence} all
produce the predicted one-dimensional kernel, but the production transcript
remains the central missing artifact.

After charging the one-rank retry cushion, the time-minimizing structured AXN
rows retain $10.75$, $31.88$, and $54.44$ bits of nominal arithmetic gap at
Levels I, III, and V.  Because the modeled row-reduction work is quadratic in
the Macaulay column count, those budgets tolerate column-count growth factors
of about $41.7$, $6.2\times10^4$, and $1.6\times10^8$, respectively.  The
Level-III and Level-V time rows remain below their nominal references if every
block degree rises by one; the Level-I rows tolerate any two single-block
degree increments.  These are sensitivity statements inside the same
arithmetic model, not production timings.

\subsection{An exact three-plus-one decomposition}
\label{sec:three-plus-one}

The four-block system has ten unordered equation families: four diagonal
families and six cross families, each containing $m_1$ equations.  The final
block need not participate in the nonlinear Macaulay solve.  Write
$\mathcal E_{ab}=0$ for the $m_1$ dehomogenized equations supported on blocks
$a$ and $b$, where $1\le a\le b\le4$.  Affine and constant terms are included;
when one block is fixed, a cross family is affine-linear in the other block.

\begin{theorem}[Exact three-plus-one completion]
\label{thm:three-plus-one}
Let $\mathcal I_{123}$ be the ideal generated by the six families
$\mathcal E_{ab}$ with $1\le a\le b\le3$.  Suppose a certified first-stage
solver returns a finite set $\mathcal S\subseteq\Omega^{3s}$ containing every
point of $V(\mathcal I_{123})$.  For each
$p=(x_1,x_2,x_3)\in\mathcal S$, substitute $p$ into
$\mathcal E_{14},\mathcal E_{24},\mathcal E_{34}$ and write the resulting
$3m_1$ affine-linear equations as
\begin{equation}\label{eq:three-plus-one-linear}
 A_p x_4=b_p,
 \qquad A_p\in\Omega^{3m_1\times s}.
\end{equation}
Apply the following certified rule:
\begin{enumerate}[label=(\roman*),itemsep=2pt]
\item discard $p$ if \eqref{eq:three-plus-one-linear} is inconsistent;
\item if $\rank A_p=s$, solve for its unique $x_4$ and retain the completion
      exactly when $\mathcal E_{44}=0$;
\item if \eqref{eq:three-plus-one-linear} is consistent but
      $\rank A_p<s$, declare the trial incomplete unless a separately priced
      solver enumerates every completion.
\end{enumerate}
Whenever no incomplete case occurs, the retained set contains every root of
the full four-block system.  Base-field descent, affine reconstruction, and
the original verifier remove every extraneous point.  The resulting algorithm
is Las Vegas.
\end{theorem}

\begin{proof}
Let $(x_1,x_2,x_3,x_4)$ be a full root.  Its first three blocks lie in
$V(\mathcal I_{123})$ and therefore occur in $\mathcal S$.  At that partial
root, the true fourth block satisfies \eqref{eq:three-plus-one-linear}.  If the
trial is certified complete, the system cannot be a consistent rank-deficient
case; it is therefore full-column-rank and its unique solution is the true
$x_4$.  The diagonal check retains it.  Conversely, every retained completion
satisfies all ten equation families.  The final public-verifier check gives
zero-error output.
\end{proof}

The first-stage Hilbert series is the three-block specialization
\begin{equation}\label{eq:hilbert-series-three}
 H^{(3)}_{m_1,s}(t_1,t_2,t_3)=
 \frac{\prod_{a=1}^{3}(1-t_a^2)^{m_1}
       \prod_{1\le a<b\le3}(1-t_at_b)^{m_1}}
      {\prod_{a=1}^{3}(1-t_a)^{s+1}}.
\end{equation}
For $\mathbf d=(d_1,d_2,d_3)$, define
\begin{align}
 M^{(3)}_s(\mathbf d)
 &=\prod_{a=1}^{3}\binom{s+d_a}{d_a},
 \label{eq:three-plus-one-columns}\\
 R^{(3)}_s(\mathbf d)
 &=m_1\sum_{a=1}^{3}
   \binom{s+d_a-2}{d_a-2}\prod_{c\ne a}\binom{s+d_c}{d_c}
 \notag\\
 &\quad{}+m_1\sum_{1\le a<b\le3}
   \binom{s+d_a-1}{d_a-1}\binom{s+d_b-1}{d_b-1}
   \prod_{c\notin\{a,b\}}\binom{s+d_c}{d_c}.
 \label{eq:three-plus-one-rows}
\end{align}
The same streamed normal-operator proof as
\Cref{prop:streaming-macaulay} gives the per-trial bound
\begin{equation}\label{eq:three-plus-one-stream-work}
 16R^{(3)}_s(\mathbf d)(s+1)^2M^{(3)}_s(\mathbf d)
\end{equation}
signed $\F_{19^4}$ multiply--accumulates with eight packed
$\F_{19^8}$ arrays.  The retry exponent remains $\tau=K-4s$: the decomposition
changes how a root-containing trial is solved, not the probability that the
full slice contains an accepted singleton.  In the robust rows, $\sigma\ge1/8$
now denotes the joint conditional probability that the first-stage degree
supplies a complete finite quotient and that every potentially extending
partial root passes the certified linear-completion rule of
\Cref{thm:three-plus-one}.

The scalar ledger still uses the published one-dimensional-kernel prediction
for the first-stage planted system.  If a larger quotient is returned, every
additional nullspace, border-basis, enumeration, and completion operation---and
its resident state---must fit the displayed $1+2^{-10}$ reserve.  Otherwise the
trial is declared incomplete or the extra work is priced separately.  Thus
$\sigma$ is a certificate-within-budget probability, not merely the probability
that the first-stage ideal is finite.

\begin{table}[H]
\centering
\caption{Exact three-plus-one streamed profiles.  ``Compact'' uses
$\bar\eta_L\le1/19$ and $\sigma=1$; ``robust'' uses
$\bar\eta_L\le1/2$ and $\sigma\ge1/8$.  Both columns charge the exact
square-row acceptance density, expected normal-preconditioner restarts, and a
$1+2^{-10}$ extraction reserve.  State is an eight-array packed
$\F_{19^8}$ envelope.}
\label{tab:three-plus-one}
\scriptsize
\setlength{\tabcolsep}{2.8pt}
\begin{tabular}{@{}c l c c r r r r r@{}}
\toprule
Level & profile & $(s,\tau)$ & $\mathbf d$ & $M^{(3)}$ & $R^{(3)}/M^{(3)}$ & compact & robust & state\\
\midrule
I   & time boundary & $(9,14)$  & $(4,4,5)$ & $1{,}023{,}472{,}450$ & 2.891 & 144.25 & 148.17 & 38.1 GiB\\
III & time minimum  & $(11,26)$ & $(3,3,6)$ & $1{,}639{,}770{,}496$ & 2.614 & 197.02 & 200.94 & 61.1 GiB\\
V   & time minimum  & $(13,38)$ & $(3,4,4)$ & $3{,}172{,}064{,}000$ & 2.312 & 250.32 & 254.24 & 118.2 GiB\\
V   & robust low-state & $(12,42)$ & $(2,3,5)$ & $256{,}214{,}140$ & 2.183 & 259.48 & 263.40 & 9.55 GiB\\
\bottomrule
\end{tabular}
\end{table}

The state column uses five-byte packed representatives of $\F_{19^8}$.  A
simpler implementation storing one 64-bit word per coordinate multiplies those
figures by $8/5$; in particular, the Level-V low-state row occupies
$15.27$~GiB rather than $9.55$~GiB and remains a commodity-memory profile.

\begin{proposition}[Certified three-plus-one frontier]
\label{prop:three-plus-one-frontier}
Restrict to $d_i\ge1$, $M^{(3)}_s(\mathbf d)<19^8$, and every slice dimension
available to both structured shapes at a level.  The first three rows of
\Cref{tab:three-plus-one} globally minimize compact and robust streamed work
within this class, uniquely up to block permutation.  The Level-I minimum is
above $2^{143}$.  At Level~V, the last row globally minimizes
$M^{(3)}_s(\mathbf d)$ among robust rows below $2^{272}$.
\end{proposition}

\begin{proof}
For $d_i\ge1$, the three cross-family contributions imply
\[
 \frac{R^{(3)}_s(\mathbf d)}{M^{(3)}_s(\mathbf d)}
 \ge \frac{3m_1}{(s+1)^2}.
\]
Hence \eqref{eq:three-plus-one-stream-work} is at least
$48m_1(M^{(3)}_s(\mathbf d))^2$.  The incumbent work bound and
$M^{(3)}<19^8$ give a finite monomial-count cap for each $s$.  For sorted
degrees $d_1\le d_2\le d_3$,
\[
 M^{(3)}_s(\mathbf d)
 \ge (s+1)^2\binom{s+d_3}{d_3},
\]
so every remaining degree tail closes monotonically.  Exact integer expansion
of \eqref{eq:hilbert-series-three} over the resulting finite boxes gives the
displayed optima.  The machine-readable certificate records every box,
coefficient, row count, and tail closure.
\end{proof}

The decomposition does not improve Level~I, where the retry exponent is too
large.  At Level~III it supplies a substantially faster finite-state point,
although the four-block robust frontier has slightly lower state.  At
Level~V it is decisive: the time row leaves $17.76$ robust bits of margin, and
the low-state row improves the four-block robust point from $2^{264.03}$ and
$12.8$~GiB to $2^{263.40}$ and $9.55$~GiB.  The selected-degree caveat has not disappeared from this XL branch; it has
been narrowed to a three-block finite quotient plus a publicly checkable linear-completion
condition, all obtained within the stated work--state reserve.

\subsection{A selected-degree-free Frobenius-channel route}
\label{sec:channel-homotopy}

The preceding XL routes exploit the four eigenblocks after extending scalars.
The Frobenius normal form of \Cref{sec:frobenius-channels} gives a different
attack over $A=\F_{19^4}$ itself.  It is complementary rather than uniformly
better: it replaces the production solving-degree assumption by a symbolic
homotopy theorem, but exposes a soft-$O$ implementation factor and a much
larger straightforward data volume.

Write the transformed quotient equations on an $A$-linear slice as
\[
 F_0(x)=z_0\in A^{m_1},\qquad
 F_1(x)=z_1\in A^{m_1},\qquad
 F_2(x)=z_2\in\F_{19^2}^{m_1}.
\]
By \Cref{cor:frobenius-channel-shape}, $F_0$ has ordinary degree at most two,
$F_1$ has degree at most $20$, and $F_2$ has degree at most $362$.  Choose
$a$ whole coordinates of $F_0$ and $b$ whole coordinates of $F_1$, where
\begin{equation}\label{eq:channel-core-count}
 a+b=s,\qquad 0\le a,b\le m_1.
\end{equation}
These give a square system of $s$ equations in $s$ variables over $A$.  Solve
that core completely, then check every unused coordinate of all three channels,
base-field descent, affine reconstruction, the rejection rule, and the original
verifier.  No equation is discarded: unused equations are filters.

\paragraph{A public separability certificate.}
Let $J_{a,b}(x)$ be the Jacobian of the square core, and let
$D_{a,b}^{\rm top}(x)$ be the highest-degree homogeneous part of
$\det J_{a,b}(x)$.  A degree-two row contributes degree one to the Jacobian;
a degree-twenty row has highest part $x_jx_k^{19}$, whose formal derivative in
characteristic $19$ has degree $19$.  Hence
\begin{equation}\label{eq:channel-jacobian-degree}
 \deg\det J_{a,b}\le a+19b.
\end{equation}
Translation to the target-dependent affine coset changes only lower-degree
Jacobian terms.  Therefore one exact public witness
$x_0\in A^s$ with $D_{a,b}^{\rm top}(x_0)\ne0$ certifies, once per chosen
slice and core, that the Jacobian determinant is not the zero polynomial for
every target trial.  Such a witness is a cheap preflight, not a solver
heuristic.

\begin{theorem}[Accepted nonsingular channel roots]
\label{thm:accepted-nonsingular-singleton}
Assume the setting of \Cref{thm:accepted-singleton}, with an $A$-linear slice
of $A$-dimension $s$, and suppose the square channel core satisfies the public
leading-Jacobian certificate above.  Let $\alpha$ be the density of signatures
accepted by the public rejection rule, put
\[
 \delta_{a,b}=\frac{a+19b}{|A|}=\frac{a+19b}{19^4},
 \qquad \mu=19^{4s-K},
\]
and suppose $\bar\eta_L\le\eta$.  Then
\begin{equation}\label{eq:accepted-nonsingular-bound}
 \Pr[\text{the trial has an accepted full root nonsingular for the core}]
 \ge \mu\frac{(\alpha-\delta_{a,b})^2}{1+\eta}.
\end{equation}
In particular, if $\alpha>\delta_{a,b}$, the expected number of target--slice
trials to such a root is at most
\begin{equation}\label{eq:channel-retry}
 19^{K-4s}\frac{1+\eta}{(\alpha-\delta_{a,b})^2}.
\end{equation}
The root need not be unique, and no independence between singularity,
acceptance, and the unused filter equations is assumed.
\end{theorem}

\begin{proof}
The nonzero Jacobian determinant has at most
$(a+19b)|A|^{s-1}$ zeros in $A^s$ by the finite-field Schwartz--Zippel bound.
For a uniform full quotient target, each fixed slice point satisfies all $K$
scalar residual equations with probability $19^{-K}$.  If $Y$ counts accepted
full roots that are nonsingular for the core, then
\[
 \mathbb EY\ge\mu(\alpha-\delta_{a,b}).
\]
Furthermore $Y(Y-1)\le Z(Z-1)$, so
\Cref{thm:coset-success} gives
$\mathbb E[Y(Y-1)]\le\mu\eta$ and hence
$\mathbb E[Y^2]\le\mu(1+\eta)$.  Cauchy--Schwarz yields
$\Pr[Y>0]\ge(\mathbb EY)^2/\mathbb E[Y^2]$, which is
\eqref{eq:accepted-nonsingular-bound}.
\end{proof}

The core can now be solved by a theorem-backed algorithm.  We state the precise
finite-field form needed here because the published proposition phrases its
field hypothesis as a characteristic bound.

\begin{lemma}[Finite-cardinality multihomogeneous symbolic homotopy]
\label{lem:finite-cardinality-homotopy}
Let $k$ be a perfect finite field.  Partition $N=\sum_{j=1}^m n_j$
variables into blocks $X_1,\ldots,X_m$ of sizes $n_1,\ldots,n_m$, and let
$f_1,\ldots,f_N$ be nonconstant polynomials represented by a straight-line
program of length $L$.  Suppose $f_i$ has multidegree at most
$\mathbf d_i=(d_{i,1},\ldots,d_{i,m})$.  Define
\begin{align}
 C_{\mathbf n}(\mathbf d)
 &= [\vartheta_1^{n_1}\cdots\vartheta_m^{n_m}]
    \prod_{i=1}^N\left(\sum_{j=1}^m d_{i,j}\vartheta_j\right),
 \label{eq:multihom-C}\\
 C_{\mathbf n'}(\mathbf d')
 &=\operatorname{coeffsum}\!\left(
   \prod_{i=1}^N\left(\vartheta_0+
        \sum_{j=1}^m d_{i,j}\vartheta_j\right)
   \bmod\langle\vartheta_0^2,
       \vartheta_1^{n_1+1},\ldots,\vartheta_m^{n_m+1}\rangle\right),
 \label{eq:multihom-Cplus}
\end{align}
and put $B=C_{\mathbf n}(\mathbf d)$ and
\begin{equation}\label{eq:finite-cardinality-threshold}
 E=\max\left\{
   \max_{1\le j\le m}\sum_{i=1}^N d_{i,j},
   8B^2
 \right\}.
\end{equation}
If $|k|\ge E$, a randomized symbolic-homotopy algorithm returns a
zero-dimensional parametrization containing every nonsingular root with
probability at least $7/8$; on every run it returns a subset of the
nonsingular roots or reports failure.  Its arithmetic cost is
\begin{equation}\label{eq:homotopy-soft-cost}
 \widetilde O\!\left(
 B C_{\mathbf n'}(\mathbf d')
 \left(L+\sum_{i,j}d_{i,j}+N^2\right)N
 \right)
\end{equation}
operations in $k$.
\end{lemma}

\begin{proof}
Safey El Din and Schost prove the same algorithm and cost when the
characteristic is zero or at least
$\max\{\max_j\sum_i d_{i,j},8(N-1)B^2\}$
\cite[Proposition~5]{SafeyElDinSchost2018}.  We replace the two uses of that
large-characteristic hypothesis by field-cardinality arguments.

For block $j$, their start system uses
$D_j=\sum_i d_{i,j}$ affine moment-curve forms.  Choose arbitrary distinct
$\alpha_{j,0},\ldots,\alpha_{j,D_j-1}\in k$ and set
\[
 \kappa_{j,u}(X_j)=X_{j,1}+\alpha_{j,u}X_{j,2}+
 \cdots+\alpha_{j,u}^{n_j-1}X_{j,n_j}+\alpha_{j,u}^{n_j}.
\]
Any $n_j$ of these forms have an invertible Vandermonde coefficient matrix,
whereas any $n_j+1$ give an inconsistent affine system because the augmented
Vandermonde matrix is invertible.  The proof of their start-system lemma
therefore applies verbatim: the product system has exactly $B$ roots, all
simple, and they are enumerated within the same soft-linear cost.  The first
term in \eqref{eq:finite-cardinality-threshold} supplies the distinct nodes.

It remains to choose the separating linear form.  The cited proof identifies
at most $B^2$ nonzero vectors $v\in\mathcal A^N$, where $\mathcal A$ is a
domain containing $k$, on none of which the chosen linear form may vanish.
Choose its coefficient vector $\lambda$ uniformly from $k^N$.  For fixed
$v\ne0$, the $k$-linear map
$\lambda\mapsto\lambda\mathbin{\cdot}v$ is nonzero, so its kernel contains at
most $|k|^{N-1}$ coefficient vectors.  Hence
\[
 \Pr_{\lambda\leftarrow k^N}[\lambda\mathbin{\cdot}v=0]\le |k|^{-1},
 \qquad
 \Pr[\lambda\text{ is not well separating}]\le B^2/|k|\le1/8.
\]
This replaces the one-parameter family in their separating-form lemma and
removes the factor $N-1$ from the field-size threshold.  Their auxiliary
algorithm accepts an arbitrary well-separating linear form, so no later step
changes.

The homotopy curve, Newton lifting, specialization, squarefree cleaning, and
zero-dimensional parametrization use only field operations and inversion at
nonsingular points.  Since finite fields are perfect, the remainder of the
correctness and arithmetic-cost proof is unchanged.  This is a
finite-cardinality adaptation proved here; the cited proposition itself is
stated under its original characteristic hypothesis.
\end{proof}


For the channel core, $N=s$, the degrees are $a$ copies of two and $b$ copies
of twenty, and therefore
\begin{equation}\label{eq:channel-bezout}
 B=2^a20^b,
 \qquad
 B^+=B\left(1+\frac a2+\frac b{20}\right).
\end{equation}
A conservative straight-line program first computes the
$s(s+1)/2$ unordered products for the degree-two rows.  For the twisted rows it
computes every $x_j^{19}$ with at most six multiplications and every ordered
product $x_jx_k^{19}$.  Charging one coefficient multiplication and one
addition per possible term gives
\begin{equation}\label{eq:channel-slp}
 L_{a,b,s}=
 \mathbf1_{a>0}\!\left[\binom{s+1}{2}+2a\left(\binom{s+1}{2}+s+1\right)\right]
 +\mathbf1_{b>0}\!\left[6s+s^2+2b(s^2+2s+1)\right].
\end{equation}
Choose the smallest $r$ for which $(19^4)^r\ge E$ and run the homotopy over an
$r$-degree extension $\widetilde A/A$.  Multiplication in that extension can
be reduced to $c_r=2r-1$ pointwise multiplications in $A$ by evaluation and
interpolation at $2r-1$ distinct $A$-points.  This counts only the leading
pointwise products.  We therefore expose, rather than hide, a multiplier
$\kappa_{\rm hom}$ containing every soft-$O$ factor, inversion, evaluation and
interpolation operation, filter and factorization cost, memory traffic, and
implementation constant.

Define the unit-normalized leading AXN work
\begin{equation}\label{eq:channel-leading-work}
 W^{\rm lead}_{a,b,s}(\eta)=
 \frac{8}{7}\,
 19^{K-4s}\frac{1+\eta}{(\alpha-\delta_{a,b})^2}\,
 c_r g_4^{\rm AXN}\,
 B B^+\bigl(L_{a,b,s}+2a+20b+s^2\bigr)s.
\end{equation}
The factor $8/7$ repeats the homotopy until its good random choice occurs.
The actual implementation ledger is
$\kappa_{\rm hom}W^{\rm lead}_{a,b,s}$; \eqref{eq:channel-leading-work} is not
asserted to upper-bound the suppressed factor.

\begin{table}[H]
\centering
\caption{Balanced Frobenius-channel homotopy profiles for all six $\ell=4$
Version~2.3 rows.  Compact uses $\eta=1/19$; robust uses $\eta=1/2$.
The square-row lower bound $\alpha=(1-19^{-4})^n$ is charged for both shapes at
each level.  ``$+7$'' is the sensitivity value obtained by setting
$\kappa_{\rm hom}=2^7$, not a claim that $2^7$ bounds the suppressed factor.
The last column is a straightforward dense coefficient-volume ceiling, not a
proved resident-state lower bound.}
\label{tab:channel-homotopy}
\scriptsize
\setlength{\tabcolsep}{2.8pt}
\begin{tabular}{@{}c c r r c r r r r@{}}
\toprule
Level & $(s;a,b;\tau)$ & $B$ & $B^+$ & $r$ & compact & robust & robust $+7$ & dense ceiling\\
\midrule
I   & $(9;5,4;14)$  & $5{,}120{,}000$ & $18{,}944{,}000$ & 3 & 135.00 & 135.51 & 142.51 & 6.03 PiB\\
III & $(10;7,3;30)$ & $1{,}024{,}000$ & $4{,}761{,}600$  & 3 & 199.06 & 199.57 & 206.57 & 341.5 TiB\\
V   & $(11;9,2;46)$ & $204{,}800$     & $1{,}146{,}880$  & 3 & 263.03 & 263.55 & 270.55 & 17.94 TiB\\
\bottomrule
\end{tabular}
\end{table}

\begin{proposition}[Certified channel-core ledger]
\label{prop:channel-homotopy-ledger}
For the three rows of \Cref{tab:channel-homotopy}, the Jacobian singular-density
bounds are respectively
\[
 \frac{81}{19^4},\qquad \frac{64}{19^4},\qquad \frac{47}{19^4}.
\]
The compact unit-normalized leading exponents are
$134.99722$, $199.06332$, and $263.03426$; the robust exponents are
$135.50818$, $199.57429$, and $263.54523$.  Thus the robust break-even values
for $\log_2\kappa_{\rm hom}$ are $7.49182$, $7.42571$, and $8.45477$ bits.
Among every whole-channel square core satisfying \eqref{eq:channel-core-count},
the displayed rows minimize the straightforward dense coefficient-volume
ceiling subject to retaining seven robust sensitivity bits below the numerical
reference.  All statements are reproduced by exact integer enumeration in the
companion certificate.
\end{proposition}

\begin{proof}
Substitute \eqref{eq:channel-bezout} and \eqref{eq:channel-slp} into
\eqref{eq:channel-leading-work}, use the exact accepted densities from
\Cref{sec:rejection}, and enumerate
$1\le s\le2m_1$ and
$\max(0,s-m_1)\le a\le\min(m_1,s)$ with $b=s-a$.  For fixed $s$, replacing a
degree-twenty equation by a degree-two equation decreases $B$, $B^+$, the
degree sum, and the SLP bound; the exhaustive certificate nevertheless checks
every mix.  It also records the extension-cardinality threshold, packed output
size, dense coefficient-volume ceiling, and every break-even factor.
\end{proof}

The final zero-dimensional parametrizations in the three balanced rows contain
at most $(s+1)B$ extension-field elements, namely about $342$~MiB,
$75.2$~MiB, and $16.4$~MiB in packed form.  A straightforward dense
parametric-homotopy representation can be much larger, as the final column
shows; \Cref{lem:finite-cardinality-homotopy} is a time theorem, not a
low-memory theorem.  For comparison, minimizing that dense ceiling subject
only to a compact normalized exponent below the reference gives
$2^{142.68}$ with $13.7$~TiB at Level~I, $2^{206.79}$ with $786$~GiB at
Level~III, and $2^{270.06}$ with $29.8$~GiB at Level~V.  Those boundary rows
leave essentially no room for suppressed factors.

This branch changes the logical status of the paper.  A reviewer may reject
the selected-degree extrapolation used by streamed XL, or may instead regard
the symbolic-homotopy constants and data movement as too large for a concrete
claim.  Those objections are no longer the same objection: the homotopy branch
has theorem-backed root completeness and no selected degree, while the XL
branch has the explicit low-state operator frontier.  The unconditional
quotient, exact filtering reduction, retry theorem, and final-verifier check
are common to both.


\subsection{Low-output cores from the original quadratic families}
\label{sec:eigenblock-cores}

The Frobenius-channel construction raises some equations to degree $20$ in
order to work directly over $A$.  A faster selected-degree-free route keeps the
original degree-two eigenblock equations.  It solves only two or three
of the four scalar-extension blocks nonlinearly, uses Frobenius descent to
reconstruct the omitted blocks, and filters every unused equation.

Fix $b\in\{2,3\}$ consecutive eigenblocks, each with $s$ variables.  Select
$l_a$ equations from diagonal family $(a,a)$ and $e_{ab}$ equations from
cross family $(a,b)$, with
\begin{equation}\label{eq:eigenblock-square-count}
 0\le l_a,e_{ab}\le m_1,
 \qquad
 \sum_a l_a+\sum_{a<b}e_{ab}=bs.
\end{equation}
The resulting square core has exact multihomogeneous B\'ezout number
\begin{equation}\label{eq:eigenblock-bezout}
 B=
 [t_1^s\cdots t_b^s]
 \prod_{a=1}^{b}(2t_a)^{l_a}
 \prod_{1\le a<c\le b}(t_a+t_c)^{e_{ac}}.
\end{equation}
Its companion quantity $B^+$ is
\eqref{eq:multihom-Cplus}.  Both are exact integers, not semi-regularity
predictions.

A base-field point becomes, after scalar extension, a Frobenius orbit
$(x,x^{19},x^{19^2},x^{19^3})$.  Thus every genuine residual root restricts
to a root of the selected $b$-block core.  Conversely, after the homotopy
enumerates the core roots, the attacker checks the Frobenius relations among
the selected blocks, reconstructs every omitted block from one selected block,
and checks all unused residual equations, rejection, affine reconstruction,
and the original verifier.  No linear-completion or selected-Macaulay-degree
assumption is involved.

\paragraph{A public leading-Jacobian certificate.}
Let $\gamma_a$ be the degree of the core-Jacobian determinant in block $a$;
explicitly,
\begin{equation}\label{eq:eigenblock-jacobian-multidegree}
 \gamma_a=
 2l_a+\sum_{c\ne a}e_{\min\{a,c\},\max\{a,c\}}-s.
\end{equation}
On the descent locus, its leading part pulls back to a polynomial over $A$ of
degree at most
\begin{equation}\label{eq:eigenblock-descent-degree}
 \Delta=\sum_{a=0}^{b-1}\gamma_a19^a.
\end{equation}
The attacker evaluates this leading determinant at one fixed public descent
point.  A nonzero value certifies that the determinant polynomial is nonzero
for every target translate.  The companion artifact supplies exact allowable
symmetric-form assignments at which this determinant is nonzero for each
recommended profile; hence the preflight is a proper algebraic condition, not
an assumed possibility.

\begin{theorem}[Accepted nonsingular eigenblock-core roots]
\label{thm:eigenblock-core-roots}
Let a selected core satisfy the public leading-Jacobian certificate, let
$\alpha$ be the density accepted by the public rejection rule, and suppose
$\bar\eta_L\le\eta$.  Put $h=4s$, $\mu=19^{h-K}$, and
$\delta=\Delta/19^4$.  Then a random target--coset trial contains an accepted
full residual root that is nonsingular for the selected core with probability
at least
\begin{equation}\label{eq:eigenblock-core-success}
 \max\left\{
   \mu(\alpha-\eta-\delta),
   \mu\frac{(\alpha-\delta)^2}{1+\eta}
 \right\},
\end{equation}
where a nonpositive first term is ignored.  Finite-cardinality symbolic
homotopy followed by the exact filters above is therefore a Las Vegas forger.
It uses neither a production solving-degree prediction nor a one-dimensional
Macaulay-kernel assumption.
\end{theorem}

\begin{proof}
The fixed leading-Jacobian certificate and
\eqref{eq:eigenblock-descent-degree} bound the density of descent points
singular for the core by $\delta$.  Let $Y$ count accepted full roots that are
nonsingular for the core and let $Z$ count all full roots.  The first moment is
at least $\mu(\alpha-\delta)$.  Since $Y(Y-1)\le Z(Z-1)$, the aggregate
collision theorem gives $\mathbb E[Y(Y-1)]\le\mu\eta$.  The singleton
inequality $\Pr[Y=1]\ge\mathbb EY-\mathbb E[Y(Y-1)]$ gives the first term of
\eqref{eq:eigenblock-core-success}; Cauchy--Schwarz with
$\mathbb E[Y^2]\le\mu(1+\eta)$ gives the second.  Conditional on such a root,
\Cref{lem:finite-cardinality-homotopy} returns it with probability at least
$7/8$.  Every later operation is an exact public filter.
\end{proof}

The random-XOF model also makes the leading-Jacobian preflight quantitative.
Group the fresh $V'\times V'$ coefficients by public base form.  If the
nonzero determinant polynomial has degree at most $d_i$ in form block $i$,
recursive application of \Cref{lem:biased-affine-subspace} gives
\begin{equation}\label{eq:block-anticoncentration}
 \Pr[\det J_{\rm core}\ne0]\ge
 \beta_{\rm core}=\prod_{i=1}^{m_1}(1-d_i\rho).
\end{equation}
The exact witness fixes the nonzero coefficient needed by that recursion.  The
core event and the seeded-spectrum event use the same internal coefficient
region, so they are combined by the union bound
$\beta_{\rm core}-\varepsilon_{\rm spec}$, not by an independence claim.
They are independent of the complete structural preflight in
\Cref{thm:structured-preflight-density}.  Consequently the certified
random-key density is
\begin{equation}\label{eq:eigenblock-vulnerable-density}
 p_{\rm eig}
 \ge p_{\rm pre}(\beta_{\rm core}-\varepsilon_{\rm spec}).
\end{equation}

\begin{table}[H]
\centering
\caption{Recommended original-family eigenblock cores.  Compact uses
$\eta=1/19$ and robust uses $\eta=1/2$.  ``Normalized'' pays the inverse exact
lower bound \eqref{eq:eigenblock-vulnerable-density}, so it is a
work-per-success ledger over a random generated key in the random-XOF
idealization.  ``$+8$'' multiplies the leading homotopy term by $2^8$ as a
sensitivity point.  Output is a packed final-parametrization ceiling, not a
peak-memory theorem.}
\label{tab:eigenblock-core}
\scriptsize
\setlength{\tabcolsep}{3.0pt}
\begin{tabular}{@{}c c c r c r r r r r@{}}
\toprule
Level & core & $(l_a;e_{ab})$ & $B$ & $r$ & robust/key & $p_{\rm eig}$ & normalized & normalized margin & output\\
\midrule
I & $3\times10$ & $(5,5,5;5,5,5)$ & $73{,}793{,}536$ & 4 & 134.077 & 0.129015 & 137.031 & 5.969 & 19.793 GiB\\
III & $3\times11$ & $(7,7,0;7,7,5)$ & $6{,}881{,}280$ & 3 & 196.246 & 0.116451 & 199.348 & 7.652 & 1.570 GiB\\
V & $2\times13$ & $(9,8;9)$ & $16{,}515{,}072$ & 3 & 247.432 & 0.200117 & 249.753 & 22.247 & 3.015 GiB\\
\bottomrule
\end{tabular}
\end{table}

For the three rows, the exact per-form degree vectors in
\eqref{eq:block-anticoncentration} are
\[
 (6,6,6,6,6),\qquad
 (5,5,5,5,5,4,4),\qquad
 (3,3,3,3,3,3,3,3,2),
\]
giving $\beta_{\rm core}=0.1369122816$, $0.1235790793$, and
$0.2123661755$.  The robust spectrum-failure probabilities are below
$2^{-41.48}$, $2^{-109.45}$, and $2^{-160.42}$, respectively.  Before paying
key density, the robust break-even budgets for
$\log_2\kappa_{\rm hom}$ are $8.923$, $10.754$, and $24.568$ bits; setting
$\kappa_{\rm hom}=2^8$ gives $142.077$, $204.246$, and $255.432$ per
certified key.  The normalized $+8$ values are $145.031$, $207.348$, and
$257.753$, so the first two levels still need an implementation measurement to
turn the normalized sensitivity row into a below-reference claim.

The released search is exhaustive within the stated original-family class:
it checks every admissible one-, two-, and three-block square core and records
the complete time/output Pareto frontier.  It examines $15{,}663$ count
vectors at Level I, $87{,}643$ at Level III, and $333{,}841$ at Level V, of
which $11{,}119$, $61{,}442$, and $232{,}706$ have nonzero B\'ezout number.
Thus the displayed profiles are not hand-selected examples.  Smaller-output
points include $872.3$~MiB at Level I, $42.1$~MiB at Level III, and
$477.8$~MiB or $47.3$~MiB at Level V, with the exact arithmetic tradeoffs in
the artifact.


\subsection{A Frobenius-orbit-complete homotopy frontier}
\label{sec:basis-complete-homotopy}

We now remove the preferred-core premise entirely.  Let the four eigenblocks be
indexed cyclically by $\mathbb Z/4\mathbb Z$.  For a stable $A$-linear slice of
$A$-dimension $s$, write $h=4s$.  If the full restricted residual Jacobian has
rank $h$ at a descended root, some $h$-row minor is nonzero.  When $h=K-2$,
each family-count pattern is obtained by deleting two rows from the ten full
multidegree families, giving
\[
  \binom{10+1}{2}=55
\]
patterns.

\begin{lemma}[Frobenius transport of minors]
\label{lem:frobenius-minor-transport}
At a point satisfying Frobenius descent, $19$-Frobenius cyclically rotates the
four eigenblocks.  It sends the determinant of every square restricted
Jacobian minor to the determinant of the cyclically rotated minor, raised to
the nineteenth power.  In particular, one minor is nonzero if and only if
every minor in its cyclic orbit is nonzero.
\end{lemma}

\begin{proof}
The residual equations and the eigenblock coordinates are defined over
$\F_{19}$ before scalar extension.  Frobenius therefore commutes with
restriction and differentiation, while cyclically permuting the eigenblocks.
Applying Frobenius to a determinant raises its value to the nineteenth power,
which preserves zero and nonzero values.
\end{proof}

Burnside enumeration gives exactly $16$ orbits of the $55$ patterns under this
cyclic action.  The attacker solves one representative from each orbit by the
finite-cardinality symbolic-homotopy algorithm, and filters every candidate
through all omitted residual equations, all descent equations, rejection,
affine reconstruction, and the original verifier.  By
\Cref{lem:frobenius-minor-transport}, every descended root at which the full
restricted Jacobian has rank $h$ is nonsingular for at least one solved
representative.

\begin{proposition}[Orbit-complete fixed-core-free frontier]
\label{prop:selected-homotopy-frontier}
For Levels I, III, and V, take respectively $(h,K)=(48,50),(68,70),(88,90)$.
Under the robust aggregate-spectrum and exact acceptance bounds, the sum over
the $16$ orbit representatives has per-certified-key AXN exponents
\[
  138.32403,\qquad 179.75835,\qquad 220.83240.
\]
The complete structural preflight of \Cref{thm:structured-preflight-density}
succeeds in the random-XOF idealization with probability greater than
$0.9423222679$ for every official $\ell=4$ shape.  Charging its inverse gives
the random-key-normalized exponents
\[
  138.40973,\qquad 179.84405,\qquad 220.91810.
\]
The corresponding margins below $143/207/272$ are $4.59027$, $27.15595$, and
$51.08190$ bits.
\end{proposition}

\begin{proof}
The artifact exhaustively enumerates the $55$ patterns, quotients them by the
cyclic Frobenius action, and obtains $16$ orbits.  For every representative it
computes the exact multihomogeneous B\'ezout number, the companion homotopy
quantity, extension degree, straight-line-program bound, candidate-filter
cost, and robust retry factor.  Summing these ledgers gives the per-key
exponents.  The final values add
$-\log_2(0.9423222679)<0.08571$ bits and the negligible seeded-spectrum tail.
\end{proof}

This theorem-backed branch has no selected Macaulay degree, no exact-corank
assumption, and no preferred core-Jacobian condition.  Its straightforward
largest individual parametrizations remain extremely large: approximately
$1.890$~PiB, $2.076$~ZiB, and $2416.552$~YiB.  These are dense sequential
output ceilings rather than peak-memory lower bounds, but they make clear that
the purpose of this branch is logical completeness.  The original-equation
eigenblock cores of \Cref{sec:eigenblock-cores} provide much smaller
output ceilings, while the streamed-XL branch provides the low-state frontier.

\subsection{The \texorpdfstring{$\ell=2$}{l=2} attacks}

For Levels I and III, use the fixed-skew lift of \Cref{prop:fixed-skew}.  The
public preflight gives quotient rank $K$, affine rank $M-K$, and residual
systems of dimensions $48/112$ and $72/168$.  CFXL is a hybrid strategy
that specializes some variables and then applies XL-style linear algebra;
Hashimoto refines this optimization for underdetermined MQ.\@  Minimizing these
variants in Beullens's printed Appendix-A convention gives per-solve
exponents $129.44$ and $183.44$ modeled bit operations
\cite{Beullens2024SNOVA,Hashimoto2023}.  Replacing $b_1$ by
$g_1^{\mathrm{AXN}}$ gives per-solve arithmetic exponents $133.75$ and
$187.74$.  Every recovered root is accepted.  We charge
$\Theta_{\rm pol}\le19^{-1}$ from \Cref{cor:aggregate-full-space}; the
minimum ranks $r_*\ge49$ and $r_*\ge73$ are sufficient.  The resulting
headline exponents are $129.51/133.82$ and $183.51/187.82$ in the Model/AXN
columns.  The same fixed-skew construction at Level~V has per-solve exponents
$237.84/242.15$ and headline exponents $237.92/242.22$ under the same aggregate condition; $r_*\ge97$ is sufficient.
The stronger thresholds $96$, $144$, and $192$ remove target retries entirely.

The optimizer includes Hashimoto's published boundary value $a=1$.  For
$\operatorname{MQ}(19,1,1)$, the semi-regular Hilbert-series rule used by CFXL
never produces the stopping degree at which its linearization would be costed.
The artifact therefore solves that elementary system by 19 direct evaluations
and deliberately overcharges each evaluation.  This
lowers six generic quotient baselines and the Level-V fixed-skew fallback; it
does not change the selected Level-I or Level-III fixed-skew costs.

The Level-V headline instead uses $R_0=I,R_1=0$ and zero offsets.  The quotient
has rank $K$, target filtering enforces the $M-K$ public consistency
conditions, and the residual system has 96 equations in all 256 variables.
Target search occurs before equation solving, so its cost adds to---rather than
multiplying---the cost of a solve.  The generic-MQ per-solve AXN arithmetic
exponent is $239.50$, while exact-uniform consistency filtering is smaller by
more than 70 exponent bits.  Under the headline aggregate condition $\Psi_{\rm pol}\le19^{-1}$,
\eqref{eq:filter-spectrum-cushion} adds less than $0.075$ bits, giving Model/AXN
exponents $235.27/239.57$.  The stronger condition $r_*\ge240$ gives an
accepted root for every target.  The fixed-skew fallback is slower but makes
every recovered root acceptable by construction.

\begin{table}[tbp]
\centering
\caption{$\ell=2$ systems, aggregate headline conditions, and sufficient
minimum-rank certificates.  Model and AXN charge
$\Theta_{\rm pol}\le19^{-1}$ for fixed skew and
$\Psi_{\rm pol}\le19^{-1}$ for the filter.  ``Budget'' is the smallest
minimum rank preserving positive nominal AXN nominal arithmetic gap; ``suff.'' is a
minimum rank implying the aggregate headline condition; ``all target''
guarantees an accepted root for every target.}
\label{tab:l2-costs}
\footnotesize
\setlength{\tabcolsep}{2.15pt}
\begin{tabular}{@{}l l r r r r r r@{}}
\toprule
Parameters & mode & eqs./vars. & model & AXN & budget & suff. & \shortstack{all\\target}\\
\midrule
$(48,16,2,2)$ & fixed-skew & $48/112$ & 129.51 & 133.82 & 46  & 49  & 96\\
$(72,24,2,2)$ & fixed-skew & $72/168$ & 183.51 & 187.82 & 68  & 73  & 144\\
$(96,32,2,2)$ & filter     & $96/256$ & 235.27 & 239.57 & 177 & 194 & 240\\
$(96,32,2,2)$ & fixed-skew & $96/224$ & 237.92 & 242.22 & 89  & 97  & 192\\
\bottomrule
\end{tabular}
\end{table}

For the filter row, the headline uses $\Psi_{\rm pol}\le19^{-1}$; the
displayed minimum rank $194$ is sufficient by
\eqref{eq:filter-small-cushion}, and the all-target threshold includes
rejection through \eqref{eq:l2-rejection}.  For fixed skew, the headline uses $\Theta_{\rm pol}\le19^{-1}$; the
displayed $r_*\ge K+1$ is sufficient, while $r_*\ge2K$ makes every
target solvable.  The quotient matrix depends only on the
fixed public outer map and the relation, not on a generated key.  Exhaustive
enumeration of all $19^2=361$ relations $R_1=aI+bS$ proves quotient rank $K$
for every relation on all three parameter shapes.  Exhaustive optimization of
the natural two-block special-XL model gives bit-operation exponents $132.31$,
$184.27$, and $236.36$, all worse than the selected generic-MQ modes.

\section{Implementation and Evidence}\label{sec:evidence}

The artifact implements the public key generator, verifier-side feature maps,
symmetry quotient, affine reduction, structured lift, rank tests, and cost
searches over $\F_{19}$.  It does \emph{not} report a production-size
end-to-end forgery, a production-size Macaulay solve, or a wall-clock attack
timing.  Its full-size experiments concern public preprocessing and rank; its
solver experiments use smaller tractable profiles.  The checks therefore have
different logical force.  The table
below is the intended reading guide for every experimental statement in this
section.

\begin{table}[H]
\centering
\caption{Evidence hierarchy.  ``Exact'' here means exact for the concrete
object checked, not an exhaustive theorem about all production directions.}
\label{tab:evidence-hierarchy}
\small
\renewcommand{\arraystretch}{1.12}
\begin{tabularx}{0.96\textwidth}{@{}>{\raggedright\arraybackslash\bfseries}p{0.24\textwidth}>{\raggedright\arraybackslash}p{0.30\textwidth}>{\raggedright\arraybackslash}X@{}}
\toprule
Check & What is computed & What it does---and does not---establish\\
\midrule
Public preflight & Exact ranks and rejection conditions for the chosen key,
relation, offsets, and slice & The exact reduction is valid for that concrete
attack instance.\\
Seeded spectrum theorem & Exact min-entropy calculation in the random-XOF
idealization & Bounds the aggregate spectrum for the reserved-line construction;
not a deterministic claim about one already fixed SHAKE seed.\\
Rank campaign & Exact ranks for many selected nonzero output or direction
vectors & Finite fixed-key evidence and regression coverage; no longer the
logical basis of the structured aggregate-spectrum claim.\\
Macaulay experiment & Exact kernels at smaller tractable dimensions & Evidence
for the production solver model; not a production timing or rank proof.\\
Four-block and 3+1 frontiers & Exact Hilbert coefficients, row and column
counts, normal-operator ledgers, linear-completion rules, and monotone tail
certificates & Globally certified optima inside the stated solver classes;
not a production solve or wall-clock trace.\\
Frobenius-channel certificate & Exact rank-ten coordinate map, every square
core mix, B'ezout and SLP quantities, extension thresholds, Jacobian losses,
and break-even factors & Proves the channel algebra and numerical ledger; the
soft-$O$ implementation factor and resident memory remain explicit.\\
Basis-complete/square homotopy & Exact family-count enumeration, multihomogeneous
B\'ezout quantities, extension thresholds, full-square ledger, filters, and
spectrum-tail calculations & Removes the fixed-core and selected-degree
premises in the stated operation-count model; $\kappa_{\rm hom}$ and large
parametrization volume remain explicit.\\
Circuit validation & Exhaustive or independent checking of the released field
primitive & Exact cost and functionality of that primitive, but not of a
complete solver.\\
\bottomrule
\end{tabularx}
\end{table}

A ``projective direction'' below means a nonzero coefficient vector modulo
multiplication by a nonzero field scalar; all scalar multiples define the same
polar rank.  Sampling many distinct directions remains useful for regression
and fixed-key diagnostics.  The structured headline now has a proof over
the random-XOF idealization of \Cref{thm:seeded-spectrum}; an exact spectrum
computation remains the deterministic route for one fixed public key.  A
global minimum-rank proof is a still stronger sufficient certificate.

\subsection{Public preflight and verifier regression}

For a structured $\ell=4$ instance, the preflight checks
\[
 \rank Q_R=K,\qquad \rank C_{R,\mathbf v}=M-K,\qquad
 \rank\mathcal O_{R,\mathbf v}=m_1\ell^2,
\]
together with $\mathcal F_{R,\mathbf v}H_{R,\mathbf v}=0$, the predicted $A$-kernel
dimension, $A$-invariance, the restricted quadratic rank, and rejection
compatibility.  All 121 regenerated full-size instances tested in the artifact pass,
including the official Level-I known-answer key.  The first deterministic
offset candidate passes for 114 instances; the second passes for the remaining
seven.

For fixed skew, the preflight checks the relation and rejection identity and
the quotient, affine, and augmented ranks.  All 60 full-size keys tested in the artifact pass.  The artifact also reproduces the official Level-I public key byte for byte,
verifies its known-answer signature against the official hash target, and
checks the fixed output-order permutation described in
\Cref{sec:affine-columns}.  Eighteen additional regressions cover all nine
Version~2.3 shapes.  An independent Gauss--Jordan implementation agrees with
the rank kernel on 2,036 matrices, including 36 production-sized matrices.

\subsection{Rank evidence}

\paragraph{Seeded aggregate spectrum.}
The machine-readable certificate evaluates
\eqref{eq:seeded-eta-mean} and \eqref{eq:seeded-eta-tail} to 180 decimal digits for both
structured shapes at each level.  For the compact profiles it obtains
$\log_2$ bad-$1/19$ bounds $-55.22$, $-123.19$, and $-191.16$; for the robust
profiles it obtains bad-$1/2$ bounds $-41.48$, $-109.45$, and $-177.41$.
The calculation retains the actual maximum symbol atom $14/256$, rather than
pretending the byte-to-field conversion is uniform.  It is a theorem-driven
calculation from public dimensions, not a sampled-rank experiment.

\paragraph{Structured slices.}
Every selected slice is tested on all coefficient-basis directions (one
nonzero output coefficient at a time) and 1,024 deterministic dense directions
(many nonzero coefficients).  Selected slices also receive complete normalized
weight-two tests (two nonzero coefficients, modulo overall scaling) or
additional dense tests.  All 102,796 tested
cross-polar maps have maximal rank $h=4s$; the per-shape counts are in
\Cref{tab:cross-polar}.  The dual campaign evaluates 21,992 directional maps
$T_y$ on the same 13 slices, and every one is surjective of rank $K$.
All 124,788 coefficient and direction vectors used in these two campaigns are
projectively distinct within their respective slice records.

The directional campaign directly tests the simple sufficient certificate
for the aggregate headline condition.  The sufficient ranks are only
$h+1=41,61,77$; every tested $T_y$ instead has rank $K=50,70,90$.  Thus the
sampled directional margins are $9$, $9$, and $13$ ranks, respectively.
These finite tests certify neither the global minimum nor the weighted
aggregate, but they probe a condition substantially weaker than surjectivity.

\begin{table}[tbp]
\centering
\caption{Exactly computed cross-polar ranks for the tested structured $\ell=4$ slice--direction pairs.}
\label{tab:cross-polar}
\small
\setlength{\tabcolsep}{5.2pt}
\begin{tabular}{@{}l r r r r@{}}
\toprule
Parameters & slices & $h=4s$ & rank tests & sampled minimum\\
\midrule
$(28,5,4,4)$  & 3 & 40 & 26,296 & 40\\
$(28,4,4,5)$  & 2 & 40 & 24,198 & 40\\
$(40,7,4,4)$  & 2 & 60 & 45,658 & 60\\
$(38,5,4,5)$  & 2 & 60 & 2,188  & 60\\
$(50,9,4,4)$  & 2 & 76 & 2,228  & 76\\
$(52,6,4,6)$  & 2 & 76 & 2,228  & 76\\
\midrule
Total          & 13 & -- & 102,796 & maximal\\
\bottomrule
\end{tabular}
\end{table}

\paragraph{Fixed-skew systems.}
A broad campaign tests every coefficient-basis direction together with
combinations having two, three, or many nonzero output coefficients.  A second campaign targets
the 21 exceptional projective directions that make the local operator on
$\Mat_2(\F_{19})$ rank two rather than four.  It tests all 381 local
projective directions in one block, every exceptional direction in every
remaining block, and deterministic two-block exceptional combinations.
Together the campaigns perform 3,888 exact rank evaluations; their sampled
minima and the budget, sufficient, and all-target thresholds appear in
\Cref{tab:fixed-skew-evidence}.

\begin{table}[tbp]
\centering
\caption{Full-size polar-rank evidence for the fixed-skew $\ell=2$ systems.
``Budget'' preserves positive AXN nominal arithmetic gap; ``suff.'' is the simple
minimum-rank certificate for the aggregate headline condition in
\Cref{tab:l2-costs}; ``all target'' guarantees a root for every target.}
\label{tab:fixed-skew-evidence}
\small
\setlength{\tabcolsep}{3.3pt}
\begin{tabular}{@{}l r r r r r r@{}}
\toprule
Parameters & preflights & tests & sampled min. & budget & suff. & \shortstack{all\\target}\\
\midrule
$(48,16,2,2)$ & 20 & 1,104 & 110 & 46 & 49 & 96\\
$(72,24,2,2)$ & 20 & 1,296 & 167 & 68 & 73 & 144\\
$(96,32,2,2)$ & 20 & 1,488 & 223 & 89 & 97 & 192\\
\midrule
Total & 60 & 3,888 & -- & -- & -- & --\\
\bottomrule
\end{tabular}
\end{table}

\paragraph{Zero-offset filter.}
The broad and exceptional-direction campaigns perform 29,424 exact rank
evaluations on two keys per shape.  Their sampled minima are 127, 190, and 254
in dimensions 128, 192, and 256.  For the Level-V filter, the sampled minimum
254 exceeds the simple sufficient minimum-rank value 194 by 60 and the
all-target threshold 240 by 14.  If the global minimum is at least 254, then
\eqref{eq:l2-accepted-fraction} gives an accepted fraction of at least
$0.999999999$ in every target fiber.  The experiment supports that hypothesis
but does not certify the global minimum.

\subsection{Solver and probability checks}

The special-XL search examines every cost-competitive multidegree and
recovers the three time optima in \Cref{tab:xl-optima}.  Separate exact searches
verify every candidate that can match the compact and robust state incumbents
and close all larger degrees by coordinatewise monomial-count monotonicity,
proving \Cref{prop:low-state-frontier,prop:explicit-stream-frontier,prop:robust-explicit-frontier}.
They record exact Hilbert coefficients, $M$, $R$, extension-field
preconditioner bounds, operator constants, and complete array ledgers.  A
separate exhaustive search proves \Cref{prop:three-plus-one-frontier} from the
three-block series \eqref{eq:hilbert-series-three}; its certificate also checks
the exact fourth-block equation counts and the rank-feasibility inequality
$s\le3m_1$.  No production three-block quotient or linear-completion rate is
claimed.  At tractable sizes,
the artifact builds 128 four-block Macaulay matrices: 64 random instances have
full column rank and 64 planted-solvable instances have the model's strongest
$D=1$ behavior, with the planted evaluation vector spanning the right kernel.
A separate boundary campaign constructs 48 planted four-block systems across
five negative-Hilbert profiles having at least one degree-one block; every
matrix has exact right corank one.  This confirms that degree one still retains
all variables and behaves normally in small profiles, but it is not a
production-degree theorem.
  A further reproducible campaign keeps the
\emph{production-selected multidegrees} $(3,3,3,4)$, $(1,2,3,5)$, and
$(2,2,2,3)$ and the corresponding values $m_1=5,7,9$, while reducing the
block size to $s=2$ so exact sparse elimination is feasible.  The three
planted matrices have $(M,R)=(15{,}000,268{,}500)$,
$(3{,}780,69{,}678)$, and $(2{,}160,47{,}628)$ and all have exact right
corank one.  At reduced size only the first Hilbert coefficient is negative,
so this check is evidence about the selected degree shapes, not an argument
that the production solving degree is proved.  A 48-matrix two-block control contains one planted
case with a two-dimensional right kernel.  The latter is no longer a conceptual failure---finite quotients of
dimension larger than one are covered by
\Cref{thm:finite-quotient-recovery}---but we still do not promote the
$\ell=2$ special-XL estimates because the control does not establish border
completeness or production degree.  These are generic small-profile
experiments, not production SNOVA timings.

The channel artifact independently constructs
\eqref{eq:frobenius-channel-isomorphism} in
$A=\F_{19}[t]/(t^4-t-1)$; its $10\times10$ matrix has exact rank ten and
determinant eight.  A second exact script enumerates every whole-channel square
core for all three levels, recomputes $B$, $B^+$, the SLP bound, the finite-
cardinality extension degree, the Jacobian loss, both retry denominators, the
packed output size, and the dense coefficient-volume ceiling.  It proves the
globality statement in \Cref{prop:channel-homotopy-ledger}.  This is an exact
arithmetic certificate, not a symbolic-homotopy implementation or a claim that
the hidden multiplier $\kappa_{\rm hom}$ is small.

Exhaustive checks over $\F_3$ and $\F_5$ verify the exact first and factorial
second moments, the compatible-collision formula, the rank-spectrum upper bound,
and the singleton and planted-distribution inequalities on 110 complete
quadratic maps.  Seventeen maps have nonmaximal cross-polar rank; in twelve,
compatibility makes $\eta_0$ strictly smaller than $\bar\eta$.  Finally, on six actual cost-driving $\ell=2$ cores, the homogeneous
quadratic outputs are linearly independent, the common polar radical---the set
of directions annihilated by every polar form---is zero, and all 384 sampled
homogeneous Jacobians (matrices of first derivatives of the homogeneous
quadratic parts) attain the maximum rank permitted by their dimensions.  These
are genericity diagnostics, not solving-degree measurements.

\subsection{Scope of the estimates}

Once a public preflight passes, the symmetry quotient, affine reduction,
stable kernel, fixed-skew identity, and rejection handling are exact.  For a
structured slice, the root quantity that matters is $\bar\eta_L$.
\Cref{thm:seeded-spectrum} proves the required aggregate bound with explicit
failure probability in the random-XOF idealization, provided the candidate is
constructed from the decoupled cross data.  The complete projective rank
histogram remains an exact deterministic certificate for a fixed key, and
direct enumeration is dominated by the Level-III/V compact solves.  Thus the
structured rank spectrum is no longer an unproved sampled minimum; the model
boundary is whether one accepts the random-XOF idealized-expansion theorem or asks
for a fixed-seed certificate.

The remaining production assumption for the streamed-XL branch is solver
behavior at the selected Macaulay degree.  The compact scalar-Wiedemann rows use $\sigma=1$ in
\Cref{thm:las-vegas-wrapper}; the robust explicit rows charge only
$\sigma\ge1/8$.  A complete commuting border basis of dimension $D>1$ still
gives zero-error recovery by \Cref{thm:finite-quotient-recovery}, but the extra
nullspace/block work needed to obtain it must fit the displayed operator
ledger.  The $D_{30}$ values in \Cref{tab:quotient-extraction} price only
FGLM-style postprocessing once the multiplication tables have been supplied.

The fixed-core-free branch also does not share that selected-degree premise.
The same aggregate spectrum that controls collisions bounds full-Jacobian
singularity by \eqref{eq:jacobian-defect-mass}.  The basis-complete rows run
every possible multidegree-family count; the Level-III row solves the complete
square residual system.  Thus their root-completeness statement needs neither
a production Macaulay degree nor a preselected nonzero Jacobian determinant.
Their quantitative boundary is instead the explicit $\kappa_{\rm hom}$ and
the potentially enormous zero-dimensional parametrization.

The Frobenius-channel branch does not share that selected-degree premise.  Its
exact preflight is a nonzero evaluation of the highest homogeneous Jacobian
determinant.  Conditional on that public certificate, \Cref{thm:accepted-nonsingular-singleton}
accounts for singular roots without independence, and
\Cref{lem:finite-cardinality-homotopy} supplies all nonsingular core roots with
constant success probability.  Its remaining quantitative boundary is
instead the explicit $\kappa_{\rm hom}$ multiplying a soft-$O$ arithmetic
bound, together with the fact that the cited theorem gives no low-memory
implementation guarantee.  The paper therefore treats the channel exponents
as break-even sensitivity ledgers rather than production timings.

The state and operator accounting is no longer implicit.  The streamed normal
operator of \Cref{prop:streaming-macaulay} avoids Macaulay materialization.
The compact four-block frontier charges its exact row count over
$\F_{19^8}$ and has packed eight-array envelopes $17.7$~GiB,
$2.21$~GiB, and $1.01$~GiB.  The three-plus-one route gives explicit Level-III/V time points
$2^{197.02}$ and $2^{250.32}$ and a robust Level-V low-state point
$2^{263.40}$ with $9.55$~GiB.  Taking the best joint-stress state point from either XL decomposition gives
packed envelopes $1.53$~TiB, $55.4$~GiB, and $9.55$~GiB.  The larger Level-I field affects only random normal
preconditioning.  These are explicit state and arithmetic bounds, not
production timings.

For $\ell=2$, the headline uses the aggregate bounds
$\Theta_{\rm pol}\le19^{-1}$ and $\Psi_{\rm pol}\le19^{-1}$; minimum ranks
$49$, $73$, and $194$ suffice, versus all-target values $96$, $144$, and
$240$.  The sampled ranks and small-profile solves support these conditions but
do not certify the complete spectra or the selected production solving
degrees.  No production-size end-to-end solve is claimed.

\paragraph{What would turn the estimates into end-to-end certificates.}
For a fixed slice, the exact target is the weighted sum
$\bar\eta_L=\sum_{y\ne0}19^{-\rank T_y}$.  The reserved-line theorem already
controls this sum under the idealized key distribution.  When a deterministic fixed-seed record
is desired, \Cref{prop:projective-rank-certificate} gives a finite exact
certificate: disjoint normalized projective charts, finite-field equations,
determinantal ideals, and quotient dimensions recover the complete rank
histogram.  Unit-ideal transcripts for the $(h+1)\times(h+1)$ minors give the
stronger minimum-rank certificate.  Direct enumeration is also exact and is
already dominated by the Level-III and Level-V compact solves.  The
random-slice identity in \Cref{prop:random-A-slice} gives another route on the
stable kernel.  For
the full-space $\ell=2$ rows, the same projective construction gives
$\Theta_{\rm pol}$ and $\Psi_{\rm pol}$ exactly; the displayed
minimum-rank bounds remain stronger sufficient certificates.

For the channel route, a fixed-key record need only add the exact
$10\times10$ channel transform, the selected whole-channel core, one nonzero
leading-Jacobian witness, and the output of a finite-field symbolic-homotopy
implementation followed by all public filters.  This would directly measure
$\kappa_{\rm hom}$ and resident memory; it would not require any Macaulay
selected-degree extrapolation.

For the fixed-core-free route, the corresponding record is even simpler
logically: publish the ordinary or stable slice, the random within-family
projection seeds for every pattern, the homotopy outputs, and the final public
filter transcript.  No leading-Jacobian witness is part of the premise; the
successful pattern's nonsingularity is checked directly at the returned root.

For the streamed-XL solver, a full production run is sufficient but not logically
necessary.  Row-reduction certificates showing the border relations of
\Cref{thm:finite-quotient-recovery} at the selected degree---or at a nearby
degree still within the sensitivity budgets---would give a machine-checkable
end-to-end root-extraction certificate.  The natural production experiment is
now sharply specified: use the streamed operator of
\Cref{prop:streaming-macaulay}, record the preconditioner seed, resident block
size, black-box products, border relations, quotient dimension, and
verifier-checked roots.  For the three-plus-one route it should additionally
record every partial root, the consistency and rank of
\eqref{eq:three-plus-one-linear}, and the final diagonal checks.  A complete
circuit comparison would additionally need
to price control flow, indexing, memory traffic, communication, and data
movement.

\section{Countermeasures}\label{sec:countermeasures}

The design lesson is that security must be based on distinct quadratic
polynomials, not on the number of ordered labels used to name them.  A security
analysis for any symmetric redesign should compute the actual quotient rank
$\rho_{\rm quad}$ and use the unconditional cap
\begin{equation}\label{eq:countermeasure-K}
 \rho_{\rm quad}\le K_d
 =\min\left\{M,m_1\binom{d+1}{2}\right\},
 \qquad d=\dim_{\F_q}\F_q[S].
\end{equation}
Equality with $K_d$ is a public rank condition, not a consequence of symmetry
alone.  Outer randomization does not restore the alternating directions removed
before the outer map.

\begin{proposition}[Necessary floor for restoring the ordered-label cap]
\label{prop:repair-floor}
Let an existing ordered-label analysis intend the quadratic-image cap
\[
 K_{\rm ord}=\min\{M,m_1d^2\}.
\]
Any redesign that retains symmetric public forms and uses $m_1'$ base-form
groups and $M'$ verifier outputs can restore even this cap only if
\begin{equation}\label{eq:repair-necessary}
 M'\ge K_{\rm ord},
 \qquad
 m_1'\ge
 \left\lceil\frac{K_{\rm ord}}{\binom{d+1}{2}}\right\rceil.
\end{equation}
If it also retains Version~2.3's coupling
\[
 x'=o'r',\qquad M'=dx',\qquad m_1'=\left\lceil\frac{x'}d\right\rceil,
\]
then necessarily
\begin{equation}\label{eq:repair-x-floor}
 x'\ge
 \max\!\left\{
 \left\lceil\frac{K_{\rm ord}}d\right\rceil,
 d\left(
 \left\lceil\frac{K_{\rm ord}}{\binom{d+1}{2}}\right\rceil-1
 \right)+1
 \right\}.
\end{equation}
These are necessary feature-count conditions only.
\end{proposition}

\begin{proof}
The quadratic image of any symmetric redesign has dimension at most
$\min\{M',m_1'\binom{d+1}{2}\}$.  Requiring this upper bound to reach
$K_{\rm ord}$ forces both inequalities in \eqref{eq:repair-necessary}.  The
least integer $x'$ satisfying $\lceil x'/d\rceil\ge m_1'$ is
$d(m_1'-1)+1$; combining this with $dx'\ge K_{\rm ord}$ gives
\eqref{eq:repair-x-floor}.
\end{proof}

For every published row, $K_{\rm ord}=M$.  Applying
\Cref{prop:repair-floor} gives the following unconditional minimum changes for
a redesign that chooses to restore that cap while retaining the current
coupling.  A secure parameter set may require still larger changes, and the
table does not prescribe how $x'=o'r'$ should be split.

\begin{table}[H]
\centering
\caption{Necessary feature-count floor for restoring the ordered-label cap under
$m_1=\lceil or/\ell\rceil$.  This restores only the quadratic-image cap
implicit in the ordered-label analysis; it is not a security estimate.}
\label{tab:repair-floor}
\scriptsize
\setlength{\tabcolsep}{3.2pt}
\begin{tabular}{@{}c l c c c c@{}}
\toprule
Level & parameters & $m_1\to m_1'$ & $or\to o'r'$ & $M\to M'$ & increase\\
\midrule
I   & $(28,5,4,4)$  & $5\to8$   & $20\to29$ & $80\to116$  & $45.00\%$\\
I   & $(48,16,2,2)$ & $16\to22$ & $32\to43$ & $64\to86$   & $34.38\%$\\
I   & $(28,4,4,5)$  & $5\to8$   & $20\to29$ & $80\to116$  & $45.00\%$\\
\midrule
III & $(40,7,4,4)$  & $7\to12$  & $28\to45$ & $112\to180$ & $60.71\%$\\
III & $(72,24,2,2)$ & $24\to32$ & $48\to63$ & $96\to126$  & $31.25\%$\\
III & $(38,5,4,5)$  & $7\to10$  & $25\to37$ & $100\to148$ & $48.00\%$\\
\midrule
V   & $(50,9,4,4)$  & $9\to15$  & $36\to57$ & $144\to228$ & $58.33\%$\\
V   & $(96,32,2,2)$ & $32\to43$ & $64\to85$ & $128\to170$ & $32.81\%$\\
V   & $(52,6,4,6)$  & $9\to15$  & $36\to57$ & $144\to228$ & $58.33\%$\\
\bottomrule
\end{tabular}
\end{table}

\begin{proposition}[Necessary output floor for removing the full fresh slice]
\label{prop:fresh-slice-repair-floor}
Consider an $\ell=4$ symmetric redesign that retains the reserved public
vinegar line and the same public vinegar count $v$.  If
\begin{equation}\label{eq:fresh-slice-dimension-condition}
 M\le \ell(v-1),
\end{equation}
then every full-rank affine reduction satisfies
$\dim_{\F_q}(V'\cap\ker C_{R,\mathbf v})\ge K$ and therefore retains the
$K$-dimensional ordinary square-slice route of
\Cref{cor:full-ordinary-slices}.  Eliminating this route by dimensions alone
requires
\begin{equation}\label{eq:fresh-slice-output-floor}
 M>\ell(v-1).
\end{equation}
If $M$ remains a multiple of $\ell$, then $M\ge\ell v$.
\end{proposition}

\begin{proof}
By \eqref{eq:D0-lower-bound},
\[
 \dim D_0\ge\ell(v-1)-(M-K)
 =K+\ell(v-1)-M.
\]
Condition \eqref{eq:fresh-slice-dimension-condition} therefore implies
$\dim D_0\ge K$.  Negating that dimension guarantee gives
\eqref{eq:fresh-slice-output-floor}; divisibility by $\ell$ rounds the strict
inequality to $M\ge\ell v$.
\end{proof}

For the six public $\ell=4$ rows, removing the route by this dimension test
alone forces the following minimum output increases.  These are necessary
conditions only: satisfying the table does not prove security, and the actual
intersection $D_0$ may exceed its lower bound.

\begin{table}[H]
\centering
\caption{Necessary output increase to eliminate the guaranteed
$K$-dimensional fresh ordinary slice by dimensions alone, while retaining
$v$ and $M\equiv0\pmod\ell$.}
\label{tab:fresh-slice-repair-floor}
\small
\setlength{\tabcolsep}{4.0pt}
\begin{tabular}{@{}c l r r r@{}}
\toprule
Level & parameters & current $M$ & necessary $M'$ & increase\\
\midrule
I   & $(28,5,4,4)$  & 80  & 112 & $40.00\%$\\
I   & $(28,4,4,5)$  & 80  & 112 & $40.00\%$\\
III & $(40,7,4,4)$  & 112 & 160 & $42.86\%$\\
III & $(38,5,4,5)$  & 100 & 152 & $52.00\%$\\
V   & $(50,9,4,4)$  & 144 & 200 & $38.89\%$\\
V   & $(52,6,4,6)$  & 144 & 208 & $44.44\%$\\
\bottomrule
\end{tabular}
\end{table}

Changing $q$ or rejecting symmetric signature blocks does not change
\eqref{eq:fresh-slice-dimension-condition}.  A redesign can instead alter the
affine geometry, remove the reserved-line separation, reduce the public
vinegar space, or increase outputs beyond this floor; each such change must be
reanalyzed rather than assumed secure.

This table does not claim that the team's $q=19$ choice was intended to
preserve $K_{\rm ord}$: the 2025 presentation explicitly switched to the
$o\binom{\ell}{2}$ alternating-form criterion
\cite{WangEtAl2025Round2Security}.  It answers a narrower but unavoidable
countermeasure question.  A response that retains symmetric public forms and
the current coupling cannot recover even the old ordered-label image cap with
smaller changes; increasing $q$ or rejecting symmetric signatures does not
recreate the missing alternating directions.

The direct repairs are to increase parameters after accounting for
\eqref{eq:countermeasure-K},
remove public-matrix symmetry, or redesign the left and right power actions so
that exchanging their labels changes the quadratic polynomial.  The
symmetric-block rejection rule is not sufficient: the fixed-skew $\ell=2$
lift makes every block nonsymmetric by construction, while an injective local
skew map for the square $\ell=4$ rows leaves accepted density at least
$(1-19^{-4})^n$ and adds under $0.0014$ bits to the priced retry factor.  Returning
unchanged to the submitted $q=16$ design would instead restore the known
characteristic-two attack surface.

A revised security analysis should publish the quotient rank, affine or filter
rank, stable-kernel rank when applicable, the polar or cross-polar condition
used for target availability, and solver logs for the actual cost-driving
system.

\section{Conclusion}\label{sec:conclusion}

The proposed odd-characteristic Version~2.3 construction has an unconditional
quadratic-image defect.  Its ordered power-pair labels do not represent
independent polynomials: symmetry identifies $(\xi,\eta)$ with
$(\eta,\xi)$ before the outer map, so every public base-form group factors
through $\Sym^2(A)$.  For $\ell=4$, sixteen labels become ten; for $\ell=2$,
four become three.  Exact affine elimination turns this identity into complete
public-key forgery reductions, and fixed-skew, filtered, or injective-skew
lifts defeat the proposed symmetric-signature rejection rule with exact
acceptance bounds.

The former global-rank objection is no longer a sampled premise.  A reserved
public vinegar line fixes the attack geometry from cross data disjoint from
fresh internal public coefficients.  The random-XOF theorem bounds the full
aggregate spectrum, and its excess form separates the unavoidable
full-rank baseline from genuine rank defects.  The same spectrum bounds the
full-Jacobian singular locus projectively.  Consequently, an accepted
nonsingular root occurs with the second-moment probability in
\eqref{eq:accepted-nonsingular-success}; no singleton assumption and no
independence among rejection, singularity, and the MQ map is required.  A fixed
public seed can instead be accompanied by its exact finite projective rank
histogram.

This yields a selected-degree-free, fixed-core-free frontier.  Frobenius
transport collapses the $55$ near-square family-count patterns to $16$ orbits.
The robust per-certified-key costs are $2^{138.324}$, $2^{179.758}$, and
$2^{220.832}$ AXN operations.  An exact structural-preflight theorem proves
success probability above $0.9423222679$ for every official $\ell=4$ shape in
the code-faithful random-XOF model; charging the full inverse probability gives
$2^{138.410}$, $2^{179.844}$, and $2^{220.918}$.  This branch has neither a
selected-degree nor a preferred-core premise, although its direct dense output
can be enormous.  The independent original-equation eigenblock-core route has
much smaller output ceilings and robust per-certified-key exponents
$134.077$, $196.246$, and $247.432$, conditioned on an exact public nonzero
Jacobian certificate.  Both homotopy routes are operation-count theorems, not
measured wall-clock claims.

The streamed-XL route supplies the complementary low-state result.  Its
state-minimizing four-block profiles cost $2^{141.70}$, $2^{203.63}$, and
$2^{269.35}$ AXN operations with packed eight-array envelopes $17.7$~GiB,
$2.21$~GiB, and $1.01$~GiB.  The exact three-plus-one decomposition gives a
robust Level-V point $2^{263.40}$ with $9.55$~GiB.  This branch still depends
on production selected-degree certificate behavior; correctness itself needs
only a complete finite commuting border quotient, not exact corank one.

The countermeasure burden is therefore not a small rejection tweak.  Merely
restoring the old ordered-label quadratic-image cap while retaining symmetry
and Version~2.3's coupling forces $31.25\%$--$60.71\%$ more verifier outputs
across the nine rows.  On the six $\ell=4$ rows, eliminating the guaranteed
full-dimensional fresh ordinary slice by dimensions alone forces
$38.89\%$--$52.00\%$ more outputs.  These are unconditional necessary floors,
not sufficient repaired parameters.  Removing public-matrix symmetry,
changing the left/right actions or affine geometry, or paying a substantial
parameter increase are the direct responses; increasing $q$ or rejecting
symmetric signature blocks does not recreate the missing alternating
features.

The remaining empirical questions are sharply isolated.  A homotopy
implementation would measure $\kappa_{\rm hom}$ and resident output state; a
streamed border-basis transcript would settle the production XL degree.  Those
questions affect the best concrete implementation claim, but not the
symmetry-quotient theorem, the exact forgery reductions, the rejection bypass,
the selected-degree-free root algorithms, or the necessary parameter floors.

\appendix
\setcounter{section}{0}
\renewcommand{\thesection}{\Alph{section}}
\renewcommand{\theHsection}{app.\Alph{section}}


\section{Proofs Deferred from the Main Text}\label{app:deferred-proofs}

The main text states the probability and structural results at the point where
an attacker uses them.  This appendix collects the longer proofs in the same
order.

\subsection{Stable-lift kernel}\label{app:proof-stable-kernel}

\begin{proof}[Proof of \Cref{thm:stable-kernel}]
Choose an $\F_q$-basis of $A$.  The spanning family in \eqref{eq:Uw} has
$m_1d^2$ elements, proving the dimension bound.  For $\zeta\in A$,
\[
 \bigl(w^\transpose\Sigma_\xi P_i\Sigma_\eta\bigr)\Sigma_\zeta
 =w^\transpose\Sigma_\xi P_i\Sigma_{\eta\zeta},
\]
so $\mathcal U_w$ is closed under right multiplication by $A$.

A cross term produced by substituting
$u_j=\Sigma_{R_j}u+\Sigma_{\Gamma_j}w$ has, when $w$ is on the left, a row
coefficient $w^\transpose\Sigma_\xi P_i\Sigma_\eta$.  When $w$ is on the
right, transpose the scalar and use $S^\transpose=S$,
$P_i^\transpose=P_i$, and commutativity of $A$; the coefficient has the same
form with the two algebra labels reversed.  Thus every row of
$\mathcal F_{R,\mathbf v}$ lies in $\mathcal U_w$, and right-$A$ stability
puts every row of $\mathcal O^{\mathcal F}_{R,\mathbf v}$ there as well.

Because \eqref{eq:orbit-matrix} contains the complete right orbit of every row
of $\mathcal F_{R,\mathbf v}$, its row space is closed under right
multiplication by every element of $A$.  Since $A\cong\F_{q^d}$ is a field,
this row space is an $A$-vector space.  Hence $d\mid\rho_{\mathcal F}$ and
\[
 \rho_{\mathcal F}\le\dim_{\F_q}\mathcal U_w\le m_1d^2.
\]
The first block of \eqref{eq:orbit-matrix} is
$\mathcal F_{R,\mathbf v}$ itself, so
$\mathcal F_{R,\mathbf v}\ker\mathcal O^{\mathcal F}_{R,\mathbf v}=0$.
For $x$ in this kernel and $\zeta\in A$, every product $S^a\zeta$ is an
$\F_q$-linear combination of $1,S,\ldots,S^{d-1}$; therefore every row
$ f\Sigma_{S^a}$ of the orbit matrix also kills $\Sigma_\zeta x$.  Thus the
kernel is $A$-linear.  Rank--nullity over $A$ gives
\[
 \dim_A H^{\mathcal F}_{R,\mathbf v}
 =n-\rho_{\mathcal F}/d\ge n-m_1d,
\]
as claimed.
\end{proof}

\subsection{Stable affine slices}\label{app:proof-stable-slices}

\begin{proof}[Proof of \Cref{cor:stable-slices}]
The first two ranks in \eqref{eq:stable-preflight} give the full-lift case of
\Cref{prop:affine-reduction}; in particular, every target determines a
nonempty affine space $a(t)+\ker C_{R,\mathbf v}$.  The orbit matrix contains
$C_{R,\mathbf v}$ as its first block, so
$H_{R,\mathbf v}\subseteq\ker C_{R,\mathbf v}$.  Hence
$a(t)+y+L\subseteq a(t)+\ker C_{R,\mathbf v}$ whenever
$y\in\ker C_{R,\mathbf v}$.  Because
$L\subseteq H_{R,\mathbf v}$, \Cref{thm:stable-kernel} gives
$\mathcal F_{R,\mathbf v}L=0$, and therefore $\mathcal F_{R,\mathbf v}u$ is constant on that
slice.

By \Cref{lem:spectral-quotient}, the unordered power-pair features can be
changed invertibly to eigenblock coordinates.  Because $Q_R$ has full column
rank, row reduction of the complete output system isolates $K$ rows whose
homogeneous quadratic parts are an invertible linear image of
$z_{\rm sym}$.  Applying the inverse image change to those rows expresses the
quadratic parts in eigenblock coordinates.  Full rank of the quadratic feature
map restricted to $L$ ensures that this restriction loses no eigenblock
component.  Translating a diagonal feature introduces only
quadratic, linear, and constant terms from one block; translating an
off-diagonal feature introduces terms from its two blocks.  If $h_a$ is a
homogenizer for block $a$, multiplication by the unique powers of the $h_a$
turns these affine polynomials into exact multidegrees $2\mathbf e_a$ and
$\mathbf e_a+\mathbf e_b$.

Distinct unordered multidegrees have disjoint monomial supports.  Each
component is generated by the $m_1$ base forms and therefore has rank at most
$m_1$.  The total restricted rank is the sum of these component ranks.  Since
there are $\binom{\ell+1}{2}$ components and the total is
$K=m_1\binom{\ell+1}{2}$, every component must have rank exactly $m_1$.
\end{proof}

\subsection{Accepted targets for the \texorpdfstring{$\ell=2$}{l=2} filter}\label{app:proof-l2-filter}

\begin{lemma}[Quadratic fiber bound]\label{lem:quadratic-fiber-bound}
Let $q$ be odd and let $g:\F_q^N\to\F_q^K$ be quadratic.  For each
$b\in\F_q^K\setminus\{0\}$, let $r(b)$ be the polar rank of the scalar
quadratic $b\cdot g$, put $r_*=\min_{b\ne0}r(b)$, and set
$\alpha=1-q^{-K}$.  Then every target $z\in\F_q^K$ satisfies
\begin{equation}\label{eq:quadratic-fiber-bound}
 \left|\,|g^{-1}(z)|-q^{N-K}\right|
 \le \alpha q^{N-r_*/2}.
\end{equation}
More generally, if $W\subseteq\F_q^N$ is an affine subspace of codimension
$c$, then
\begin{equation}\label{eq:quadratic-affine-upper}
 |\{x\in W:g(x)=z\}|
 \le q^{N-c-K}+\alpha q^{N-r_*/2}.
\end{equation}
\end{lemma}

\begin{proof}
Fix a nontrivial additive character $\psi$ of $\F_q$.  Fourier inversion gives
\[
 |g^{-1}(z)|
 =q^{-K}\sum_{b\in\F_q^K}\psi(-b\cdot z)
   \sum_{x\in\F_q^N}\psi(b\cdot g(x)).
\]
The $b=0$ term is $q^{N-K}$.  For $b\ne0$, let $r=r(b)$.  An invertible
linear change of variables splits the polar form into a nondegenerate
$r$-dimensional part and a radical of dimension $N-r$.  If the linear part of
$b\cdot g$ is nonzero on the radical, summing over the radical gives zero.
Otherwise the radical contributes $q^{N-r}$, while diagonalizing the
nondegenerate quadratic part reduces its sum to a product of one-dimensional
Gauss sums of absolute value $q^{1/2}$.  In either case the inner sum has
absolute value at most $q^{N-r/2}$~\cite{Lam2005}.  Summing the $q^K-1$
nonzero Fourier coefficients proves \eqref{eq:quadratic-fiber-bound}.

Write $W=a+U$, where $U$ has codimension $c$.  Translation does not change a
quadratic polynomial's polar form.  Restricting a bilinear form to
$U\times U$ lowers its rank by at most $2c$: after extending a basis of $U$ to
one of $\F_q^N$, deleting the last $c$ rows and columns changes matrix rank by
at most $2c$.  Thus every nonzero scalar combination of $g|_W$ has polar rank
at least $r_*-2c$.  If $r_*\ge2c$, the same Fourier calculation on the
$(N-c)$-dimensional affine space $W$ bounds each nonzero character sum by
\[
 q^{(N-c)-(r_*-2c)/2}=q^{N-r_*/2}.
\]
If $r_*<2c$, the trivial bound $q^{N-c}$ is already at most
$q^{N-r_*/2}$.  This proves \eqref{eq:quadratic-affine-upper} in both cases.
\end{proof}

\begin{proof}[Proof of \Cref{prop:l2-rejection}]
By \Cref{lem:quadratic-fiber-bound}, every target fiber has between
$q^{N-K}-\alpha q^{N-r_*/2}$ and
$q^{N-K}+\alpha q^{N-r_*/2}$ roots.  A rejected root lies on
the intersection of some $t_{\rm rej}$ block-symmetry hyperplanes.  Each such
intersection is an affine subspace of codimension $t_{\rm rej}$, so
\eqref{eq:quadratic-affine-upper} bounds the number of roots it contains by
$q^{N-t_{\rm rej}-K}+\alpha q^{N-r_*/2}$.  Subtracting the union bound
over the $B$ intersections from the full-fiber lower bound gives
\[
 q^{N-K}(1-Bq^{-t_{\rm rej}})
 -\alpha q^{N-r_*/2}(1+B).
\]
This is positive exactly under the sufficient inequality
\eqref{eq:l2-rejection}.

Every fiber has size at most
$q^{N-K}(1+\alpha q^{K-r_*/2})$.  Since
$|\mathcal X_{\rm acc}|=a_n q^N$, double-counting accepted assignments by their images
proves \eqref{eq:l2-accepted-target}.  Finally, let
$T_z=|g^{-1}(z)|$ and let $R_z$ be the number of rejected roots.  The accepted
fraction is $1-R_z/T_z$.  Applying the union-bound upper estimate to $R_z$ and
the character-sum lower estimate to $T_z$ gives
\eqref{eq:l2-accepted-fraction}.
\end{proof}

\subsection{Root collisions and the rank-spectrum bound}\label{app:proof-coset-success}

\begin{proof}[Proof of \Cref{thm:coset-success}]
Average over the $q^{\dim V-h}$ cosets and $q^K$ targets.  Every point of $V$
contributes once to the first moment, which gives \eqref{eq:EZ}.  For an
ordered pair of distinct points in one coset, write the second point as $x+y$
with $y\in L\setminus\{0\}$.  The collision equation is
\[
 g(x+y)=g(x)
 \quad\Longleftrightarrow\quad
 T_y(x)=-(g(y)-g(0)).
\]
It has $\kappa(y)q^{\dim V-r(y)}$ solutions.  Dividing the total number of
ordered collisions by the number of target--coset pairs proves
\eqref{eq:factorial-moment}.

For the rank-spectrum identity, note that
$q^{-r(y)}=q^{-K}|\ker T_y^*|$ and double-count pairs $(b,y)$ satisfying
$b\cdot T_y=0$.  The term $b=0$ contributes $q^h-1$.  For $b\ne0$, the number
of nonzero $y\in L$ annihilated by $b\cdot T_y$ is
$q^{h-e_L(b)}-1$.  Therefore
\[
 \bar\eta
 =q^{-K}\left((q^h-1)+\sum_{b\ne0}(q^{h-e_L(b)}-1)\right)
 =\mu(1+\Theta_L(g))-1.
\]
The inequality $\eta_0\le\bar\eta$ is immediate.

Using $\mathbb E[Z^2]=\mu(1+\eta_0)$ gives
\[
 \operatorname{Var}(Z)
 =\mu(1+\eta_0-\mu)
 \le\mu^2\Theta_L(g),
\]
and the second-moment inequality proves \eqref{eq:coset-success}.  Finally,
$e_*\ge h-\Delta$ implies
$\Theta_L(g)\le(q^K-1)q^{-h+\Delta}<q^{\tau+\Delta}$.
\end{proof}

\subsection{Unique-root trials and the planted distribution}\label{app:proof-singleton-coset}

\begin{proof}[Proof of \Cref{cor:singleton-coset}]
For every nonnegative integer $z$,
\[
 \mathbf 1_{z=1}\ge z-z(z-1),\qquad
 \mathbf 1_{z>0}\ge z-\binom z2,\qquad
 \mathbf 1_{z\ge2}\le\binom z2.
\]
Apply these inequalities to $Z$ and use
\eqref{eq:factorial-moment} together with $\eta_0\le\bar\eta$.
Dividing the upper bound on $\Pr[Z\ge2]$ by
\eqref{eq:bonferroni-success} proves \eqref{eq:conditional-multiple}.

The planted law size-biases a target--coset pair by $Z$.  Hence
\[
 \Pr_{\mathsf P}[Z\ge2]
 =\frac{\mathbb E[Z\mathbf 1_{Z\ge2}]}{\mu}
 \le\frac{\mathbb E[Z(Z-1)]}{\mu}=\eta_0.
\]
Conditioned on $Z=1$, both $\mathsf U$ and $\mathsf P$ are uniform on the same
set of target--coset pairs.  If two distributions have the same conditional
law on an event, their total variation distance is at most the larger mass they
assign to its complement.  Here the complement is $\{Z\ge2\}$.  The same indicator inequality gives the sharper exact-parameter lower bound
$\Pr[Z>0]\ge\mu(1-\eta_0/2)$.  Combining it with
$\Pr[Z\ge2]\le\mathbb E[Z(Z-1)]/2=\mu\eta_0/2$ gives
\[
 \Pr_{\mathsf U}[Z\ge2]
 \le \frac{\eta_0}{2-\eta_0}\le\eta_0
 \qquad(\eta_0<1).
\]
The displayed size-bias bound gives
$\Pr_{\mathsf P}[Z\ge2]\le\eta_0$.  Both are at most $\bar\eta$, proving the
total-variation claim.
\end{proof}

\subsection{Full-space target availability}\label{app:proof-full-space}

\begin{proof}[Proof of \Cref{cor:full-space}]
Take $L=V$ in \Cref{thm:coset-success}.  Then $e_L(b)$ is the polar rank and
$\Theta_L(g)\le(q^K-1)q^{-r_*}=\alpha q^{K-r_*}$.  For the last statement, \eqref{eq:quadratic-fiber-bound} gives
\[
 |g^{-1}(z)|\ge q^{N-K}-\alpha q^{N-r_*/2}>0
\]
when $r_*\ge2K$.
\end{proof}

\section{Dimension Formulas}\label{app:dimensions}

For a Version~2.3 row,
\[
 n=v+o,
 \qquad m_1=\left\lceil\frac{or}{\ell}\right\rceil,
 \qquad M=or\ell,
 \qquad N_0=n\ell,
 \qquad N_{\rm raw}=m_1\ell^2r^2,
 \qquad K=m_1\binom{\ell+1}{2}.
\]
For a full structured lift, $\rank C_{R,\mathbf v}=M-K$ and
\[
 \dim\ker C_{R,\mathbf v}=N_0-M+K.
\]
If $\rank\mathcal O_{R,\mathbf v}=m_1\ell^2$, then
\[
 \dim_{\F_q}H_{R,\mathbf v}=N_0-m_1\ell^2,
 \qquad
 \dim_A H_{R,\mathbf v}=n-m_1\ell.
\]
The selected special-XL slice has $A$-dimension $s$, base-field dimension
$h=\ell s$, and retry exponent $\tau=K-h$.

\begin{table}[tbp]
\centering
\caption{Structured-lift dimensions for the six Version~2.3 $\ell=4$ rows.}
\label{tab:dimensions}
\small
\setlength{\tabcolsep}{5pt}
\begin{tabular}{@{}l r r r r r@{}}
\toprule
Parameters & $M$ & $K$ & $\dim_A H$ & $s$ & $\tau$\\
\midrule
$(28,5,4,4)$ & 80  & 50 & 13 & 10 & 10\\
$(28,4,4,5)$ & 80  & 50 & 12 & 10 & 10\\
$(40,7,4,4)$ & 112 & 70 & 19 & 15 & 10\\
$(38,5,4,5)$ & 100 & 70 & 15 & 15 & 10\\
$(50,9,4,4)$ & 144 & 90 & 23 & 19 & 14\\
$(52,6,4,6)$ & 144 & 90 & 22 & 19 & 14\\
\bottomrule
\end{tabular}
\end{table}

\section{Cost-Search Certificate}\label{app:cost-certificates}

For each $m_1\in\{5,7,9\}$ and every admissible integer $s$, the artifact
expands \eqref{eq:hilbert-series} exactly.  Here admissible means
$1\le s\le\min\{\dim_A H,\lfloor K/4\rfloor\}$.  The $K/4$ bounds are
$12,17,22$.  They bind or tie for both $m_1=5$ rows and both $m_1=9$ rows.
For $m_1=7$, the bounds are $17$ for $(40,7,4,4)$ and $15$ for
$(38,5,4,5)$, the latter coming from $\dim_A H$.  The artifact searches to
the larger group bound and then filters the certificate separately for each
parameter set; both $m_1=7$ rows attain the same optimum at $s=15$.
The monomial count
$M_s(\mathbf d)$ is coordinatewise increasing, so any multidegree that could
match or beat the incumbent lies in a finite rectangular box.  By symmetry it
suffices to enumerate sorted multidegrees.  The search checks 3,690, 11,217,
and 27,449 candidates for $m_1=5,7,9$, respectively.
 
The optima, unique up
to block permutation, are
those in \Cref{tab:xl-optima}.  No zero signed coefficient occurs anywhere
in these bounded searches, so using a strictly negative rather than a
non-positive trigger does not change the result.  The optimal coefficients are
\[
 -14{,}677{,}784,\qquad
 -124{,}983{,}335{,}190,\qquad
 -1{,}164{,}044{,}975{,}088{,}360.
\]
The corresponding Macaulay column counts are
\[
 23{,}417{,}049{,}656,
 \quad 44{,}237{,}557{,}981{,}725{,}696,
 \quad 236{,}094{,}291{,}515{,}544{,}000{,}000.
\]
Using the sharper full-rank retry bound would give published-model exponents
128.81713, 171.68562, and 214.12301.  The strengthened headline instead uses
$\mathsf{Retry}^{+}_{\tau}$, adding $\log_2(19/18)$ and giving
128.89513, 171.76362, and 214.20101.  Under the validated signed AXN circuit,
the corresponding headline exponents are 132.24987, 175.11835, and
217.55575.  The NAND sensitivity gives 133.11831, 175.98679, and 218.42418.

\paragraph{Low-state certificate.}
A second search allows every block degree $d_i\ge1$ and minimizes
$M_s(\mathbf d)$ subject to a cost budget.  Degree one is legitimate here:
the Macaulay columns still contain the homogenizer and every degree-one
variable in that block, while only diagonal equation multiples that exceed the
target degree are omitted.  For the published-model budgets $143/207/272$,
the exact global minima are
\[
\begin{array}{c|c|c|r|r}
\text{Level}&(s,\tau)&\mathbf d&[\mathbf t^{\mathbf d}]H&M_s(\mathbf d)\\\hline
\mathrm I&(8,18)&(1,2,3,3)&-70{,}875&11{,}026{,}125\\
\mathrm{III}&(9,34)&(1,1,2,4)&-16{,}396&3{,}932{,}500\\
\mathrm V&(11,46)&(1,2,2,3)&-31{,}536&26{,}574{,}912.
\end{array}
\]
Their published-model exponents are $140.12986$, $205.46174$, and
$262.52898$.  Under the AXN budgets $143/207/272$, the exact global minima are
$(s,\tau,\mathbf d)=(9,14,(2,2,3,4))$ at Level~I,
$(10,30,(1,2,3,3))$ at Level~III, and
$(11,46,(1,2,2,3))$ at Level~V.  Their Macaulay column counts are
$475{,}832{,}500$, $59{,}383{,}896$, and $26{,}574{,}912$, and their AXN
exponents are $137.69565$, $199.96188$, and $265.88371$.

Here is the finite-tail proof of globality.  For fixed $s$,
$M_s(\mathbf d)$ and each displayed cost are coordinatewise increasing.  A
sorted four-tuple with largest coordinate $D$ has
\[
 M_s(\mathbf d)\ge
 \binom{s+D}{D}(s+1)^3.
\]
Thus, once the right-hand side exceeds the incumbent, every larger degree is
irrelevant.  At smaller $s$, even the degree-one floor lies above the relevant
cost budget; at larger $s$, that floor or the first negative-Hilbert profile
already exceeds the incumbent.  Exact enumeration of every remaining finite
box gives the displayed unique minima up to block permutation.  The companion
JSON records every per-$s$ exclusion, every tested negative coefficient, and
the exact row counts in \eqref{eq:macaulay-row-count}.  As a boundary check,
the artifact also constructs 48 planted four-block systems across five
negative-Hilbert profiles containing degree-one blocks; every exact Macaulay
matrix has one-dimensional right kernel.

\paragraph{Three-plus-one certificate.}
For the series \eqref{eq:hilbert-series-three}, the artifact performs an
independent exact search over sorted triples $1\le d_1\le d_2\le d_3$ and every
admissible $s$.  It charges \eqref{eq:three-plus-one-stream-work}, the exact
square-row acceptance density, normal-preconditioner restarts over
$\F_{19^8}$, and the extraction reserve.  The compact time minima are
\[
\begin{array}{c|c|c|r|r|r}
\text{Level}&(s,\tau)&\mathbf d&[\mathbf t^{\mathbf d}]H^{(3)}&M^{(3)}&R^{(3)}\\\hline
\mathrm I&(9,14)&(4,4,5)&-416{,}990&1{,}023{,}472{,}450&2{,}958{,}813{,}000\\
\mathrm{III}&(11,26)&(3,3,6)&-1{,}321{,}334&1{,}639{,}770{,}496&4{,}286{,}113{,}104\\
\mathrm V&(13,38)&(3,4,4)&-646{,}618&3{,}172{,}064{,}000&7{,}332{,}242{,}400.
\end{array}
\]
Their compact/robust AXN exponents are respectively
$144.24760/148.16994$, $197.02045/200.94294$, and
$250.31880/254.24141$.  The Level-V robust minimum-state row is
$(s,\tau,\mathbf d)=(12,42,(2,3,5))$, with coefficient $-416{,}615$,
$M^{(3)}=256{,}214{,}140$, $R^{(3)}=559{,}191{,}087$, exponent
$263.40011$, and an eight-array state of $9.545$~GiB.

For every $d_i\ge1$,
$R^{(3)}/M^{(3)}\ge3m_1/(s+1)^2$, so any profile within a given work incumbent
has an explicit finite cap on $M^{(3)}$.  Sorted-degree monotonicity then gives
$M^{(3)}\ge(s+1)^2\binom{s+d_3}{d_3}$ and closes every omitted tail.  The
companion script records the per-$s$ caps and checks 507, 455, and 558 sorted
degree triples for Levels I, III, and V; 8, 35, and 60 have negative Hilbert
coefficient in the certified boxes.

\section{Same-Basis Ordered-Label Estimator Baseline}\label{app:ordered-baseline}

The comparison with the official Version~2.3 table mixes cost conventions, so
we also recost a reproducible \emph{estimator baseline} under exactly the same
AXN arithmetic basis as the selected attacks.  This baseline feeds the full
ordered-label equation/variable dimensions from the earlier affine-column
accounting into the pinned generic-MQ optimizer and replaces each charged
base-field operation by $g_1^{\mathrm{AXN}}=441$.  An attacker can indeed feed
the unreduced affine-column equations to a generic solver.  What is
counterfactual is the estimator's treatment of the ordered labels as though
they represented independent generic quadratic features: on every symmetric
key, reversed labels define the same polynomial by
\Cref{thm:symmetric-square}.  Thus the table is a controlled historical
estimator comparison, not evidence for a second quadratic space of the stated
ordered-label rank.  It reports every row rather than only the range quoted in
the introduction.

\begin{table}[H]
\centering
\caption{Historical ordered-label generic-MQ estimator and selected attack,
both recosted in the same AXN arithmetic basis.  ``Estimator input'' gives the
equation/variable dimensions of the unreduced affine-column equations.  The
system is executable, but treating its reversed labels as independent generic
quadratic features is counterfactual on a symmetric key.}
\label{tab:ordered-baseline}
\footnotesize
\setlength{\tabcolsep}{4.1pt}
\renewcommand{\arraystretch}{1.08}
\begin{tabular}{@{}c l c r r r@{}}
\toprule
Level & Parameters & \shortstack{Estimator\\input} & \shortstack{Ordered-label\\AXN} &
\shortstack{Selected\\AXN} & Saving\\
\midrule
I   & $(28,5,4,4)$  & $80/132$  & 206.69 & 132.25 &  74.44\\
I   & $(48,16,2,2)$ & $64/128$  & 171.77 & 133.82 &  37.95\\
I   & $(28,4,4,5)$  & $80/128$  & 206.69 & 132.25 &  74.44\\
\midrule
III & $(40,7,4,4)$  & $112/188$ & 277.52 & 175.12 & 102.40\\
III & $(72,24,2,2)$ & $96/192$  & 242.15 & 187.82 &  54.33\\
III & $(38,5,4,5)$  & $100/172$ & 250.29 & 175.12 &  75.17\\
\midrule
V   & $(50,9,4,4)$  & $144/236$ & 345.49 & 217.56 & 127.93\\
V   & $(96,32,2,2)$ & $128/256$ & 311.08 & 239.57 &  71.51\\
V   & $(52,6,4,6)$  & $144/232$ & 347.13 & 217.56 & 129.58\\
\bottomrule
\end{tabular}
\end{table}

The exact unrounded savings are $37.94724$--$129.57738$ bits.  This comparison
still inherits the generic-MQ and special-XL solving models, but it removes the
field-operation and gate-basis mismatch from the comparison itself.

\section{Fixed-Skew Certificate}\label{app:fixed-skew}

For the official $\ell=2$ matrix, direct arithmetic gives
\[
 9I+10S=\begin{pmatrix}0&1\\1&7\end{pmatrix},
 \qquad e_0^\transpose(9I+10S)=e_1^\transpose.
\]
Together with \eqref{eq:fixed-skew-offset}, this proves
\eqref{eq:fixed-skew-identity} without a rank or genericity assumption.  The
remaining public preflight checks $\rank Q_R=K$ and
$\rank C_{R,\mathbf v}=M-K$.

The resulting residual dimensions are $48/112$, $72/168$, and $96/224$.
The reoptimized generic-MQ convention gives per-solve exponents 129.44030,
183.43803, and 237.84310.  The last optimum uses Hashimoto's $a=1$ boundary
and the direct $\operatorname{MQ}(19,1,1)$ solver described in the main text.
Replacing the coarse factor $b_1$ by the validated signed AXN circuit
$g_1^{\mathrm{AXN}}$ gives per-solve arithmetic exponents 133.74643,
187.74416, and 242.14922.  By \Cref{cor:aggregate-full-space}, the headline condition
$\Theta_{\rm pol}\le19^{-1}$ adds less than $\log_2(20/19)$ bits, giving
Model/AXN exponents
129.51430/133.82043, 183.51203/187.81816, and
237.91710/242.22322.  The minimum ranks $K+1=49,73,97$ suffice
for the aggregate condition.  Positive AXN nominal arithmetic gap only requires
ranks 46, 68, and 89; the all-target condition $r\ge2K$ gives 96, 144, and
192.
These are the values in \Cref{tab:l2-costs} and
\Cref{tab:fixed-skew-evidence}.

\section{Artifact and Reproduction}\label{app:artifact}

The companion artifact, distributed separately from this manuscript
package, contains a self-contained evidence directory with:
\begin{itemize}
\item the public-key generator, direct verifier, and official Level-I
known-answer public key and signature;
\item the quotient, affine-feature, structured-offset, stable-kernel, and
rejection checks;
\item all 121 structured preflight records, all 13 structured-slice records,
and 60 fixed-skew preflights;
\item 102,796 structured cross-polar ranks, 21,992 dual directional-derivative
ranks, 3,888 fixed-skew polar ranks, and 29,424 zero-offset filter ranks,
including the exceptional-direction campaigns;
\item exhaustive $\ell=2$ relation enumeration, both rejection certificates,
accepted-target bounds, and generic-MQ costs;
\item the complete special-XL and two-block cost searches, 128 exact
$\ell=4$ Macaulay experiments, 48 additional degree-one boundary experiments,
48 two-block controls, and the target-generation bound;
\item 110 exhaustive toy checks of the exact collision and rank-spectrum formulas,
including rank-defective and compatibility-defective cases, plus genericity
diagnostics on six actual $\ell=2$ headline cores;
\item an independent rank implementation checked on 2,036 matrices, a
projective-distinctness audit of all 124,788 structured rank-test directions,
and independent scalar evaluators for the released gate circuits and headline
cost conversions;
\item the certified compact and joint-stress four-block time--state
frontiers, including exact Hilbert coefficients, row counts, preconditioner
bounds, monotone tail certificates, and the $\F_{19^{12}}$ Level-I extension
ledger;
\item the exact three-plus-one Hilbert search, streamed operator ledger,
linear-completion theorem, and compact/robust time--state certificate underlying
\Cref{tab:three-plus-one};
\item the explicit streamed normal-operator ledgers and global
state-minimality certificates underlying
\Cref{tab:explicit-stream,tab:robust-explicit}, including the
$16R(s+1)^2M$ and $32R(s+1)^2M$ work bounds and eight-array memory reserves;
\item the exact random-XOF seeded-spectrum and excess-spectrum evaluations
of \Cref{thm:seeded-spectrum,cor:full-ordinary-slices}, retaining the
modulo-$19$ atom $14/256$;
\item the complete $55$-pattern enumerator, exact reduction to $16$
Frobenius orbits, multihomogeneous B\'ezout and companion homotopy sums,
per-orbit extension/SLP ledgers, candidate filters, structural-preflight
density certificate, and the normalized fixed-core-free frontier of
\Cref{prop:selected-homotopy-frontier};
\item the complete-square homotopy ledger for $K=50,70,90$, including field
thresholds, root-success denominators, output ceilings, and break-even factors;
\item the exact necessary repair-floor calculations of
\Cref{prop:repair-floor,tab:repair-floor,prop:fresh-slice-repair-floor,tab:fresh-slice-repair-floor};
\item three exact reduced-size Macaulay checks at the production-selected
multidegrees, with full row and column counts and right-corank transcripts; and
\item the exact projective-spectrum enumeration cost calculation of
\Cref{prop:exhaustive-rank-cost}.
\end{itemize}
All finite-field ranks are computed exactly over $\F_{19}$.  The
companion artifact records hashes of the scripts, public keys, stable kernels,
selected slices, and polar bases, together with a manifest for every file.
None of the finite sampling in that artifact is presented as a proof of a
production aggregate spectrum.

\begingroup
\footnotesize
\printbibliography
\endgroup
\end{document}
